%% Generated by Sphinx.
\def\sphinxdocclass{jupyterBook}
\documentclass[letterpaper,10pt,english]{jupyterBook}
\ifdefined\pdfpxdimen
   \let\sphinxpxdimen\pdfpxdimen\else\newdimen\sphinxpxdimen
\fi \sphinxpxdimen=.75bp\relax
\ifdefined\pdfimageresolution
    \pdfimageresolution= \numexpr \dimexpr1in\relax/\sphinxpxdimen\relax
\fi
%% let collapsible pdf bookmarks panel have high depth per default
\PassOptionsToPackage{bookmarksdepth=5}{hyperref}
%% turn off hyperref patch of \index as sphinx.xdy xindy module takes care of
%% suitable \hyperpage mark-up, working around hyperref-xindy incompatibility
\PassOptionsToPackage{hyperindex=false}{hyperref}
%% memoir class requires extra handling
\makeatletter\@ifclassloaded{memoir}
{\ifdefined\memhyperindexfalse\memhyperindexfalse\fi}{}\makeatother

\PassOptionsToPackage{booktabs}{sphinx}
\PassOptionsToPackage{colorrows}{sphinx}

\PassOptionsToPackage{warn}{textcomp}

\catcode`^^^^00a0\active\protected\def^^^^00a0{\leavevmode\nobreak\ }
\usepackage{cmap}
\usepackage{fontspec}
\defaultfontfeatures[\rmfamily,\sffamily,\ttfamily]{}
\usepackage{amsmath,amssymb,amstext}
\usepackage{polyglossia}
\setmainlanguage{english}



\setmainfont{FreeSerif}[
  Extension      = .otf,
  UprightFont    = *,
  ItalicFont     = *Italic,
  BoldFont       = *Bold,
  BoldItalicFont = *BoldItalic
]
\setsansfont{FreeSans}[
  Extension      = .otf,
  UprightFont    = *,
  ItalicFont     = *Oblique,
  BoldFont       = *Bold,
  BoldItalicFont = *BoldOblique,
]
\setmonofont{FreeMono}[
  Extension      = .otf,
  UprightFont    = *,
  ItalicFont     = *Oblique,
  BoldFont       = *Bold,
  BoldItalicFont = *BoldOblique,
]



\usepackage[Bjarne]{fncychap}
\usepackage[,numfigreset=1,mathnumfig]{sphinx}

\fvset{fontsize=\small}
\usepackage{geometry}


% Include hyperref last.
\usepackage{hyperref}
% Fix anchor placement for figures with captions.
\usepackage{hypcap}% it must be loaded after hyperref.
% Set up styles of URL: it should be placed after hyperref.
\urlstyle{same}

\addto\captionsenglish{\renewcommand{\contentsname}{Lattice Boltzmann Method - Application to Different Partial Differential Equation}}

\usepackage{sphinxmessages}



        % Start of preamble defined in sphinx-jupyterbook-latex %
         \usepackage[Latin,Greek]{ucharclasses}
        \usepackage{unicode-math}
        % fixing title of the toc
        \addto\captionsenglish{\renewcommand{\contentsname}{Contents}}
        \hypersetup{
            pdfencoding=auto,
            psdextra
        }
        % End of preamble defined in sphinx-jupyterbook-latex %
        

\title{Ricardo Leite Martins Bazarin - Academic Documentation}
\date{Oct 08, 2026}
\release{}
\author{Ricardo Leite Martins Bazarin}
\newcommand{\sphinxlogo}{\vbox{}}
\renewcommand{\releasename}{}
\makeindex
\begin{document}

\pagestyle{empty}
\sphinxmaketitle
\pagestyle{plain}
\sphinxtableofcontents
\pagestyle{normal}
\phantomsection\label{\detokenize{index::doc}}


\noindent{\sphinxincludegraphics[width=250\sphinxpxdimen]{{personalphoto}.png}\hspace*{\fill}}

\sphinxAtStartPar
Mechanical Engineer, graduated from the Federal Technological University of Paraná (UTFPR), Guarapuava/🇧🇷 campus. He holds both
a master’s degree and a doctorate in the Thermal Sciences area from the Graduate Program in Mechanical and Materials Engineering
at UTFPR, Curitiba campus. He has participated in research groups on Porous Media and Rheology, where he developed work in
numerical modeling applied to Thermal Sciences and Fluid Mechanics, with emphasis on the use and development of the Lattice
Boltzmann Method. Currently, he is a researcher at the Federal University of Santa Catarina (UFSC), specializing in digital
rock simulation using the Lattice Boltzmann Method on High\sphinxhyphen{}Performance Computing (HPC) systems.
\begin{itemize}
\item {} 
\sphinxAtStartPar
🔗 \sphinxhref{https://buscatextual.cnpq.br/buscatextual/visualizacv.do?id=K8733837Z4}{Curriculum Lattes (Brazilian Academic Curriculum)}

\item {} 
\sphinxAtStartPar
💻 \sphinxhref{https://github.com/Ricardoleite}{GitHub Page}

\item {} 
\sphinxAtStartPar
👓 \sphinxhref{https://scholar.google.com.br/citations?user=ybe2TbQAAAAJ\&hl=pt-BR}{Google Scholar}

\end{itemize}

\begin{DUlineblock}{0em}
\item[] \sphinxstylestrong{\Large Articles Published in Periodics}
\end{DUlineblock}
\begin{itemize}
\item {} 
\sphinxAtStartPar
\sphinxhref{https://doi.org/10.1007/s11242-020-01530-w}{2021\sphinxhyphen{}Boundary Effects on the Tortuosity and Permeability of Idealized Porous Media}
\begin{itemize}
\item {} 
\sphinxAtStartPar
Authors: Bazarin, De Lai, Naaktgeboren, Junqueira.

\item {} 
\sphinxAtStartPar
Area: (Porous media fluid\sphinxhyphen{}flow) Single\sphinxhyphen{}phase porous properties characterization in fractal porous media

\end{itemize}

\item {} 
\sphinxAtStartPar
\sphinxhref{https://doi.org/10.1016/j.compfluid.2021.105142}{2021\sphinxhyphen{}Moments\sphinxhyphen{}based method for boundary conditions in the lattice Boltzmann framework: A comparative analysis for the lid driven cavity flow}
\begin{itemize}
\item {} 
\sphinxAtStartPar
Authors: Bazarin, Philippi, Randles, Hegele Jr.

\item {} 
\sphinxAtStartPar
Area: (Boundary conditions modeling)

\end{itemize}

\item {} 
\sphinxAtStartPar
\sphinxhref{https://doi.org/10.1103/PhysRevE.110.015303}{2024\sphinxhyphen{}Improved lattice Boltzmann model for immiscible multicomponent systems with high viscosity gradients at the interface}
\begin{itemize}
\item {} 
\sphinxAtStartPar
Authors: Bazarin, Naaktgeboren, Junqueira, Philippi, Hegele Jr.

\item {} 
\sphinxAtStartPar
Area: (Multicomponent/multiphase modeling) Numerical modeling for multicomponent fluid\sphinxhyphen{}flow

\end{itemize}

\item {} 
\sphinxAtStartPar
\sphinxhref{https://doi.org/10.1063/5.0264878}{2025\sphinxhyphen{}Lattice Boltzmann method assessment for the Buckley–Leverett equation: Extension to dynamic capillary effect}
\begin{itemize}
\item {} 
\sphinxAtStartPar
Authors: Bazarin, dos Santos, Siebert.

\item {} 
\sphinxAtStartPar
Area: (Porous media fluid\sphinxhyphen{}flow \sphinxhyphen{} Multicomponent/multiphase modeling) Numerical modeling of multicomponent/multiphase flow in porous\sphinxhyphen{}continuous porous media.

\item {} 
\sphinxAtStartPar
Suplemental Materials \sphinxhyphen{} \sphinxhref{https://github.com/poro-labcc/Buckley-Leverett-Equation-Non-Linear-Convective-Diffusive-Lattice-Boltzmann-Method}{C++ code of lattice Boltzmann simulation presented.}

\end{itemize}

\end{itemize}

\begin{DUlineblock}{0em}
\item[] \sphinxstylestrong{\large Academic Material}
\end{DUlineblock}

\sphinxAtStartPar
The full table of contents is generated automatically from the \sphinxcode{\sphinxupquote{\_toc.yml}} file
and appears in the sidebar (HTML) or in the table of contents (PDF).

\sphinxstepscope


\part{Lattice Boltzmann Method \sphinxhyphen{} Application to Different Partial Differential Equation}

\sphinxstepscope


\chapter{Diffusive Equation}
\label{\detokenize{LBM-Study/Diffusive-Equation:diffusive-equation}}\label{\detokenize{LBM-Study/Diffusive-Equation::doc}}


\sphinxAtStartPar
The diffusive equation is described by
\begin{equation}\label{equation:LBM-Study/Diffusive-Equation:df-eq-lb}
\begin{split}
\partial_t \phi - \partial_{\alpha}\left( \nu \partial_{\alpha}\phi \right) = 0
\end{split}
\end{equation}
\sphinxAtStartPar
where \(\phi(\alpha,t)\) is a time and space dependent parameter and \(\nu\) is diffusive coeficient.


\section{Numerical Discretization}
\label{\detokenize{LBM-Study/Diffusive-Equation:numerical-discretization}}
\begin{figure}[htbp]
\centering
\capstart

\noindent\sphinxincludegraphics[scale=2.5]{{mesh-1-diffusive}.png}
\caption{1D regular grid example for diffusive equation.}\label{\detokenize{LBM-Study/Diffusive-Equation:mesh-1-diffusive}}\end{figure}


\subsection{FDM Discretization}
\label{\detokenize{LBM-Study/Diffusive-Equation:fdm-discretization}}
\sphinxAtStartPar
Discretizing the Eq. \eqref{equation:LBM-Study/Diffusive-Equation:df-eq-lb} through finite diference method, we have for forward time and central space (FTCS) discretizations:
\begin{equation*}
\begin{split}
\partial_t \phi - \partial_{\alpha}\left( \nu \partial_{\alpha}\phi \right) = \frac{\phi(x,t + \Delta_{t}) - \phi(x,t) }{\Delta_{t}} - \nu\frac{ \phi(x - \Delta_{x},t) -2\phi(x,t) + \phi(x + \Delta_{x},t) }{\Delta_{x}^{2}} =0,
\end{split}
\end{equation*}
\sphinxAtStartPar
for \(\nu\) constant and isolating the term \(u(x,t + \Delta_{t})\), we can describe the \(\phi\) time evolution:
\begin{equation*}
\begin{split}
\phi(x,t + \Delta_{t})   = \phi(x,t) + \frac{\nu\Delta_{t}}{\Delta_{x}^{2}} \left( \phi(x - \Delta_{x},t) -2\phi(x,t) + \phi(x + \Delta_{x},t) \right) .
\end{split}
\end{equation*}

\subsection{Lattice Boltzmann Equation}
\label{\detokenize{LBM-Study/Diffusive-Equation:lattice-boltzmann-equation}}
\sphinxAtStartPar
Describing the problem through the BGK lattice Boltzmann equation (BGK\sphinxhyphen{}LB):
\begin{equation}\label{equation:LBM-Study/Diffusive-Equation:LB-df-Eq}
\begin{split}
f_i( x_{\alpha} + c_{i,\alpha} \delta t, t+\delta t) = f_i(x_{\alpha}, t)  -\left( \frac{ f_{i} - f_{i}^{eq} }{ \tau } \right), 
\end{split}
\end{equation}
\sphinxAtStartPar
where the equilibrium distribution function is defined by
\begin{equation*}
\begin{split}
f^{eq}_{i} = w_{i}\phi(x_{\alpha}, t),
\end{split}
\end{equation*}
\sphinxAtStartPar
and the equilibrium moments are given by
\begin{equation}\label{equation:LBM-Study/Diffusive-Equation:moments-fi-df}
\begin{split}
\displaystyle\sum_{i=0} f_{i}^{eq}=\phi, \quad \quad  \displaystyle\sum_{i}f_{i}^{eq} c_{i,\alpha} =0 \quad \quad \textrm{and} \quad \quad  \displaystyle\sum_{i}f_{i}^{eq} c_{i,\alpha} c_{i,\beta} =c_{s}^{2} \phi \delta_{\alpha\beta}.
\end{split}
\end{equation}
\sphinxAtStartPar
Through Chapman–Enskog analysis, it is demonstrated that the first\sphinxhyphen{}order non\sphinxhyphen{}equilibrium moment describes the pressure gradient and, consequently, the average velocity:
\begin{equation*}
\begin{split}
m^{neq}_{\alpha} = \displaystyle\sum_{i} c_{i,\alpha}\left( f_i-f_i^{eq}\right)  = -(c_{s}^{2} \tau \delta_{t})\partial_{\alpha}\phi  \quad \quad \textrm{and} \quad\quad \tau = \displaystyle\frac{\nu}{c_{s}^{2} \delta_{t}} +\frac{1}{2}.
\end{split}
\end{equation*}

\subsubsection{Lattice Direction Moments}
\label{\detokenize{LBM-Study/Diffusive-Equation:lattice-direction-moments}}
\begin{sphinxuseclass}{cell}
\begin{sphinxuseclass}{tag_hide-input}
\end{sphinxuseclass}
\end{sphinxuseclass}
\begin{sphinxuseclass}{cell}
\begin{sphinxuseclass}{tag_hide-input}\begin{sphinxVerbatimOutput}

\begin{sphinxuseclass}{cell_output}\begin{equation*}
\begin{split}\displaystyle \begin{gathered}\underbrace{\sum_{i=0} f_{i}^{eq} =\sum_{i=0} f_{i} }_{\textrm{Zero-Order Moment}} =\phi,\quad \quad \underbrace{\sum_{i=0} f_{i}^{eq}c_{i,x} }_{\textrm{x-First-Order Moment}} =0,\quad \quad \underbrace{\sum_{i=0} f_{i}^{eq}c_{i,y} }_{\textrm{y-First-Order Moment}} =0,\\ \quad \quad \underbrace{\sum_{i=0} f_{i}^{eq}c_{i,x}c_{i,x} }_{\textrm{xx-Second-Order Moment}} =\frac{\phi}{3},\quad \quad \underbrace{\sum_{i=0} f_{i}^{eq}c_{i,x}c_{i,y} }_{\textrm{xy-Second-Order Moment}} =0 \quad \quad and \quad \quad \underbrace{\sum_{i=0} f_{i}^{eq}c_{i,y}c_{i,y} }_{\textrm{yy-Second-Order Moment}} =\frac{\phi}{3}\end{gathered}\end{split}
\end{equation*}
\end{sphinxuseclass}\end{sphinxVerbatimOutput}

\end{sphinxuseclass}
\end{sphinxuseclass}

\subsubsection{Chapmann\sphinxhyphen{}Enskog Analysis}
\label{\detokenize{LBM-Study/Diffusive-Equation:chapmann-enskog-analysis}}
\begin{sphinxuseclass}{toggle}
\sphinxAtStartPar
Applying the Chapmann\sphinxhyphen{}Enskog procedure to LB equation, we expand the term \(f_{i}\left(\boldsymbol{x}+\boldsymbol{e}_i \Delta t, t+\Delta t\right)\) in a Taylor series to recover the derivative form of the equation, i.e.,
\begin{equation}\label{equation:LBM-Study/Diffusive-Equation:EqExp-df-Eq}
\begin{split}
f_{i}\left( x_{\alpha} + c_{i,\alpha} \Delta t, t+\Delta t\right)- f_{i}( x_{\alpha}, t)=\displaystyle\sum_{j=1}^{\infty}\frac{\Delta t^{j}}{j!}D_{t}^{j}f_{i}= -\left( \frac{ f_{i} - f_{i}^{eq} }{ \tau } \right).
\end{split}
\end{equation}
\sphinxAtStartPar
Rescaling the dimensionless form of the Eq. \eqref{equation:LBM-Study/Diffusive-Equation:EqExp-df-Eq} in terms of the Knudsen number (\(Kn\)), we have
\begin{equation*}
\begin{split}
\displaystyle \sum^{\infty}_{j=1}  \frac{Kn^{(j-1)}}{j!} D_{t}^{j} f_{i} = - \frac{1}{Kn}\left( \frac{ f_{i} - f_{eq,i} }{ \tau } \right) \quad \quad \rightarrow \quad \quad \displaystyle \sum^{\infty}_{j=1}  \frac{Kn^{(j)}}{j!} D_{t}^{j} f_{i} = -  \frac{ f_{i} - f_{i}^{eq} }{ \tau },
\end{split}
\end{equation*}
\sphinxAtStartPar
applying the asymptotic expansion in both the distribution function (\(f_{i}=f_{i}^{(0)}+Kn f_{i}^{(1)} + Kn^{2} f_{i}^{(2)}+\cdots\)) and time partial derivative (\(\partial_{t}=\partial_{t}^{(0)}+ Kn \partial_{t}^{(1)}+Kn^{2} \partial_{t}^{(2)}+\cdots\)), and separating the equation in orders up to the order \(Kn^{2}\):
\begin{equation}\label{equation:LBM-Study/Diffusive-Equation:Chap-Kn-df-Eq}
\begin{split}
\begin{array}{ll}
    (Kn^{(0)}):& f_{i}^{(0)}= f_{i}^{eq},\\
    (Kn^{(1)}):& \left( \partial_{t}^{(0)} + c_{i,\alpha}\partial_{\alpha} \right)f_{i}^{(0)}= - \displaystyle\frac{ f_{i}^{(1)} }{ \tau } , \\
    (Kn^{(2)}):& \partial_{t}^{(1)} f_{i}^{(0)} + \left( \partial_{t}^{(0)} + c_{i,\alpha}\partial_{\alpha} \right)f_{i}^{(1)} + \displaystyle\frac{\left( \partial_{t}^{(0)} + c_{i,\alpha}\partial_{\alpha} \right)^{2} f_{i}^{(0)}}{2}= - \displaystyle\frac{ f_{i}^{(2)} }{ \tau },\\
    \textrm{ou}&  \\
    (Kn^{(2)}):& \partial_{t}^{(1)} f_{i}^{(0)} + \left( \partial_{t}^{(0)} + c_{i,\alpha}\partial_{\alpha} \right)\displaystyle\left(1 - \frac{1}{2\tau}\right) f_{i}^{(1)} =  - \displaystyle\frac{ f_{i}^{(2)} }{ \tau } \quad\quad \rightarrow \quad\quad  f_{i}^{(1)}\textrm{ Formulation} \\
    \textrm{ou}& \\
    (Kn^{(2)}):& \partial_{t}^{(1)} f_{i}^{(0)} + \left( \partial_{t}^{(0)} + c_{i,\alpha}\partial_{\alpha} \right)^{2} \displaystyle\left(\frac{1}{2} - \tau\right) f_{i}^{(0)} =  - \displaystyle\frac{ f_{i}^{(2)} }{ \tau } \quad\quad \rightarrow  \quad\quad f_{i}^{(0)}\textrm{ Formulation}.
 \end{array}
\end{split}
\end{equation}
\sphinxAtStartPar
\sphinxstylestrong{Zero\sphinxhyphen{}Order Moment Balance}

\sphinxAtStartPar
To retrieve the balance equation, we sum the Eq. \eqref{equation:LBM-Study/Diffusive-Equation:Chap-Kn-df-Eq} for \(Kn^{(1)}\) and \(Kn^{(2)}\) over \(\sum_{i=0} \):
\begin{equation*}
\begin{split}
\begin{array}{l}
(Kn^{(1)}): \displaystyle\sum_{i}\left( \partial_{t}^{(0)} + c_{i,\alpha}\partial_{\alpha} \right)f_{i}^{(0)}= \displaystyle\sum_{i} \left(- \displaystyle\frac{ f_{i}^{(1)} }{ \tau } \right),\\
(Kn^{(1)}): \partial_{t}^{(0)}  \phi =0,
\end{array}
\end{split}
\end{equation*}
\sphinxAtStartPar
and
\begin{equation*}
\begin{split}
\begin{array}{l}
(Kn^{(2)}): \displaystyle\sum_{i}\left( \partial_{t}^{(1)} f_{i}^{(0)} + \left( \partial_{t}^{(0)} + c_{i,\alpha}\partial_{\alpha} \right)^{2} \displaystyle\left( \frac{1}{2}-\tau\right) f_{i}^{(0)} \right) =  \displaystyle\sum_{i}\left( - \displaystyle\frac{ f_{i}^{(2)} }{ \tau } \right), \\
(Kn^{(2)}): \partial_{t}^{(1)}\phi + \left( \displaystyle\frac{1}{2}-\tau\right) \partial_{\alpha}\partial_{\beta} \left(c_{s}^{2} \phi \delta_{\alpha\beta} \right) =0. 
\end{array}
\end{split}
\end{equation*}
\sphinxAtStartPar
or
\begin{equation}\label{equation:LBM-Study/Diffusive-Equation:df-Kn1-0-B0}
\begin{split}
\begin{array}{ll}
(Kn^{(2)}):& \displaystyle\sum_{i}\left[ \partial_{t}^{(1)} f_{i}^{(0)} + \left(1-\displaystyle\frac{1}{2\tau}\right)\left( \partial_{t}^{(0)} + c_{i,\alpha}\partial_{\alpha} \right) f_{i}^{(1)}   \right]=  \displaystyle\sum_{i}\left( - \displaystyle\frac{f_{i}^{(2)}}{\tau} \right), \nonumber \\
(Kn^{(2)}):& \partial_{t}^{(1)}\phi + \left(1-\displaystyle\frac{1}{2\tau}\right) m^{(1)}_{\alpha} = 0.
\end{array}
\end{split}
\end{equation}
\sphinxAtStartPar
where \(m^{(1)}_{\alpha}\) is the first\sphinxhyphen{}order moment of \(f_{i}^{(1)}\) and \(\sum_{i=1}f_{i}^{(j)}=0\) for \(j\geq 1\) due to imposition of \(\phi\) conservation. To compute the moment \(m^{(1)}_{\alpha}\), we multiply Eq. \eqref{equation:LBM-Study/Diffusive-Equation:Chap-Kn-df-Eq} by \(c_{i,\alpha}\) and sum over all lattice directions:
\begin{equation}\label{equation:LBM-Study/Diffusive-Equation:df-Kn0-B1}
\begin{split}
\begin{array}{l}
(Kn^{(1)}): \displaystyle\sum_{i}c_{i,\alpha}\left( \partial_{t}^{(0)} + e_{\beta,i}\partial_{\beta} \right)f_{i}^{(0)}= \displaystyle\sum_{i} c_{i,\alpha} \left( -\frac{f_{i}^{(1)}}{\tau}  \right),\\
(Kn^{(1)}): m^{(1)}_{\alpha} =-\tau \partial_{\beta}\left( c_{s}^{2} \phi \delta_{\alpha\beta}  \right).
\end{array}
\end{split}
\end{equation}
\sphinxAtStartPar
By substituting Eq. \eqref{equation:LBM-Study/Diffusive-Equation:df-Kn0-B1} into Eq. \eqref{equation:LBM-Study/Diffusive-Equation:df-Kn1-0-B0}, we recover the Eq. \eqref{equation:LBM-Study/Diffusive-Equation:df-eq-lb}, where \(\nu=c_{s}^{2}(\tau-1/2)\). In the regularized formulation of the lattice Boltzmann equation, the first\sphinxhyphen{}order correction is approximated as \(f_{i}^{(1)}\approx (f_{i}-f_{i}^{(0)})=f_{i}^{neq}\), where higher\sphinxhyphen{}order Knudsen moments are filtered out and the collision term is reformulated in terms of \(f_{i}^{neq}\).

\end{sphinxuseclass}

\subsubsection{Boundary Conditions}
\label{\detokenize{LBM-Study/Diffusive-Equation:boundary-conditions}}
\sphinxAtStartPar
The boundary conditions for the lattices can be derived by solving a linear system of known moments.


\begin{savenotes}\sphinxattablestart
\sphinxthistablewithglobalstyle
\centering
\begin{tabulary}{\linewidth}[t]{TTTTTT}
\sphinxtoprule
\sphinxstyletheadfamily 
\sphinxAtStartPar
D1Q2:
&\sphinxstyletheadfamily 
\sphinxAtStartPar
Boundaries
&
\sphinxAtStartPar

&\sphinxstyletheadfamily 
\sphinxAtStartPar
Layers
&
\sphinxAtStartPar

&
\sphinxAtStartPar

\\
\sphinxmidrule
\sphinxtableatstartofbodyhook
\sphinxAtStartPar

&
\sphinxAtStartPar

&
\sphinxAtStartPar

&
\sphinxAtStartPar
West
&
\sphinxAtStartPar

&
\sphinxAtStartPar
East
\\
\sphinxhline
\sphinxAtStartPar

&
\sphinxAtStartPar
Unknown \(f_{i}\)
&
\sphinxAtStartPar

&
\sphinxAtStartPar
\(f_2=\phi_{e} - f_1\)
&
\sphinxAtStartPar

&
\sphinxAtStartPar
\(f_1=\phi_{w} - f_2\)
\\
\sphinxhline
\sphinxAtStartPar
D1Q3:
&
\sphinxAtStartPar
Boundaries
&
\sphinxAtStartPar

&
\sphinxAtStartPar
Layers
&
\sphinxAtStartPar

&
\sphinxAtStartPar

\\
\sphinxhline
\sphinxAtStartPar

&
\sphinxAtStartPar

&
\sphinxAtStartPar

&
\sphinxAtStartPar
West
&
\sphinxAtStartPar

&
\sphinxAtStartPar
East
\\
\sphinxhline
\sphinxAtStartPar

&
\sphinxAtStartPar
Unknown \(f_{i}\)
&
\sphinxAtStartPar

&
\sphinxAtStartPar
\(f_2=\phi_{e} - f_1 - f_0\)
&
\sphinxAtStartPar

&
\sphinxAtStartPar
\(f_1=\phi_{w} - f_2 - f_0\)
\\
\sphinxbottomrule
\end{tabulary}
\sphinxtableafterendhook\par
\sphinxattableend\end{savenotes}




\section{1st Benchmark: Single sine mode on a finite rod (Dirichlet)}
\label{\detokenize{LBM-Study/Diffusive-Equation:st-benchmark-single-sine-mode-on-a-finite-rod-dirichlet}}
\sphinxAtStartPar
PDE: \((\phi_t = \partial_{x}(\nu\partial_{x}\phi), (0<x<L,\ t>0)\)

\sphinxAtStartPar
BC: \((\phi(0,t)=\phi(L,t)=0)\)

\sphinxAtStartPar
IC: \(\phi(x,0)=\sin \left(\frac{\pi x}{L}\right)\)

\sphinxAtStartPar
\sphinxstylestrong{Solution}
\begin{equation*}
\begin{split}
\phi(x,t)=\sin\left(\frac{\pi x}{L}\right) \exp \left(-\nu \frac{\pi^2}{L^2} t\right).
\end{split}
\end{equation*}

\subsection{LBM D1Q3 Solution (\protect\(\tau=1\protect\))}
\label{\detokenize{LBM-Study/Diffusive-Equation:lbm-d1q3-solution-tau-1}}
\begin{sphinxuseclass}{cell}
\begin{sphinxuseclass}{tag_hide-input}
\begin{sphinxuseclass}{tag_hide-output}
\end{sphinxuseclass}
\end{sphinxuseclass}
\end{sphinxuseclass}

\subsection{FDM Solution}
\label{\detokenize{LBM-Study/Diffusive-Equation:fdm-solution}}
\begin{sphinxuseclass}{cell}
\begin{sphinxuseclass}{tag_hide-input}
\begin{sphinxuseclass}{tag_hide-output}
\end{sphinxuseclass}
\end{sphinxuseclass}
\end{sphinxuseclass}

\subsection{Numerical and Analytical Comparisson}
\label{\detokenize{LBM-Study/Diffusive-Equation:numerical-and-analytical-comparisson}}
\begin{sphinxuseclass}{cell}
\begin{sphinxuseclass}{tag_hide-input}\begin{sphinxVerbatimOutput}

\begin{sphinxuseclass}{cell_output}
\noindent\sphinxincludegraphics{{55234d5692dce1c246184742da092377313fdd7752e743ec6a2e20557d7e0a75}.png}

\end{sphinxuseclass}\end{sphinxVerbatimOutput}

\end{sphinxuseclass}
\end{sphinxuseclass}


\begin{sphinxuseclass}{cell}
\begin{sphinxuseclass}{tag_hide-input}
\begin{sphinxuseclass}{tag_hide-output}
\end{sphinxuseclass}
\end{sphinxuseclass}
\end{sphinxuseclass}


\begin{sphinxuseclass}{cell}
\begin{sphinxuseclass}{tag_hide-input}
\begin{sphinxuseclass}{tag_hide-output}
\end{sphinxuseclass}
\end{sphinxuseclass}
\end{sphinxuseclass}
\sphinxstepscope


\chapter{Poisson Equation (Steady\sphinxhyphen{}State Diffusive Equation)}
\label{\detokenize{LBM-Study/Poisson-Equation:poisson-equation-steady-state-diffusive-equation}}\label{\detokenize{LBM-Study/Poisson-Equation::doc}}
\sphinxAtStartPar
The Poisson is partial differential equation employed to represent problems as the pressure distribution over steady\sphinxhyphen{}state in space. It is characterized by a steady\sphinxhyphen{}state version of the difusive equation, i.e, the parameter is independent of time. The partial differential equation is given by
\begin{equation}\label{equation:LBM-Study/Poisson-Equation:poisson-eq-lb}
\begin{split}
\partial_{\alpha}\left( \nu \partial_{\alpha}p \right) = 0
\end{split}
\end{equation}
\sphinxAtStartPar
where \(p\) is the pressure and \(\nu\) diffusive control paremeter.


\section{Numerical Discretization}
\label{\detokenize{LBM-Study/Poisson-Equation:numerical-discretization}}
\begin{figure}[htbp]
\centering
\capstart

\noindent\sphinxincludegraphics[scale=2.5]{{mesh-1-diffusive-steady}.png}
\caption{1D regular grid example for diffusive equation.}\label{\detokenize{LBM-Study/Poisson-Equation:mesh-1-diffusive-steady}}\end{figure}


\subsection{FDM Discretization}
\label{\detokenize{LBM-Study/Poisson-Equation:fdm-discretization}}
\sphinxAtStartPar
Discretizing the Eq. \eqref{equation:LBM-Study/Diffusive-Equation:df-eq-lb} through finite diference method, we have for forward time and central space (FTCS) discretizations:
\begin{equation*}
\begin{split}
\partial_{\alpha}\left( \nu \partial_{\alpha}p \right) = \nu\frac{ p(x - \Delta_{x},t) -2p(x,t) + p(x + \Delta_{x},t) }{\Delta_{x}^{2}} =0,
\end{split}
\end{equation*}
\sphinxAtStartPar
for \(\nu\) constant and isolating the term \(u(x,t + \Delta_{t})\), we can describe the \(p\) time evolution:
\begin{equation*}
\begin{split}
p(x,t)   =  \frac{ p(x - \Delta_{x},t) + p(x + \Delta_{x},t) }{2} .
\end{split}
\end{equation*}

\subsection{Lattice Boltzmann Equation}
\label{\detokenize{LBM-Study/Poisson-Equation:lattice-boltzmann-equation}}
\sphinxAtStartPar
Describing the problem through the BGK lattice Boltzmann equation (BGK\sphinxhyphen{}LB) for the lattice D1Q3:
\begin{equation}\label{equation:LBM-Study/Poisson-Equation:LB-Poisson-Eq}
\begin{split}
f_i( x_{\alpha} + c_{i,\alpha} \delta t, t'+\delta t) = f_i(x_{\alpha}, t')  -\left( \frac{ f_{i} - f_{i}^{eq} }{ \tau } \right), 
\end{split}
\end{equation}
\sphinxAtStartPar
the equilibrium distribution function is defined by
\begin{equation*}
\begin{split}
f^{eq}_{i} = \left\{ \begin{array}{ll}
   \left(w_{0} - 1\right)\widehat{p}(x_{\alpha}, t')  & i=0 \\
   w_{i}\widehat{p}(x_{\alpha}, t')  & i\neq 0
\end{array} \right. ,
\end{split}
\end{equation*}
\sphinxAtStartPar
where \(\widehat{p}\) is a shifted pressure, shifted by a reference constant pressure \(p_{0}\), i.e., \(p=p_{0}+\widehat{p}\) (This assumption is necessary due to pressure by definition can not be negative, but negative pressure values can be found depending of initial simulation set). The equilibrium moments are given by
\begin{equation}\label{equation:LBM-Study/Poisson-Equation:moments-fi}
\begin{split}
\displaystyle\sum_{i=0} f_{i}^{eq}=0, \quad \quad \frac{1}{1-w_{0}}\displaystyle\sum_{i=1} f_{i}^{eq} = \frac{1}{1-w_{0}} \displaystyle\sum_{i=1} f_{i}=\widehat{p}, \quad \quad  \displaystyle\sum_{i}f_{i}^{eq} c_{i,\alpha} =0 \quad \quad \textrm{and} \quad \quad  \displaystyle\sum_{i}f_{i}^{eq} c_{i,\alpha} c_{i,\beta} =c_{s}^{2} \widehat{p} \delta_{\alpha\beta} .
\end{split}
\end{equation}
\sphinxAtStartPar
Through Chapman–Enskog analysis, it is demonstrated that the first\sphinxhyphen{}order non\sphinxhyphen{}equilibrium moment describes the pressure gradient and, consequently, the average velocity:
\begin{equation*}
\begin{split}
m^{neq}_{\alpha} = \displaystyle\sum_{i} c_{i,\alpha}\left( f_i-f_i^{eq}\right)  = -(c_{s}^{2} \tau \delta_{t})\partial_{\alpha}\widehat{p} \quad \quad \textrm{and} \quad\quad \tau = \displaystyle\frac{\nu}{c_{s}^{2} \delta_{t}} +\frac{1}{2}.
\end{split}
\end{equation*}

\subsubsection{Lattice Direction Moments}
\label{\detokenize{LBM-Study/Poisson-Equation:lattice-direction-moments}}
\begin{sphinxuseclass}{cell}
\begin{sphinxuseclass}{tag_hide-input}
\end{sphinxuseclass}
\end{sphinxuseclass}
\begin{sphinxuseclass}{cell}
\begin{sphinxuseclass}{tag_hide-input}\begin{sphinxVerbatimOutput}

\begin{sphinxuseclass}{cell_output}\begin{equation*}
\begin{split}\displaystyle \begin{gathered}\underbrace{\sum_{i=0} f_{i}^{eq} =\sum_{i=0} f_{i} }_{\textrm{Zero-Order Moment}} =0,\quad \quad \underbrace{\frac{1}{1-w_{0}}\sum_{i=1} f_{i}^{eq} = \frac{1}{1-w_{0}}\sum_{i=1} f_{i} }_{\textrm{Zero-Order Moment}} =\widehat{p},\quad \quad \underbrace{\sum_{i=0} f_{i}^{eq}c_{i,x} }_{\textrm{x-First-Order Moment}} =0,\quad \quad \underbrace{\sum_{i=0} f_{i}^{eq}c_{i,y} }_{\textrm{y-First-Order Moment}} =0\end{gathered}\end{split}
\end{equation*}\begin{equation*}
\begin{split}\displaystyle \begin{gathered}\underbrace{\sum_{i=0} f_{i}^{eq}c_{i,x}c_{i,x} }_{\textrm{xx-Second-Order Moment}} =\frac{\widehat{p}}{3},\quad \quad \underbrace{\sum_{i=0} f_{i}^{eq}c_{i,x}c_{i,y} }_{\textrm{xy-Second-Order Moment}} =0 \quad \quad and \quad \quad \underbrace{\sum_{i=0} f_{i}^{eq}c_{i,y}c_{i,y} }_{\textrm{yy-Second-Order Moment}} =\frac{\widehat{p}}{3}\end{gathered}\end{split}
\end{equation*}
\end{sphinxuseclass}\end{sphinxVerbatimOutput}

\end{sphinxuseclass}
\end{sphinxuseclass}

\subsubsection{Chapmann\sphinxhyphen{}Enskog Analysis}
\label{\detokenize{LBM-Study/Poisson-Equation:chapmann-enskog-analysis}}
\begin{sphinxuseclass}{toggle}
\sphinxAtStartPar
Applying the Chapmann\sphinxhyphen{}Enskog procedure to LB equation, we expand the term \(f_{i}\left(\boldsymbol{x}+\boldsymbol{e}_i \Delta t, t+\Delta t\right)\) in a Taylor series to recover the derivative form of the equation, i.e.,
\begin{equation}\label{equation:LBM-Study/Poisson-Equation:EqExp-Poisson-Eq}
\begin{split}
f_{i}\left(x_{\alpha} + c_{i,\alpha} \Delta t, t+\Delta t\right)- f_{i}(\boldsymbol{x}, t)=\displaystyle\sum_{j=1}^{\infty}\frac{\Delta t^{j}}{j!}D_{t}^{j}f_{i}= -\left( \frac{ f_{i} - f_{i}^{eq} }{ \tau } \right).
\end{split}
\end{equation}
\sphinxAtStartPar
Rescaling the dimensionless form of the Eq. \eqref{equation:LBM-Study/Poisson-Equation:EqExp-Poisson-Eq} in terms of the Knudsen number (\(Kn\)), we have
\begin{equation*}
\begin{split}
\displaystyle \sum^{\infty}_{j=1}  \frac{Kn^{(j-1)}}{j!} D_{t}^{j} f_{i} = - \frac{1}{Kn}\left( \frac{ f_{i} - f_{eq,i} }{ \tau } \right) \quad \quad \rightarrow \quad \quad \displaystyle \sum^{\infty}_{j=1}  \frac{Kn^{(j)}}{j!} D_{t}^{j} f_{i} = -  \frac{ f_{i} - f_{i}^{eq} }{ \tau },
\end{split}
\end{equation*}
\sphinxAtStartPar
applying the asymptotic expansion in both the distribution function (\(f_{i}=f_{i}^{(0)}+Kn f_{i}^{(1)} + Kn^{2} f_{i}^{(2)}+\cdots\)) and time partial derivative (\(\partial_{t}=\partial_{t}^{(0)}+ Kn \partial_{t}^{(1)}+Kn^{2} \partial_{t}^{(2)}+\cdots\)), and separating the equation in orders up to the order \(Kn^{2}\):
\begin{equation}\label{equation:LBM-Study/Poisson-Equation:Chap-Kn-Poisson-Eq}
\begin{split}
\begin{array}{ll}
    (Kn^{(0)}):& f_{i}^{(0)}= f_{i}^{eq},\\
    (Kn^{(1)}):& \left( \partial_{t}^{(0)} + c_{i,\alpha}\partial_{\alpha} \right)f_{i}^{(0)}= - \displaystyle\frac{ f_{i}^{(1)} }{ \tau } , \\
    (Kn^{(2)}):& \partial_{t}^{(1)} f_{i}^{(0)} + \left( \partial_{t}^{(0)} + c_{i,\alpha}\partial_{\alpha} \right)f_{i}^{(0)} + \displaystyle\frac{\left( \partial_{t}^{(0)} + c_{i,\alpha}\partial_{\alpha} \right)^{2} f_{i}^{(0)}}{2}= - \displaystyle\frac{ f_{i}^{(2)} }{ \tau },\\
    \textrm{ou}& \\
    (Kn^{(2)}):& \partial_{t}^{(1)} f_{i}^{(0)} + \left( \partial_{t}^{(0)} + c_{i,\alpha}\partial_{\alpha} \right)\displaystyle\left(1 - \frac{1}{2\tau}\right) f_{i}^{(1)} =  - \displaystyle\frac{ f_{i}^{(2)} }{ \tau } \quad\quad \rightarrow \quad\quad  f_{i}^{(1)}\textrm{ Formulation} \\
    \textrm{ou}& \\
    (Kn^{(2)}):& \partial_{t}^{(1)} f_{i}^{(0)} + \left( \partial_{t}^{(0)} + c_{i,\alpha}\partial_{\alpha} \right)^{2} \displaystyle\left(\frac{1}{2} - \tau\right) f_{i}^{(0)} =  - \displaystyle\frac{ f_{i}^{(2)} }{ \tau } \quad\quad \rightarrow  \quad\quad f_{i}^{(0)}\textrm{ Formulation}.
 \end{array}
\end{split}
\end{equation}
\sphinxAtStartPar
\sphinxstylestrong{Zero\sphinxhyphen{}Order Moment Balance}

\sphinxAtStartPar
To retrieve the balance equation, we sum the Eq. \eqref{equation:LBM-Study/Poisson-Equation:Chap-Kn-Poisson-Eq} for \(Kn^{(1)}\) and \(Kn^{(2)}\) over \(\sum_{i=0} \):
\begin{equation*}
\begin{split}
\begin{array}{l}
(Kn^{(1)}): \displaystyle\sum_{i}\left( \partial_{t}^{(0)} + c_{i,\alpha}\partial_{\alpha} \right)f_{i}^{(0)}= \displaystyle\sum_{i} \left(- \displaystyle\frac{ f_{i}^{(1)} }{ \tau } \right),\\
(Kn^{(1)}): 0  =0,
\end{array}
\end{split}
\end{equation*}
\sphinxAtStartPar
and
\begin{equation*}
\begin{split}
\begin{array}{l}
(Kn^{(2)}): \displaystyle\sum_{i}\left( \partial_{t}^{(1)} f_{i}^{(0)} + \left( \partial_{t}^{(0)} + c_{i,\alpha}\partial_{\alpha} \right)^{2} \displaystyle\left( \frac{1}{2}-\tau\right) f_{i}^{(0)} \right) =  \displaystyle\sum_{i}\left( - \displaystyle\frac{ f_{i}^{(2)} }{ \tau } \right), \\
(Kn^{(2)}): \left( \displaystyle\frac{1}{2}-\tau\right) \partial_{\alpha}\partial_{\beta} \left(c_{s}^{2} \widehat{p} \delta_{\alpha\beta} \right) =0. 
\end{array}
\end{split}
\end{equation*}
\sphinxAtStartPar
or
\begin{equation}\label{equation:LBM-Study/Poisson-Equation:Kn1-0-B0}
\begin{split}
\begin{array}{ll}
(Kn^{(2)}):& \displaystyle\sum_{i}\left[ \partial_{t}^{(1)} f_{i}^{(0)} - \left(1-\displaystyle\frac{1}{2\tau}\right)\left( \partial_{t}^{(0)} + c_{i,\alpha}\partial_{\alpha} \right) f_{i}^{(1)}   \right]=  \displaystyle\sum_{i}\left( - \displaystyle\frac{f_{i}^{(2)}}{\tau} \right), \nonumber \\
(Kn^{(2)}):& - \left(1-\displaystyle\frac{1}{2\tau}\right) m^{(1)}_{\alpha} = 0.
\end{array}
\end{split}
\end{equation}
\sphinxAtStartPar
where \(m^{(1)}_{\alpha}\) is the first\sphinxhyphen{}order moment of \(f_{i}^{(1)}\) and \(1/(1-w_{0})\sum_{i=1}f_{i}^{(j)}=0\) for \(j\geq 1\) due to imposition of \(\widehat{p}\) conservation. To compute the moment \(m^{(1)}_{\alpha}\), we multiply Eq. \eqref{equation:LBM-Study/Poisson-Equation:Chap-Kn-Poisson-Eq} by \(c_{i,\alpha}\) and sum over all lattice directions:
\begin{equation}\label{equation:LBM-Study/Poisson-Equation:Kn0-B1}
\begin{split}
\begin{array}{l}
(Kn^{(1)}): \displaystyle\sum_{i}c_{i,\alpha}\left( \partial_{t}^{(0)} + e_{\beta,i}\partial_{\beta} \right)f_{i}^{(0)}= \displaystyle\sum_{i} c_{i,\alpha} \left( -\frac{f_{i}^{(1)}}{\tau}  \right),\\
(Kn^{(1)}): m^{(1)}_{\alpha} =-\tau \partial_{\beta}\left( c_{s}^{2} \widehat{p} \delta_{\alpha\beta}  \right).
\end{array}
\end{split}
\end{equation}
\sphinxAtStartPar
By substituting Eq. \eqref{equation:LBM-Study/Poisson-Equation:Kn0-B1} into Eq. \eqref{equation:LBM-Study/Poisson-Equation:Kn1-0-B0}, we recover the Eq. \eqref{equation:LBM-Study/Poisson-Equation:poisson-eq-lb}, where \(\nu=c_{s}^{2}(\tau-1/2)\). In the regularized formulation of the lattice Boltzmann equation, the first\sphinxhyphen{}order correction is approximated as \(f_{i}^{(1)}\approx (f_{i}-f_{i}^{(0)})=f_{i}^{neq}\), where higher\sphinxhyphen{}order Knudsen moments are filtered out and the collision term is reformulated in terms of \(f_{i}^{neq}\).

\end{sphinxuseclass}

\subsubsection{Boundary Conditions}
\label{\detokenize{LBM-Study/Poisson-Equation:boundary-conditions}}
\sphinxAtStartPar
The boundary conditions for the lattice D2Q5 can be derived by solving a linear system of known moments, as illustrated in \DUrole{xref,std,std-doc}{appendix/bcs\sphinxhyphen{}D2Q5\sphinxhyphen{}poisson\sphinxhyphen{}equation}. For the D2Q5 lattice, the unknown distribution functions at each boundary are determined as summarized in Table below.


\begin{savenotes}\sphinxattablestart
\sphinxthistablewithglobalstyle
\centering
\begin{tabulary}{\linewidth}[t]{TTTTT}
\sphinxtoprule
\sphinxstyletheadfamily 
\sphinxAtStartPar
Boundaries
&
\sphinxAtStartPar

&\sphinxstyletheadfamily 
\sphinxAtStartPar
Layers
&
\sphinxAtStartPar

&
\sphinxAtStartPar

\\
\sphinxmidrule
\sphinxtableatstartofbodyhook
\sphinxAtStartPar

&
\sphinxAtStartPar

&
\sphinxAtStartPar
North
&
\sphinxAtStartPar

&
\sphinxAtStartPar
South
\\
\sphinxhline
\sphinxAtStartPar
Unknown \(f_{i}\)
&
\sphinxAtStartPar

&
\sphinxAtStartPar
\(f_2=\widehat{p}_{n}(1-w_{0}) - (f_1+f_3+f_4)\)
&
\sphinxAtStartPar

&
\sphinxAtStartPar
\(f_4=\widehat{p}_{s}(1-w_{0}) - (f_1+f_2+f_3)\)
\\
\sphinxhline
\sphinxAtStartPar

&
\sphinxAtStartPar

&
\sphinxAtStartPar
East
&
\sphinxAtStartPar

&
\sphinxAtStartPar
West
\\
\sphinxhline
\sphinxAtStartPar

&
\sphinxAtStartPar

&
\sphinxAtStartPar
\(f_3=\widehat{p}_{e}(1-w_{0}) - (f_1+f_2+f_4)\)
&
\sphinxAtStartPar

&
\sphinxAtStartPar
\(f_1=\widehat{p}_{w}(1-w_{0}) - (f_2+f_3+f_4)\)
\\
\sphinxhline
\sphinxAtStartPar
Boundaries
&
\sphinxAtStartPar

&
\sphinxAtStartPar
Concave Corners
&
\sphinxAtStartPar

&
\sphinxAtStartPar

\\
\sphinxhline
\sphinxAtStartPar

&
\sphinxAtStartPar

&
\sphinxAtStartPar
Northwest
&
\sphinxAtStartPar

&
\sphinxAtStartPar
Northeast
\\
\sphinxhline
\sphinxAtStartPar
Unknown \(f_{i}\)
&
\sphinxAtStartPar

&
\sphinxAtStartPar
\(f_1=-(\partial_{x}\widehat{p}+\partial_{y}\widehat{p})\tau c_{s}^{2}/2 - f_2 + \widehat{p}_{nw}/3\)
&
\sphinxAtStartPar

&
\sphinxAtStartPar
\(f_3=(\partial_{x}\widehat{p}-\partial_{y}\widehat{p})\tau c_{s}^{2}/2 - f_2 + \widehat{p}_{ne}/3\)
\\
\sphinxhline
\sphinxAtStartPar

&
\sphinxAtStartPar

&
\sphinxAtStartPar
\(f_4=(\partial_{x}\widehat{p}+\partial_{y}\widehat{p})\tau c_{s}^{2}/2 - f_3 + \widehat{p}_{nw}/3\)
&
\sphinxAtStartPar

&
\sphinxAtStartPar
\(f_4=-(\partial_{x}\widehat{p}-\partial_{y}\widehat{p})\tau c_{s}^{2}/2 - f_1 + \widehat{p}_{ne}/3\)
\\
\sphinxhline
\sphinxAtStartPar

&
\sphinxAtStartPar

&
\sphinxAtStartPar
Southwest
&
\sphinxAtStartPar

&
\sphinxAtStartPar
Southeast
\\
\sphinxhline
\sphinxAtStartPar

&
\sphinxAtStartPar

&
\sphinxAtStartPar
\(f_1=-(\partial_{x}\widehat{p}-\partial_{y}\widehat{p})\tau c_{s}^{2}/2 - f_4 + \widehat{p}_{sw}/3\)
&
\sphinxAtStartPar

&
\sphinxAtStartPar
\(f_2=-(\partial_{x}\widehat{p}+\partial_{y}\widehat{p})\tau c_{s}^{2}/2 - f_1 + \widehat{p}_{se}/3\)
\\
\sphinxhline
\sphinxAtStartPar

&
\sphinxAtStartPar

&
\sphinxAtStartPar
\(f_2=(\partial_{x}\widehat{p}-\partial_{y}\widehat{p})\tau c_{s}^{2}/2 - f_3 + \widehat{p}_{sw}/3\)
&
\sphinxAtStartPar

&
\sphinxAtStartPar
\(f_3=(\partial_{x}\widehat{p}+\partial_{y}\widehat{p})\tau c_{s}^{2}/2 - f_4 + \widehat{p}_{se}/3\)
\\
\sphinxbottomrule
\end{tabulary}
\sphinxtableafterendhook\par
\sphinxattableend\end{savenotes}


\section{1st Benchmark: Single sine mode on a finite rod (Dirichlet)}
\label{\detokenize{LBM-Study/Poisson-Equation:st-benchmark-single-sine-mode-on-a-finite-rod-dirichlet}}
\sphinxAtStartPar
PDE: \((\phi_t = \partial_{x}(\nu\partial_{x}\phi), (0<x<L,\ t>0)\)

\sphinxAtStartPar
BC: \((\phi(0,t)=\phi(L,t)=0)\)

\sphinxAtStartPar
IC: \(\phi(x,0)=\sin \left(\frac{\pi x}{L}\right)\)

\sphinxAtStartPar
\sphinxstylestrong{Solution}
\begin{equation*}
\begin{split}
\phi(x,t)=\sin\left(\frac{\pi x}{L}\right) \exp \left(-\nu \frac{\pi^2}{L^2} t\right).
\end{split}
\end{equation*}

\subsection{FDM Solution:}
\label{\detokenize{LBM-Study/Poisson-Equation:fdm-solution}}
\begin{sphinxuseclass}{cell}
\begin{sphinxuseclass}{tag_hide-input}
\begin{sphinxuseclass}{tag_hide-output}
\end{sphinxuseclass}
\end{sphinxuseclass}
\end{sphinxuseclass}

\subsection{LBM D1Q3 Solution With Steady\sphinxhyphen{}Scheme:}
\label{\detokenize{LBM-Study/Poisson-Equation:lbm-d1q3-solution-with-steady-scheme}}
\begin{sphinxuseclass}{cell}
\begin{sphinxuseclass}{tag_hide-input}
\begin{sphinxuseclass}{tag_hide-output}
\end{sphinxuseclass}
\end{sphinxuseclass}
\end{sphinxuseclass}

\subsection{LBM D1Q2 Solution (Not Accept the Steady\sphinxhyphen{}Scheme):}
\label{\detokenize{LBM-Study/Poisson-Equation:lbm-d1q2-solution-not-accept-the-steady-scheme}}
\begin{sphinxuseclass}{cell}
\begin{sphinxuseclass}{tag_hide-input}
\begin{sphinxuseclass}{tag_hide-output}
\end{sphinxuseclass}
\end{sphinxuseclass}
\end{sphinxuseclass}

\subsection{LBM D1Q3 Solution (Transient):}
\label{\detokenize{LBM-Study/Poisson-Equation:lbm-d1q3-solution-transient}}
\begin{sphinxuseclass}{cell}
\begin{sphinxuseclass}{tag_hide-input}
\begin{sphinxuseclass}{tag_hide-output}
\end{sphinxuseclass}
\end{sphinxuseclass}
\end{sphinxuseclass}

\subsection{Numerical and Analytical Comparisson:}
\label{\detokenize{LBM-Study/Poisson-Equation:numerical-and-analytical-comparisson}}
\begin{sphinxuseclass}{cell}
\begin{sphinxuseclass}{tag_hide-input}\begin{sphinxVerbatimOutput}

\begin{sphinxuseclass}{cell_output}
\noindent\sphinxincludegraphics{{a26a8ff05423c3e044d97ad61a5b25eb5cee398256b7ea91e983801bf49fecf7}.png}

\end{sphinxuseclass}\end{sphinxVerbatimOutput}

\end{sphinxuseclass}
\end{sphinxuseclass}

\section{2nd Benchmark: 2D Poisson Pressure Field}
\label{\detokenize{LBM-Study/Poisson-Equation:nd-benchmark-2d-poisson-pressure-field}}
\sphinxAtStartPar
Analyzing the accuracy of Poisson equation solver, we apply it to a two\sphinxhyphen{}dimensional analytical case. The geometry is described by a square domain of length \(L\) initialized with a constant dimensionless pressure \(\widehat{p}(x,y,t=0)=0\). The boundary conditions of the geometry are given by
\begin{equation*}
\begin{split}
\widehat{p}(x=0,y,t)=\displaystyle\frac{sinh\left(\pi(1-y)\right)}{sinh\left(\pi\right)}, \quad \quad \widehat{p}(x=L,y,t)= -\frac{sinh\left(\pi(1-y)\right)}{sinh\left(\pi\right)},
\end{split}
\end{equation*}\begin{equation*}
\begin{split}
\widehat{p}(x,y=0,t)=cos(\pi x), \quad \quad \widehat{p}(x,y=L,t)=0.
\end{split}
\end{equation*}
\sphinxAtStartPar
The pressure analytical solution for Eq. \eqref{equation:LBM-Study/Poisson-Equation:poisson-eq-lb} pressure is given
\begin{equation*}
\begin{split}
    \widehat{p}(x,y,t)=cos(\pi x) \displaystyle\frac{sinh\left(\pi(1-y)\right)}{sinh\left(\pi\right)}
\end{split}
\end{equation*}
\sphinxAtStartPar
where \(Dp_{0}\) is considered equal to 1, consequently, velocity values can be analytically obtained by the equation:
\begin{equation*}
\begin{split}
    u_{x}=D\partial_{x}p= \pi sin(\pi x) \displaystyle\frac{sinh\left(\pi(1-y)\right)}{sinh\left(\pi\right)}, \quad \quad u_{y}=D\partial_{y}p= \pi cos(\pi x) \displaystyle\frac{cosh\left(\pi(1-y)\right)}{sinh\left(\pi\right)}.
\end{split}
\end{equation*}

\subsection{LBM D2Q5 Code}
\label{\detokenize{LBM-Study/Poisson-Equation:lbm-d2q5-code}}
\begin{sphinxuseclass}{cell}
\begin{sphinxuseclass}{tag_skip-execution}
\begin{sphinxuseclass}{tag_hide-input}
\begin{sphinxuseclass}{tag_hide-output}
\end{sphinxuseclass}
\end{sphinxuseclass}
\end{sphinxuseclass}
\end{sphinxuseclass}
\begin{sphinxuseclass}{cell}
\begin{sphinxuseclass}{tag_skip-execution}
\begin{sphinxuseclass}{tag_hide-input}
\end{sphinxuseclass}
\end{sphinxuseclass}
\end{sphinxuseclass}
\begin{sphinxuseclass}{cell}
\begin{sphinxuseclass}{tag_skip-execution}
\begin{sphinxuseclass}{tag_hide-input}
\begin{sphinxuseclass}{tag_hide-output}
\end{sphinxuseclass}
\end{sphinxuseclass}
\end{sphinxuseclass}
\end{sphinxuseclass}
\begin{sphinxuseclass}{cell}
\begin{sphinxuseclass}{tag_skip-execution}
\begin{sphinxuseclass}{tag_hide-input}\begin{sphinxVerbatimOutput}

\begin{sphinxuseclass}{cell_output}
\noindent\sphinxincludegraphics{{a5e85de83974bd58d89d09327c8c61612ad7b8ccfe4314bd01e99be1bf2255f3}.png}

\end{sphinxuseclass}\end{sphinxVerbatimOutput}

\end{sphinxuseclass}
\end{sphinxuseclass}
\end{sphinxuseclass}
\begin{sphinxuseclass}{cell}
\begin{sphinxuseclass}{tag_skip-execution}
\begin{sphinxuseclass}{tag_hide-input}
\end{sphinxuseclass}
\end{sphinxuseclass}
\end{sphinxuseclass}
\begin{sphinxuseclass}{cell}
\begin{sphinxuseclass}{tag_skip-execution}
\begin{sphinxuseclass}{tag_hide-input}\begin{sphinxVerbatimOutput}

\begin{sphinxuseclass}{cell_output}
\noindent\sphinxincludegraphics{{a5d85b5c52d191d4ab06164ace1f069946d3dab1cf0de340584764cba31ea8ef}.png}

\end{sphinxuseclass}\end{sphinxVerbatimOutput}

\end{sphinxuseclass}
\end{sphinxuseclass}
\end{sphinxuseclass}
\begin{sphinxuseclass}{cell}
\begin{sphinxuseclass}{tag_hide-input}
\begin{sphinxuseclass}{tag_skip-execution}\begin{sphinxVerbatimOutput}

\begin{sphinxuseclass}{cell_output}
\noindent\sphinxincludegraphics{{7e80630a18555e5cb98b6a6f3f5b2318741369b23d32c5053499d2207f72380c}.png}

\end{sphinxuseclass}\end{sphinxVerbatimOutput}

\end{sphinxuseclass}
\end{sphinxuseclass}
\end{sphinxuseclass}
\begin{sphinxuseclass}{cell}
\begin{sphinxuseclass}{tag_skip-execution}
\begin{sphinxuseclass}{tag_hide-input}
\begin{sphinxuseclass}{tag_hide-output}
\end{sphinxuseclass}
\end{sphinxuseclass}
\end{sphinxuseclass}
\end{sphinxuseclass}
\begin{sphinxuseclass}{cell}
\begin{sphinxuseclass}{tag_skip-execution}
\begin{sphinxuseclass}{tag_hide-input}
\end{sphinxuseclass}
\end{sphinxuseclass}
\end{sphinxuseclass}
\begin{sphinxuseclass}{cell}
\begin{sphinxuseclass}{tag_skip-execution}
\begin{sphinxuseclass}{tag_hide-input}
\begin{sphinxuseclass}{tag_hide-output}
\end{sphinxuseclass}
\end{sphinxuseclass}
\end{sphinxuseclass}
\end{sphinxuseclass}
\begin{sphinxuseclass}{cell}
\begin{sphinxuseclass}{tag_skip-execution}
\begin{sphinxuseclass}{tag_hide-input}
\begin{sphinxuseclass}{tag_hide-output}
\end{sphinxuseclass}
\end{sphinxuseclass}
\end{sphinxuseclass}
\end{sphinxuseclass}
\begin{sphinxuseclass}{cell}
\begin{sphinxuseclass}{tag_skip-execution}
\begin{sphinxuseclass}{tag_hide-input}
\begin{sphinxuseclass}{tag_hide-output}
\end{sphinxuseclass}
\end{sphinxuseclass}
\end{sphinxuseclass}
\end{sphinxuseclass}
\begin{sphinxuseclass}{cell}
\begin{sphinxuseclass}{tag_skip-execution}
\begin{sphinxuseclass}{tag_hide-input}\begin{sphinxVerbatimOutput}

\begin{sphinxuseclass}{cell_output}
\noindent\sphinxincludegraphics{{cc6a0c3bf3366904a5cdd8aa74d750096a61b6db71c78ad537ee4396c67d7f05}.png}

\end{sphinxuseclass}\end{sphinxVerbatimOutput}

\end{sphinxuseclass}
\end{sphinxuseclass}
\end{sphinxuseclass}
\sphinxstepscope


\chapter{Convective Equation (Transport Equation)}
\label{\detokenize{LBM-Study/Convective-Equation:convective-equation-transport-equation}}\label{\detokenize{LBM-Study/Convective-Equation::doc}}
\sphinxAtStartPar
The convective equation also known as transport equation is a partial diferential equation (PDE) usually described by
\begin{equation}\label{equation:LBM-Study/Convective-Equation:cv-eq-lb}
\begin{split}
\partial_t u + \partial_{\alpha}B_{\alpha}  = 0
\end{split}
\end{equation}
\sphinxAtStartPar
where \(u(\alpha,t)\) is a time and space dependent parameter that is being transporte and \(B_{\alpha}(\alpha,t)\) is a second time and space dependent parameter that can also be depende of \(u(\alpha,t)\) in non\sphinxhyphen{}linear convective diffusive equations.


\section{Numerical Discretization}
\label{\detokenize{LBM-Study/Convective-Equation:numerical-discretization}}
\sphinxAtStartPar
Discretizing the above equation in regular grid of 1D domain, we have

\begin{figure}[htbp]
\centering
\capstart

\noindent\sphinxincludegraphics[scale=2.5]{{mesh-1-convective}.png}
\caption{1D regular grid.}\label{\detokenize{LBM-Study/Convective-Equation:mesh-1-convective}}\end{figure}


\subsection{FDM Discretization}
\label{\detokenize{LBM-Study/Convective-Equation:fdm-discretization}}
\sphinxAtStartPar
The Discretization through finite difference method can assume distinct forms that describe the problem with different accuracy, rate convergence and stability. In the sequence of this section, we will explore different variation that can work with the convective formulation.


\subsubsection{Forward\sphinxhyphen{}Time and Forward\sphinxhyphen{}Space (FTFS)}
\label{\detokenize{LBM-Study/Convective-Equation:forward-time-and-forward-space-ftfs}}
\sphinxAtStartPar
Discretizing the Eq. \eqref{equation:LBM-Study/Convective-Equation:cv-eq-lb} through finite diference method, we have for forward\sphinxhyphen{}time and forward\sphinxhyphen{}space (FTFS) discretizations:
\begin{equation*}
\begin{split}
\partial_t u + \partial_{x}B_{\alpha}  = \frac{u(x,t + \Delta_{t}) - u(x,t) }{\Delta_{t}} + \frac{B_{\alpha}(x,t) - B_{\alpha}(x + \Delta_{x},t) }{\Delta_{x}} =0,
\end{split}
\end{equation*}
\sphinxAtStartPar
isolating the term \(u(x,t + \Delta_{t})\), we can describe the \(u\) time evolution:
\begin{equation*}
\begin{split}
u(x,t + \Delta_{t})   = u(x,t) - \frac{\Delta_{t}}{\Delta_{x}} \left( B_{\alpha}(x,t) - B_{\alpha}(x + \Delta_{x},t) \right) .
\end{split}
\end{equation*}

\subsubsection{Forward\sphinxhyphen{}Time and Central\sphinxhyphen{}Space with Small Diffusivity (FTCS\sphinxhyphen{}SD)}
\label{\detokenize{LBM-Study/Convective-Equation:forward-time-and-central-space-with-small-diffusivity-ftcs-sd}}
\sphinxAtStartPar
Give that discretize convective equation with pure forwar\sphinxhyphen{}time and central\sphinxhyphen{}space (FTCS) can be unstable for several cases, employ a forwar\sphinxhyphen{}time and central\sphinxhyphen{}space with small diffusivity (FTCS\sphinxhyphen{}SD) is a alternative to improve stability and ensure a propertie conservation.
\begin{equation*}
\begin{split}
\begin{array}{ll}
\partial_t u + \partial_{x}B_{\alpha} - \partial_{\alpha}\left( \nu \partial_{\alpha}u \right) = & \displaystyle\frac{u(x,t + \Delta_{t}) - u(x,t) }{\Delta_{t}} + \frac{B_{\alpha}(x+\Delta_{x},t) - B_{\alpha}(x - \Delta_{x},t) }{2\Delta_{x}} \\ & - \nu \displaystyle\frac{ u(x - \Delta_{x},t) -2u(x,t) + u(x + \Delta_{x},t) }{\Delta_{x}^{2}} =0,
\end{array}
\end{split}
\end{equation*}
\sphinxAtStartPar
isolating the term \(u(x,t + \Delta_{t})\), we can describe the \(u\) time evolution:
\begin{equation*}
\begin{split}
u(x,t + \Delta_{t})   = u(x,t) + \frac{\Delta_{t}}{2\Delta_{x}} \left( B_{\alpha}(x-\Delta_{x},t) - B_{\alpha}(x + \Delta_{x},t) \right) + \frac{\nu\Delta_{t}}{\Delta_{x}^{2}} \left( u(x - \Delta_{x},t) -2u(x,t) + u(x + \Delta_{x},t) \right) .
\end{split}
\end{equation*}

\subsection{Lattice Boltzmann Equation}
\label{\detokenize{LBM-Study/Convective-Equation:lattice-boltzmann-equation}}
\sphinxAtStartPar
Describing the problem through the BGK lattice Boltzmann equation, we assume a equilibrium state for the present problem, where \( f_{i}= f_{i}^{eq}\) and the BGK lattice Boltzmann equation reduces to:
\begin{equation}\label{equation:LBM-Study/Convective-Equation:LB-cv-Eq}
\begin{split}
f_{i}\left(x_{\alpha} + c_{i,\alpha} \Delta t, t+\Delta t\right)- f_{i}(x_{\alpha}, t)= -\left( \frac{ f_{i} - f_{i}^{eq} }{ \tau } \right) \quad \quad \rightarrow \quad \quad f_i( x_{\alpha} + c_{i,\alpha} \delta t, t+\delta t) = f_{i}^{eq} , 
\end{split}
\end{equation}
\sphinxAtStartPar
where the equilibrium distribution function is defined by
\begin{equation*}
\begin{split}
f^{eq}_{i} = w_{i}\left( u + \frac{c_{i,\alpha} B_{\alpha}}{c_{s}^{2}} \right).
\end{split}
\end{equation*}

\subsubsection{Lattice Direction Moments}
\label{\detokenize{LBM-Study/Convective-Equation:lattice-direction-moments}}
\begin{sphinxuseclass}{cell}
\begin{sphinxuseclass}{tag_hide-input}
\end{sphinxuseclass}
\end{sphinxuseclass}

\paragraph{Raw Moments}
\label{\detokenize{LBM-Study/Convective-Equation:raw-moments}}
\begin{sphinxuseclass}{cell}
\begin{sphinxuseclass}{tag_hide-input}\begin{sphinxVerbatimOutput}

\begin{sphinxuseclass}{cell_output}\begin{equation*}
\begin{split}\displaystyle \begin{gathered}\underbrace{\sum_{i=0} f_{i}^{eq} =\sum_{i=0} f_{i} }_{\textrm{Zero-Order Moment}} =u,\quad \quad \underbrace{\sum_{i=0} f_{i}^{eq}c_{i,x} }_{\textrm{x-First-Order Moment}} =B_{\alpha},\quad \quad \underbrace{\sum_{i=0} f_{i}^{eq} c_{i,x}c_{i,x} }_{\textrm{xx-Second-Order Moment}} =\frac{u}{3},\quad \quad \underbrace{\sum_{i=0} f_{i}^{eq} c_{i,x}c_{i,x}c_{i,x} }_{\textrm{xxx-Third-Order Moment}} =B_{\alpha}\end{gathered}\end{split}
\end{equation*}
\end{sphinxuseclass}\end{sphinxVerbatimOutput}

\end{sphinxuseclass}
\end{sphinxuseclass}

\paragraph{Hermite Moments}
\label{\detokenize{LBM-Study/Convective-Equation:hermite-moments}}
\begin{sphinxuseclass}{cell}
\begin{sphinxuseclass}{tag_hide-input}\begin{sphinxVerbatimOutput}

\begin{sphinxuseclass}{cell_output}\begin{equation*}
\begin{split}\displaystyle \begin{gathered}\underbrace{\sum_{i=0} f_{i}^{eq} =\sum_{i=0} f_{i} }_{\textrm{Zero-Order Hermite Moment}} =u,\quad \quad \underbrace{\sum_{i=0} f_{i}^{eq}c_{i,x} }_{\textrm{x-First-Order Hermite Moment}} =B_{\alpha}\end{gathered}\end{split}
\end{equation*}\begin{equation*}
\begin{split}\displaystyle \begin{gathered}\underbrace{\sum_{i=0} f_{i}^{eq} \left(c_{i,x}c_{i,x} - c_{s}^{2}\right) }_{\textrm{xx-Second-Order Hermite Moment}} =0,\quad \quad \underbrace{\sum_{i=0} f_{i}^{eq} \left(c_{i,x}c_{i,x} - 3c_{s}^{2} \right) c_{i,x} }_{\textrm{xxx-Third-Order Hermite Moment}} =0\end{gathered}\end{split}
\end{equation*}
\end{sphinxuseclass}\end{sphinxVerbatimOutput}

\end{sphinxuseclass}
\end{sphinxuseclass}

\subsubsection{Chapmann\sphinxhyphen{}Enskog Analysis (Considering Equilibrium State \sphinxhyphen{} Consequently Disregarding Inherited Diffusive Effects)}
\label{\detokenize{LBM-Study/Convective-Equation:chapmann-enskog-analysis-considering-equilibrium-state-consequently-disregarding-inherited-diffusive-effects}}
\begin{sphinxuseclass}{toggle}
\sphinxAtStartPar
Applying the Chapmann\sphinxhyphen{}Enskog procedure to LB equation, we expand the term \(f_{i}\left(\boldsymbol{x}+\boldsymbol{e}_i \Delta t, t+\Delta t\right)=f_{i}^{'}\) in a Taylor series to recover the derivative form of the equation, i.e.,
\begin{equation}\label{equation:LBM-Study/Convective-Equation:EqExp-cv-Eq}
\begin{split}
f_{i}^{'} - f_{i}(\boldsymbol{x}, t)=\displaystyle\sum_{j=1}^{\infty}\frac{\Delta t^{j}}{j!}D_{t}^{j}f_{i}= -\left( \frac{ f_{i} - f_{i}^{eq} }{ \tau } \right) \quad \quad \overbrace{\rightarrow}^{\textrm{Equilibrium State}} \quad \quad f_{i}^{'} = \displaystyle\sum_{j=0}^{\infty}\frac{\Delta t^{j}}{j!}D_{t}^{j}f_{i}= f_{i}^{eq}(\boldsymbol{x}, t).
\end{split}
\end{equation}
\sphinxAtStartPar
Rescaling the dimensionless form of the Eq. \eqref{equation:LBM-Study/Convective-Equation:EqExp-cv-Eq} in terms of the Knudsen number (\(Kn\)), we have
\begin{equation*}
\begin{split}
\displaystyle \sum^{\infty}_{j=0}  \frac{Kn^{(j-1)}}{j!} D_{t}^{j} f_{i} = f_{i}^{eq}(x_{\alpha}, t),
\end{split}
\end{equation*}
\sphinxAtStartPar
applying the asymptotic expansion in both the distribution function (\(f_{i}=f_{i}^{(0)}+Kn f_{i}^{(1)} + Kn^{2} f_{i}^{(2)}+\cdots\)) and time partial derivative (\(\partial_{t}=\partial_{t}^{(0)}+ Kn \partial_{t}^{(1)}+Kn^{2} \partial_{t}^{(2)}+\cdots\)), and separating the equation in orders up to the order \(Kn^{2}\):
\begin{equation}\label{equation:LBM-Study/Convective-Equation:Chap-Kn-cv-Eq}
\begin{split}
\begin{array}{ll}
    (Kn^{(0)}):& f_{i}^{(0)}= f_{i}^{eq},\\
    (Kn^{(1)}):& \left( \partial_{t}^{(0)} + c_{i,\alpha}\partial_{\alpha} \right)f_{i}^{(0)}= 0, \\
    (Kn^{(2)}):& \partial_{t}^{(1)} f_{i}^{(0)} + \left( \partial_{t}^{(0)} + c_{i,\alpha}\partial_{\alpha} \right)f_{i}^{(1)} + \displaystyle\frac{\left( \partial_{t}^{(0)} + c_{i,\alpha}\partial_{\alpha} \right)^{2} f_{i}^{(0)}}{2}= 0,\\
    \textrm{ou} &\\
    (Kn^{(2)}):& \partial_{t}^{(1)} f_{i}^{(0)} + \left( \partial_{t}^{(0)} + c_{i,\alpha}\partial_{\alpha} \right)f_{i}^{(1)} = 0,\\
    (Kn^{(3)}):& \partial_{t}^{(2)} f_{i}^{(0)} +\partial_{t}^{(1)} f_{i}^{(1)} + \partial_{t}^{(1)}\left( \left( \partial_{t}^{(0)} + c_{i,\alpha}\partial_{\alpha} \right)f_{i}^{(0)} \right) + \left( \partial_{t}^{(0)} + c_{i,\alpha}\partial_{\alpha} \right)f_{i}^{(2)} + \displaystyle\frac{\left( \partial_{t}^{(0)} + c_{i,\alpha}\partial_{\alpha} \right)^{2} f_{i}^{(1)}}{2} \\ 
    & +  \displaystyle\frac{\left( \partial_{t}^{(0)} + c_{i,\alpha}\partial_{\alpha} \right)^{3} f_{i}^{(0)}}{6} =  0,  \\
    \textrm{ou} &\\
    (Kn^{(3)}):& \partial_{t}^{(2)} f_{i}^{(0)} +\partial_{t}^{(1)} f_{i}^{(1)} + \left( \partial_{t}^{(0)} + c_{i,\alpha}\partial_{\alpha} \right)f_{i}^{(2)} =  0.  \\
 \end{array}
\end{split}
\end{equation}
\sphinxAtStartPar
\sphinxstylestrong{Zero\sphinxhyphen{}Order Moment Balance}

\sphinxAtStartPar
To retrieve the balance equation, we sum the Eq. \eqref{equation:LBM-Study/Convective-Equation:Chap-Kn-cv-Eq} for \(Kn^{(1)}\) and \(Kn^{(2)}\) over \(\sum_{i=0} \):
\begin{equation}\label{equation:LBM-Study/Convective-Equation:cv-Kn1-B0}
\begin{split}
\begin{array}{l}
(Kn^{(1)}): \displaystyle\sum_{i}\left( \partial_{t}^{(0)} + c_{i,\alpha}\partial_{\alpha} \right)f_{i}^{(0)}= 0,\\
(Kn^{(1)}): \partial_{t}^{(0)}  u + \partial_{\alpha} B_{\alpha} =0,
\end{array}
\end{split}
\end{equation}\begin{equation}\label{equation:LBM-Study/Convective-Equation:cv-Kn2-B0}
\begin{split}
\begin{array}{ll}
(Kn^{(2)}):& \displaystyle\sum_{i}\left[ \partial_{t}^{(1)} f_{i}^{(0)} + \left( \partial_{t}^{(0)} + c_{i,\alpha}\partial_{\alpha} \right) f_{i}^{(1)}   \right]=  0, \\
(Kn^{(2)}):& \partial_{t}^{(1)} u + \partial_{\alpha} m^{(1)}_{\alpha} = 0,
\end{array}
\end{split}
\end{equation}
\sphinxAtStartPar
and
\begin{equation}\label{equation:LBM-Study/Convective-Equation:cv-Kn3-B0}
\begin{split}
\begin{array}{ll}
(Kn^{(3)}):& \displaystyle\sum_{i}\left[ \partial_{t}^{(2)} f_{i}^{(0)} +\partial_{t}^{(1)} f_{i}^{(1)} + \left( \partial_{t}^{(0)} + c_{i,\alpha}\partial_{\alpha} \right)f_{i}^{(2)}   \right]=  0, \\
(Kn^{(3)}):& \partial_{t}^{(2)} u + \partial_{\alpha} m^{(2)}_{\alpha} = 0.
\end{array}
\end{split}
\end{equation}
\sphinxAtStartPar
By summing the Eqs. \eqref{equation:LBM-Study/Convective-Equation:cv-Kn1-B0}, \eqref{equation:LBM-Study/Convective-Equation:cv-Kn2-B0}, and \eqref{equation:LBM-Study/Convective-Equation:cv-Kn3-B0}, and also extending to higher orders of Knudsen, we have
\begin{equation*}
\begin{split}
\partial_{t}  u + \partial_{\alpha} B_{\alpha} + \underbrace{\partial_{\alpha} \left( m^{(1)}_{\alpha} + m^{(2)}_{\alpha} + ...\right) }_{\textrm{Numerical Errors}} =0,
\end{split}
\end{equation*}
\sphinxAtStartPar
where these numerical error can be related to the numerical error from the approximation of a derivative term by an expansion in Taylor series.

\end{sphinxuseclass}

\subsubsection{Chapmann\sphinxhyphen{}Enskog Analysis (Full)}
\label{\detokenize{LBM-Study/Convective-Equation:chapmann-enskog-analysis-full}}
\begin{sphinxuseclass}{toggle}
\sphinxAtStartPar
Applying the Chapmann\sphinxhyphen{}Enskog procedure to LB equation, we expand the term \(f_{i}\left(\boldsymbol{x}+\boldsymbol{e}_i \Delta t, t+\Delta t\right)\) in a Taylor series to recover the derivative form of the equation, i.e.,
\begin{equation}\label{equation:LBM-Study/Convective-Equation:EqExp-cv-NotEq}
\begin{split}
f_{i}\left(x_{\alpha} + c_{i,\alpha} \Delta t, t+\Delta t\right)- f_{i}(x_{\alpha}, t)=\displaystyle\sum_{j=1}^{\infty}\frac{\Delta t^{j}}{j!}D_{t}^{j}f_{i}= -\left( \frac{ f_{i} - f_{i}^{eq} }{ \tau } \right) .
\end{split}
\end{equation}
\sphinxAtStartPar
Rescaling the dimensionless form of the Eq. \eqref{equation:LBM-Study/Convective-Equation:EqExp-cv-NotEq} in terms of the Knudsen number (\(Kn\)), we have
\begin{equation*}
\begin{split}
\displaystyle \sum^{\infty}_{j=1}  \frac{Kn^{(j-1)}}{j!} D_{t}^{j} f_{i} = - \frac{1}{Kn}\left( \frac{ f_{i} - f_{eq,i} }{ \tau } \right) \quad \quad \rightarrow \quad \quad \displaystyle \sum^{\infty}_{j=1}  \frac{Kn^{(j)}}{j!} D_{t}^{j} f_{i} = -  \frac{ f_{i} - f_{i}^{eq} }{ \tau },
\end{split}
\end{equation*}
\sphinxAtStartPar
applying the asymptotic expansion in both the distribution function (\(f_{i}=f_{i}^{(0)}+Kn f_{i}^{(1)} + Kn^{2} f_{i}^{(2)}+\cdots\)) and time partial derivative (\(\partial_{t}=\partial_{t}^{(0)}+ Kn \partial_{t}^{(1)}+Kn^{2} \partial_{t}^{(2)}+\cdots\)), and separating the equation in orders up to the order \(Kn^{2}\):
\begin{equation}\label{equation:LBM-Study/Convective-Equation:Chap-Kn-cv-NotEq}
\begin{split}
\begin{array}{ll}
    (Kn^{(0)}):& f_{i}^{(0)}= f_{i}^{eq},\\
    (Kn^{(1)}):& \left( \partial_{t}^{(0)} + c_{i,\alpha}\partial_{\alpha} \right)f_{i}^{(0)}= - \displaystyle\frac{ f_{i}^{(1)} }{ \tau } , \\
    (Kn^{(2)}):& \partial_{t}^{(1)} f_{i}^{(0)} + \left( \partial_{t}^{(0)} + c_{i,\alpha}\partial_{\alpha} \right)f_{i}^{(1)} + \displaystyle\frac{\left( \partial_{t}^{(0)} + c_{i,\alpha}\partial_{\alpha} \right)^{2} f_{i}^{(0)}}{2}= - \displaystyle\frac{ f_{i}^{(2)} }{ \tau },\\
    \textrm{ou} & \\
    (Kn^{(2)}):& \partial_{t}^{(1)} f_{i}^{(0)} + \left( \partial_{t}^{(0)} + c_{i,\alpha}\partial_{\alpha} \right)\displaystyle\left(1 - \frac{1}{2\tau}\right) f_{i}^{(1)} =  - \displaystyle\frac{ f_{i}^{(2)} }{ \tau } \quad\quad \rightarrow \quad\quad  f_{i}^{(1)}\textrm{ Formulation} \\
    \textrm{ou} & \\
    (Kn^{(2)}):& \partial_{t}^{(1)} f_{i}^{(0)} + \left( \partial_{t}^{(0)} + c_{i,\alpha}\partial_{\alpha} \right)^{2} \displaystyle\left(\frac{1}{2} - \tau\right) f_{i}^{(0)} =  - \displaystyle\frac{ f_{i}^{(2)} }{ \tau } \quad\quad \rightarrow  \quad\quad f_{i}^{(0)}\textrm{ Formulation} \\
    (Kn^{(3)}):& \partial_{t}^{(2)} f_{i}^{(0)} +\partial_{t}^{(1)} f_{i}^{(1)} + \partial_{t}^{(1)}\left( \left( \partial_{t}^{(0)} + c_{i,\alpha}\partial_{\alpha} \right)f_{i}^{(0)} \right) + \left( \partial_{t}^{(0)} + c_{i,\alpha}\partial_{\alpha} \right)f_{i}^{(2)} + \displaystyle\frac{\left( \partial_{t}^{(0)} + c_{i,\alpha}\partial_{\alpha} \right)^{2} f_{i}^{(1)}}{2} \\
    & + \displaystyle\frac{\left( \partial_{t}^{(0)} + c_{i,\alpha}\partial_{\alpha} \right)^{3} f_{i}^{(0)}}{6} = - \displaystyle\frac{ f_{i}^{(3)} }{ \tau },  \\
    \textrm{ou} & \\
    (Kn^{(3)}):& \partial_{t}^{(2)} f_{i}^{(0)} + \left(1-\displaystyle\frac{1}{\tau}\right)\partial_{t}^{(1)} f_{i}^{(1)} + \displaystyle\left(\frac{1}{2} - \frac{1}{6\tau}\right) \left( \partial_{t}^{(0)} + c_{i,\alpha}\partial_{\alpha} \right)^{2} f_{i}^{(1)} + \left( \partial_{t}^{(0)} + c_{i,\alpha}\partial_{\alpha} \right)f_{i}^{(2)} =  - \displaystyle\frac{f_{i}^{(3)}}{\tau},
 \end{array}
\end{split}
\end{equation}
\sphinxAtStartPar
\sphinxstylestrong{Zero\sphinxhyphen{}Order Moment Balance}

\sphinxAtStartPar
To retrieve the balance equation, we sum the Eq. \eqref{equation:LBM-Study/Convective-Equation:Chap-Kn-cv-NotEq} for \(Kn^{(1)}\) and \(Kn^{(2)}\) over \(\sum_{i=0} \):
\begin{equation}\label{equation:LBM-Study/Convective-Equation:cv-Kn1-B0-all}
\begin{split}
\begin{array}{l}
(Kn^{(1)}): \displaystyle\sum_{i}\left( \partial_{t}^{(0)} + c_{i,\alpha}\partial_{\alpha} \right)f_{i}^{(0)}= - \displaystyle\sum_{i}\left(\displaystyle\frac{ f_{i}^{(1)} }{ \tau } \right),\\
(Kn^{(1)}): \partial_{t}^{(0)}  u + \partial_{\alpha} B_{\alpha} =0,
\end{array}
\end{split}
\end{equation}\begin{equation}\label{equation:LBM-Study/Convective-Equation:cv-Kn2-B0-all}
\begin{split}
\begin{array}{ll}
(Kn^{(2)}):& \displaystyle\sum_{i}\left[ \partial_{t}^{(1)} f_{i}^{(0)} + \displaystyle\left(1 - \frac{1}{2\tau}\right)\left( \partial_{t}^{(0)} + c_{i,\alpha}\partial_{\alpha} \right) f_{i}^{(1)}   \right]=  - \displaystyle\sum_{i}\left(\displaystyle\frac{ f_{i}^{(2)} }{ \tau } \right), \\
(Kn^{(2)}):& \partial_{t}^{(1)} u + \displaystyle\left(1 - \frac{1}{2\tau}\right)\partial_{\alpha} m^{(1)}_{\alpha} = 0,
\end{array}
\end{split}
\end{equation}
\sphinxAtStartPar
To determine the unknow term \(m^{(1)}_{\alpha}\), we sum the Eq. \eqref{equation:LBM-Study/Convective-Equation:Chap-Kn-cv-NotEq} for \(Kn^{(1)}\) over the first\sphinxhyphen{}order moment:
\begin{equation}\label{equation:LBM-Study/Convective-Equation:cv-Kn1-B1-all}
\begin{split}
\begin{array}{l}
(Kn^{(1)}): \displaystyle\sum_{i}\left( \partial_{t}^{(0)} + c_{i,\beta}\partial_{\beta} \right)f_{i}^{(0)}c_{i,\alpha}= - \displaystyle\sum_{i}\left(\displaystyle\frac{ f_{i}^{(1)} }{ \tau } c_{i,\alpha} \right),\\
(Kn^{(1)}): \partial_{t}^{(0)}  B_{\alpha} + \partial_{\beta} \left(uc_{s}^{2}\right)\delta_{\alpha\beta} = - \displaystyle\frac{m^{(1)}_{\alpha}}{\tau} \quad \quad \rightarrow \quad \quad   m^{(1)}_{\alpha} =-\tau\left( \partial_{t}^{(0)}  B_{\alpha} + \partial_{\beta} \left(uc_{s}^{2}\right)\delta_{\alpha\beta} \right).
\end{array}
\end{split}
\end{equation}
\sphinxAtStartPar
Extending the analysis up to \(Kn^{(3)}\)
\begin{equation}\label{equation:LBM-Study/Convective-Equation:cv-Kn3-B0-all}
\begin{split}
\begin{array}{ll}
(Kn^{(3)}):& \displaystyle\sum_{i}\Bigg[ \partial_{t}^{(2)} f_{i}^{(0)} + \left(1-\displaystyle\frac{1}{\tau}\right)\partial_{t}^{(1)} f_{i}^{(1)} + \left(\frac{1}{2} - \frac{1}{6\tau}\right) \left( \partial_{t}^{(0)} + c_{i,\alpha}\partial_{\alpha} \right)^{2} f_{i}^{(1)} \\
& + \left( \partial_{t}^{(0)} + c_{i,\alpha}\partial_{\alpha} \right)f_{i}^{(2)} \Bigg] =  - \displaystyle\sum_{i}\left[\displaystyle\frac{f_{i}^{(3)}}{\tau} \right], \\
(Kn^{(3)}):& \partial_{t}^{(2)} u + \partial_{t}^{(1)}\left(\partial_{\alpha} m^{(1)}_{\alpha} + \partial_{\beta} m^{(1)}_{\beta}\right) + \partial_{\alpha}\partial_{\beta} m^{(1)}_{\alpha\beta} + \partial_{\alpha} m^{(2)}_{\alpha} = 0.
\end{array}
\end{split}
\end{equation}
\sphinxAtStartPar
By summing the Eqs. \eqref{equation:LBM-Study/Convective-Equation:cv-Kn1-B0-all}, \eqref{equation:LBM-Study/Convective-Equation:cv-Kn2-B0-all}, and \eqref{equation:LBM-Study/Convective-Equation:cv-Kn3-B0-all}, and also extending to higher orders of Knudsen, we have
\begin{equation*}
\begin{split}
\partial_{t}  u + \partial_{\alpha} B_{\alpha} - \partial_{\alpha}\left(\nu \partial_{\alpha} u \right) + \underbrace{\partial_{\alpha} m^{(2)}_{\alpha} + \partial_{\alpha}\partial_{\beta} m^{(1)}_{\alpha\beta} + \partial_{t}^{(1)}\left(\partial_{\alpha} m^{(1)}_{\alpha} + \partial_{\beta} m^{(1)}_{\beta}\right) - \left(\tau - \frac{1}{2} \right) \partial_{t}^{(0)} \partial_{\alpha} B_{\alpha} + ...  }_{\text{Numerical Errors}} =0,
\end{split}
\end{equation*}
\sphinxAtStartPar
where \(\nu=c_{s}^{2}(\tau-1/2)\). The numerical errors can be corrected using extra moments in the equilibrium distribution function and some ones can neglected due to its relevance decay exponentially with the decrease of \(\Delta x\) and \(\Delta t\).

\end{sphinxuseclass}

\subsubsection{Boundary Conditions}
\label{\detokenize{LBM-Study/Convective-Equation:boundary-conditions}}
\sphinxAtStartPar
The boundary conditions for the lattices can be derived by solving a linear system of known moments.


\begin{savenotes}\sphinxattablestart
\sphinxthistablewithglobalstyle
\centering
\begin{tabulary}{\linewidth}[t]{TTTTTT}
\sphinxtoprule
\sphinxstyletheadfamily 
\sphinxAtStartPar
D1Q3:
&\sphinxstyletheadfamily 
\sphinxAtStartPar
Boundaries
&
\sphinxAtStartPar

&\sphinxstyletheadfamily 
\sphinxAtStartPar
Layers
&
\sphinxAtStartPar

&
\sphinxAtStartPar

\\
\sphinxmidrule
\sphinxtableatstartofbodyhook
\sphinxAtStartPar

&
\sphinxAtStartPar

&
\sphinxAtStartPar

&
\sphinxAtStartPar
West
&
\sphinxAtStartPar

&
\sphinxAtStartPar
East
\\
\sphinxhline
\sphinxAtStartPar

&
\sphinxAtStartPar
Unknown \(f_{i}\)
&
\sphinxAtStartPar

&
\sphinxAtStartPar
\(f_2=u_{e} - f_1 - f_0\)
&
\sphinxAtStartPar

&
\sphinxAtStartPar
\(f_1=u_{w} - f_2 - f_0\)
\\
\sphinxhline
\sphinxAtStartPar

&
\sphinxAtStartPar

&
\sphinxAtStartPar

&
\sphinxAtStartPar

&
\sphinxAtStartPar

&
\sphinxAtStartPar

\\
\sphinxbottomrule
\end{tabulary}
\sphinxtableafterendhook\par
\sphinxattableend\end{savenotes}


\section{Benchmark \sphinxhyphen{} 1D Non\sphinxhyphen{}Linear Convective formulation of the Buckley\sphinxhyphen{}Leverett Equation}
\label{\detokenize{LBM-Study/Convective-Equation:benchmark-1d-non-linear-convective-formulation-of-the-buckley-leverett-equation}}
\sphinxAtStartPar
Considering a pure convective formulation of the Buckley\sphinxhyphen{}Leverett equation, where gravitational effects and capillary pressure are neglected, we have:
\begin{equation}\label{equation:LBM-Study/Convective-Equation:EqMassMTP-BL-Sim}
\begin{split}
\partial_{t} (s_{w})+ \partial_{\alpha}\left[\displaystyle\frac{\left<u\right>_{\alpha}}{\varphi} F(s_w)\right] = 0 ,
\end{split}
\end{equation}
\sphinxAtStartPar
with the water fraction flow given by,
\begin{equation}\label{equation:LBM-Study/Convective-Equation:EqFractionalFlow}
\begin{split}
F(s_w) = \frac{ s_{w}^{2} }{ s_{w}^{2} + M_{\mu}(1-s_{w})^{2}}.
\end{split}
\end{equation}
\sphinxAtStartPar
The simulation parameters are chosen so that the average velocity component \(\langle u \rangle_{\alpha} / \varphi\) equals \(1\,\mathrm{m/s}\) and the dynamic viscosity ratio \(M_{\mu}\) equals \(1\). We consider a one\sphinxhyphen{}dimensional porous medium of length \(L_x = 20\,\mathrm{m}\), initially saturated with oil (see \hyperref[\detokenize{LBM-Study/Convective-Equation:d-displace-bl}]{Fig.\@ \ref{\detokenize{LBM-Study/Convective-Equation:d-displace-bl}}}). Water is injected at the inlet \((x=0)\) under a Dirichlet boundary condition \(s_{w}(x=0)=1\). At the outlet \((x=L)\), a Neumann boundary condition specifying a zero gradient in water saturation is applied.

\begin{figure}[htbp]
\centering
\capstart

\noindent\sphinxincludegraphics[scale=0.6]{{1D-Displace-BL}.png}
\caption{Geometry of the non\sphinxhyphen{}linear convective problem: initial and boundary conditions.}\label{\detokenize{LBM-Study/Convective-Equation:d-displace-bl}}\end{figure}


\subsection{Transient Analytical Solution}
\label{\detokenize{LBM-Study/Convective-Equation:transient-analytical-solution}}
\sphinxAtStartPar
The non\sphinxhyphen{}linear convective problem given by the Eq. \eqref{equation:LBM-Study/Convective-Equation:EqMassMTP-BL-Sim} has an analytical solution dependent of the initial condition (Riemann problem). In the present case, the solution for the specific initial condition is given by
\begin{equation*}
\begin{split}
s_{w}(x,t)= \left\{  \begin{array}{ccc}
    t \displaystyle F' & \textrm{for} & x\leq \zeta t \\
    & & \\
    0 & \textrm{for} & x >  \zeta t \\ 
\end{array} \right.
\end{split}
\end{equation*}
\sphinxAtStartPar
where \(\zeta=F'(\tilde{s}_{w})=F(\tilde{s}_{w})/\tilde{s}_{w}\) is the shock wave speed, and \(\tilde{s}_{w}\) is the sharp front saturation determined analytically by
\begin{equation*}
\begin{split}
\tilde{s}_{w}=\displaystyle\frac{F(\tilde{s}_{w})}{F'(\tilde{s}_{w})}=\sqrt{\displaystyle\frac{M_{\mu}}{M_{\mu}+1}}, \quad \quad \textrm{where} \quad \quad F^{'}(s_{w})=\frac{dF}{ds_{w}}=\frac{2s_{w}M_{\mu}(1-s_{w})}{(s_{w}^{2}+M_{\mu}(1-s_{w})^{2})^{2}}.
\end{split}
\end{equation*}

\subsubsection{Forward Time and Forward Space \sphinxhyphen{} FTFS}
\label{\detokenize{LBM-Study/Convective-Equation:id1}}

\subsubsection{FDM Solution:}
\label{\detokenize{LBM-Study/Convective-Equation:fdm-solution}}
\begin{sphinxuseclass}{cell}
\begin{sphinxuseclass}{tag_hide-input}
\begin{sphinxuseclass}{tag_hide-output}
\end{sphinxuseclass}
\end{sphinxuseclass}
\end{sphinxuseclass}

\subsubsection{Lattice Boltzmann Framework of FTFS:}
\label{\detokenize{LBM-Study/Convective-Equation:lattice-boltzmann-framework-of-ftfs}}
\begin{sphinxuseclass}{cell}
\begin{sphinxuseclass}{tag_hide-input}
\begin{sphinxuseclass}{tag_hide-output}
\end{sphinxuseclass}
\end{sphinxuseclass}
\end{sphinxuseclass}

\subsubsection{Numerical and Analytical Comparisson}
\label{\detokenize{LBM-Study/Convective-Equation:numerical-and-analytical-comparisson}}
\begin{sphinxuseclass}{cell}
\begin{sphinxuseclass}{tag_hide-input}\begin{sphinxVerbatimOutput}

\begin{sphinxuseclass}{cell_output}
\noindent\sphinxincludegraphics{{2dceeb4bae91b609db7f314c69b1836356bdf744293821ec550a0aab1d9fa921}.png}

\end{sphinxuseclass}\end{sphinxVerbatimOutput}

\end{sphinxuseclass}
\end{sphinxuseclass}

\subsection{Forward\sphinxhyphen{}Time and Central\sphinxhyphen{}Space with Small Diffusivity (FTCS\sphinxhyphen{}SD)}
\label{\detokenize{LBM-Study/Convective-Equation:id2}}

\subsubsection{Lattice Boltzmann Method (D1Q3):}
\label{\detokenize{LBM-Study/Convective-Equation:lattice-boltzmann-method-d1q3}}
\begin{sphinxuseclass}{cell}
\begin{sphinxuseclass}{tag_hide-input}
\begin{sphinxuseclass}{tag_hide-output}
\end{sphinxuseclass}
\end{sphinxuseclass}
\end{sphinxuseclass}

\subsubsection{FDM approach of LBM Framework D1Q3:}
\label{\detokenize{LBM-Study/Convective-Equation:fdm-approach-of-lbm-framework-d1q3}}
\begin{sphinxuseclass}{cell}
\begin{sphinxuseclass}{tag_hide-input}
\begin{sphinxuseclass}{tag_hide-output}
\end{sphinxuseclass}
\end{sphinxuseclass}
\end{sphinxuseclass}

\subsubsection{Numerical and Analytical Comparisson}
\label{\detokenize{LBM-Study/Convective-Equation:id3}}
\begin{sphinxuseclass}{cell}
\begin{sphinxuseclass}{tag_hide-input}\begin{sphinxVerbatimOutput}

\begin{sphinxuseclass}{cell_output}
\noindent\sphinxincludegraphics{{0ecf32fa32dd9a1868357f9422b464befeec339ae0f3ff7364b0268ba3bddf66}.png}

\end{sphinxuseclass}\end{sphinxVerbatimOutput}

\end{sphinxuseclass}
\end{sphinxuseclass}
\sphinxstepscope


\chapter{Convective\sphinxhyphen{}Diffusive Equation}
\label{\detokenize{LBM-Study/Convective-Diffusive-Equation:convective-diffusive-equation}}\label{\detokenize{LBM-Study/Convective-Diffusive-Equation::doc}}
\sphinxAtStartPar
The convective\sphinxhyphen{}diffusive equation is a partial diferential equation (PDE) commonly written as
\begin{equation}\label{equation:LBM-Study/Convective-Diffusive-Equation:cvdif-eq-lb}
\begin{split}
\partial_t u + \partial_{\alpha}B_{\alpha} - \partial_{\alpha}\left(\nu \partial_{\alpha} C \right)  = 0
\end{split}
\end{equation}
\sphinxAtStartPar
where \(u(\alpha,t)\) is a space and time dependent scalar parameter, \(B_{\alpha}(\alpha,t)\) is a vector a \(C\) a scalar term, may be time and space dependent parameters and can also be also expressed as linear or non\sphinxhyphen{}linear functions of \(u(\alpha,t)\). The term \(\nu\) is the diffusion coefficient, a constituve term that controls the intensity of diffusive effect.


\section{Numerical Discretization}
\label{\detokenize{LBM-Study/Convective-Diffusive-Equation:numerical-discretization}}
\sphinxAtStartPar
Discretizing the above equation in regular grid of 1D domain, we have

\begin{figure}[htbp]
\centering
\capstart

\noindent\sphinxincludegraphics[scale=2.5]{{mesh-1-convective-diffusive}.png}
\caption{1D regular grid.}\label{\detokenize{LBM-Study/Convective-Diffusive-Equation:mesh-1-convective-diffusive}}\end{figure}


\subsection{FDM Discretization}
\label{\detokenize{LBM-Study/Convective-Diffusive-Equation:fdm-discretization}}
\sphinxAtStartPar
To explore different discretization approaches, we adopt forward\sphinxhyphen{}time, forward\sphinxhyphen{}space (commonly referred to as the upwind scheme), as well as central\sphinxhyphen{}difference schemes.


\subsubsection{Forward\sphinxhyphen{}Time and Forward\sphinxhyphen{}Space (FTFSCS)}
\label{\detokenize{LBM-Study/Convective-Diffusive-Equation:forward-time-and-forward-space-ftfscs}}
\sphinxAtStartPar
Discretizing the Eq. \eqref{equation:LBM-Study/Convective-Diffusive-Equation:cvdif-eq-lb} through finite diference method, we have for forward\sphinxhyphen{}time, forward\sphinxhyphen{}space for advection and central\sphinxhyphen{}space for diffusion (FTFSCS) discretizations:
\begin{equation*}
\begin{split}
\begin{array}{ll}
\partial_t u + \partial_{\alpha}B_{\alpha} - \partial_{\alpha}\left(\nu \partial_{\alpha} C \right) = & \displaystyle\frac{u(x,t + \Delta_{t}) - u(x,t) }{\Delta_{t}} + \frac{B_{\alpha}(x,t) - B_{\alpha}(x + \Delta_{x},t) }{\Delta_{x}} \\
& - \nu \displaystyle\frac{ C(x - \Delta_{x},t) -2C(x,t) + C(x + \Delta_{x},t) }{\Delta_{x}^{2}} =0,
\end{array}
\end{split}
\end{equation*}
\sphinxAtStartPar
isolating the term \(u(x,t + \Delta_{t})\), we can describe the \(u\) time evolution:
\begin{equation*}
\begin{split}
u(x,t + \Delta_{t})   = u(x,t) - \frac{\Delta_{t}}{\Delta_{x}} \left( B_{\alpha}(x,t) - B_{\alpha}(x + \Delta_{x},t) \right) + \frac{\nu\Delta_{t}}{\Delta_{x}^{2}} \left( C(x - \Delta_{x},t) -2C(x,t) + C(x + \Delta_{x},t) \right) 
\end{split}
\end{equation*}

\subsubsection{Forward\sphinxhyphen{}Time and Central\sphinxhyphen{}Space with Small Diffusivity (FTCS)}
\label{\detokenize{LBM-Study/Convective-Diffusive-Equation:forward-time-and-central-space-with-small-diffusivity-ftcs}}
\sphinxAtStartPar
Given the discretization of the convective\sphinxhyphen{}diffusive equation with pure forwar\sphinxhyphen{}time and central\sphinxhyphen{}space in advective and diffusive terms (FTCS), we have
\begin{equation*}
\begin{split}
\begin{array}{ll}
\partial_t u + \partial_{x}B_{\alpha} - \partial_{\alpha}\left( \nu \partial_{\alpha}u \right) =& \displaystyle\frac{u(x,t + \Delta_{t}) - u(x,t) }{\Delta_{t}} + \frac{B_{\alpha}(x+\Delta_{x},t) - B_{\alpha}(x - \Delta_{x},t) }{2\Delta_{x}} \\
& - \nu \displaystyle\frac{ C(x - \Delta_{x},t) -2C(x,t) + C(x + \Delta_{x},t) }{\Delta_{x}^{2}} =0,
\end{array}
\end{split}
\end{equation*}
\sphinxAtStartPar
isolating the term \(u(x,t + \Delta_{t})\), we can describe the \(u\) time evolution:
\begin{equation*}
\begin{split}
u(x,t + \Delta_{t})   = u(x,t) - \frac{\Delta_{t}}{2\Delta_{x}} \left( B_{\alpha}(x+\Delta_{x},t) - B_{\alpha}(x - \Delta_{x},t) \right) + \frac{\nu\Delta_{t}}{\Delta_{x}^{2}} \left( C(x - \Delta_{x},t) -2C(x,t) + C(x + \Delta_{x},t) \right) 
\end{split}
\end{equation*}

\subsubsection{Forward\sphinxhyphen{}Time and Central\sphinxhyphen{}Space with Small Diffusivity (FTCS\sphinxhyphen{}SD)}
\label{\detokenize{LBM-Study/Convective-Diffusive-Equation:forward-time-and-central-space-with-small-diffusivity-ftcs-sd}}
\sphinxAtStartPar
Give that discretize convective equation with pure forwar\sphinxhyphen{}time and central\sphinxhyphen{}space (FTCS), employ a small diffusivity (FTCS\sphinxhyphen{}SD) is a alternative to improve stability and ensure a propertie conservation.
\begin{equation*}
\begin{split}
\begin{array}{c}
\partial_t u + \partial_{x}B_{\alpha} - \partial_{\alpha}\left( \nu_{u} \partial_{\alpha}u \right) - \partial_{\alpha}\left( \nu \partial_{\alpha}C \right) = \displaystyle\frac{u(x,t + \Delta_{t}) - u(x,t) }{\Delta_{t}} + \frac{B_{\alpha}(x+\Delta_{x},t) - B_{\alpha}(x - \Delta_{x},t) }{\Delta_{x}} -\\
\nu_{u} \displaystyle\frac{ u(x - \Delta_{x},t) -2u(x,t) + u(x + \Delta_{x},t) }{\Delta_{x}^{2}} - \nu\frac{ C(x - \Delta_{x},t) -2C(x,t) + C(x + \Delta_{x},t) }{\Delta_{x}^{2}} =0,
\end{array}
\end{split}
\end{equation*}
\sphinxAtStartPar
isolating the term \(u(x,t + \Delta_{t})\), we can describe the \(u\) time evolution:
\begin{equation*}
\begin{split}
\begin{array}{ll}
u(x,t + \Delta_{t})   =& u(x,t) - \displaystyle\frac{\Delta_{t}}{\Delta_{x}} \left( B_{\alpha}(x+\Delta_{x},t) - B_{\alpha}(x - \Delta_{x},t) \right) + \frac{\nu_{u}\Delta_{t}}{\Delta_{x}^{2}} \left( u(x - \Delta_{x},t) -2u(x,t) + u(x + \Delta_{x},t) \right) \\ 
& + \displaystyle\frac{\nu\Delta_{t}}{\Delta_{x}^{2}} \left( C(x - \Delta_{x},t) -2C(x,t) + C(x + \Delta_{x},t) \right) .
\end{array}
\end{split}
\end{equation*}

\subsection{Lattice Boltzmann Equation}
\label{\detokenize{LBM-Study/Convective-Diffusive-Equation:lattice-boltzmann-equation}}
\sphinxAtStartPar
Describing the problem through the BGK lattice Boltzmann equation:
\begin{equation}\label{equation:LBM-Study/Convective-Diffusive-Equation:LB-cvdif-Eq}
\begin{split}
f_{i}\left(x_{\alpha}+ c_{i,\alpha} \Delta t, t+\Delta t\right)- f_{i}(x_{\alpha}, t)= -\left( \frac{ f_{i} - f_{i}^{eq} }{ \tau } \right).
\end{split}
\end{equation}
\sphinxAtStartPar
The equilibrium distribution function is defined by
\begin{equation*}
\begin{split}
f^{eq}_{i} = w_{i}\left( u + \frac{c_{i,\alpha} B_{\alpha}}{c_{s}^{2}} + \frac{ \left( c_{s}^{2}C\delta_{\alpha\beta} + \Lambda_{\alpha\beta} - c_{s}^{2}u\delta_{\alpha\beta} \right) \left( c_{i,\alpha}c_{i,\beta} - c_{s}^{2}\delta_{\alpha\beta} \right)}{2c_{s}^{4}}  \right),
\end{split}
\end{equation*}
\sphinxAtStartPar
where \(\Lambda_{\alpha\beta}\) is a correction term defined by
\begin{equation*}
\begin{split}
\Lambda_{\alpha\beta} = \displaystyle\int  B_{\alpha}^{'}B_{\beta}^{'}  du = \displaystyle\int \left(\displaystyle\frac{\partial B_{\alpha}}{\partial u} \displaystyle\frac{\partial B_{\beta}}{\partial u} \right) du.
\end{split}
\end{equation*}

\subsubsection{Lattice Direction Moments}
\label{\detokenize{LBM-Study/Convective-Diffusive-Equation:lattice-direction-moments}}
\begin{sphinxuseclass}{cell}
\begin{sphinxuseclass}{tag_hide-input}
\end{sphinxuseclass}
\end{sphinxuseclass}

\paragraph{Raw Moments}
\label{\detokenize{LBM-Study/Convective-Diffusive-Equation:raw-moments}}
\begin{sphinxuseclass}{cell}
\begin{sphinxuseclass}{tag_hide-input}\begin{sphinxVerbatimOutput}

\begin{sphinxuseclass}{cell_output}\begin{equation*}
\begin{split}\displaystyle \begin{gathered}\underbrace{\sum_{i=0} f_{i}^{eq} =\sum_{i=0} f_{i} }_{\textrm{Zero-Order Moment}} =u,\quad \underbrace{\sum_{i=0} f_{i}^{eq}c_{i,x} }_{\textrm{x-First-Order Moment}} =B_{\alpha},\quad \underbrace{\sum_{i=0} f_{i}^{eq} c_{i,x}c_{i,x} }_{\textrm{xx-Second-Order Moment}} =\frac{C}{3} + \Lambda_{\alpha\beta},\quad \underbrace{\sum_{i=0} f_{i}^{eq} c_{i,x}c_{i,x}c_{i,x} }_{\textrm{xxx-Third-Order Moment}} =B_{\alpha}.\end{gathered}\end{split}
\end{equation*}
\end{sphinxuseclass}\end{sphinxVerbatimOutput}

\end{sphinxuseclass}
\end{sphinxuseclass}

\paragraph{Hermite Moments}
\label{\detokenize{LBM-Study/Convective-Diffusive-Equation:hermite-moments}}
\begin{sphinxuseclass}{cell}
\begin{sphinxuseclass}{tag_hide-input}\begin{sphinxVerbatimOutput}

\begin{sphinxuseclass}{cell_output}\begin{equation*}
\begin{split}\displaystyle \begin{gathered}\underbrace{\sum_{i=0} f_{i}^{eq} =\sum_{i=0} f_{i} }_{\textrm{Zero-Order Hermite Moment}} =u,\quad \quad \underbrace{\sum_{i=0} f_{i}^{eq}c_{i,x} }_{\textrm{x-First-Order Hermite Moment}} =B_{\alpha},\end{gathered}\end{split}
\end{equation*}\begin{equation*}
\begin{split}\displaystyle \begin{gathered}\underbrace{\sum_{i=0} f_{i}^{eq} \left(c_{i,x}c_{i,x} - c_{s}^{2} \right) }_{\textrm{xx-Second-Order Hermite Moment}} =\frac{C}{3} + \Lambda_{\alpha\beta} - \frac{u}{3},\quad \quad \underbrace{\sum_{i=0} f_{i}^{eq} \left(c_{i,x}c_{i,x} - 3c_{s}^{2} \right) c_{i,x} }_{\textrm{xxx-Third-Order Hermite Moment}} =0.\end{gathered}\end{split}
\end{equation*}
\end{sphinxuseclass}\end{sphinxVerbatimOutput}

\end{sphinxuseclass}
\end{sphinxuseclass}

\subsubsection{Chapmann\sphinxhyphen{}Enskog Analysis}
\label{\detokenize{LBM-Study/Convective-Diffusive-Equation:chapmann-enskog-analysis}}


\sphinxAtStartPar
Applying the Chapmann\sphinxhyphen{}Enskog procedure to LB equation, we expand the term \(f_{i}\left(\boldsymbol{x}+\boldsymbol{e}_i \Delta t, t+\Delta t\right)\) in a Taylor series to recover the derivative form of the equation, i.e.,
\begin{equation}\label{equation:LBM-Study/Convective-Diffusive-Equation:EqExp-cvdif-NotEq}
\begin{split}
f_{i}\left( x_{\alpha} + c_{i,\alpha} \Delta t, t+\Delta t\right)- f_{i}(x_{\alpha}, t)=\displaystyle\sum_{j=1}^{\infty}\frac{\Delta t^{j}}{j!}D_{t}^{j}f_{i}= -\left( \frac{ f_{i} - f_{i}^{eq} }{ \tau } \right) .
\end{split}
\end{equation}
\sphinxAtStartPar
Rescaling the dimensionless form of the Eq. \eqref{equation:LBM-Study/Convective-Diffusive-Equation:EqExp-cvdif-NotEq} in terms of the Knudsen number (\(Kn\)), we have
\begin{equation*}
\begin{split}
\displaystyle \sum^{\infty}_{j=1}  \frac{Kn^{(j-1)}}{j!} D_{t}^{j} f_{i} = - \frac{1}{Kn}\left( \frac{ f_{i} - f_{eq,i} }{ \tau } \right) \quad \quad \rightarrow \quad \quad \displaystyle \sum^{\infty}_{j=1}  \frac{Kn^{(j)}}{j!} D_{t}^{j} f_{i} = -  \frac{ f_{i} - f_{i}^{eq} }{ \tau },
\end{split}
\end{equation*}
\sphinxAtStartPar
applying the asymptotic expansion in both the distribution function (\(f_{i}=f_{i}^{(0)}+Kn f_{i}^{(1)} + Kn^{2} f_{i}^{(2)}+\cdots\)) and time partial derivative (\(\partial_{t}=\partial_{t}^{(0)}+ Kn \partial_{t}^{(1)}+Kn^{2} \partial_{t}^{(2)}+\cdots\)), and separating the equation in orders up to the order \(Kn^{2}\):
\begin{equation}\label{equation:LBM-Study/Convective-Diffusive-Equation:Chap-Kn-cvdif-NotEq}
\begin{split}
\begin{array}{ll}
    (Kn^{(0)}):& f_{i}^{(0)}= f_{i}^{eq},\\
    (Kn^{(1)}):& \left( \partial_{t}^{(0)} + c_{i,\alpha}\partial_{\alpha} \right)f_{i}^{(0)}= - \displaystyle\frac{ f_{i}^{(1)} }{ \tau } , \\
    (Kn^{(2)}):& \partial_{t}^{(1)} f_{i}^{(0)} + \left( \partial_{t}^{(0)} + c_{i,\alpha}\partial_{\alpha} \right)f_{i}^{(1)} + \displaystyle\frac{\left( \partial_{t}^{(0)} + c_{i,\alpha}\partial_{\alpha} \right)^{2} f_{i}^{(0)}}{2}= - \displaystyle\frac{ f_{i}^{(2)} }{ \tau },\\
    \textrm{ou}& \\
    (Kn^{(2)}):& \partial_{t}^{(1)} f_{i}^{(0)} + \left( \partial_{t}^{(0)} + c_{i,\alpha}\partial_{\alpha} \right)\displaystyle\left(1 - \frac{1}{2\tau}\right) f_{i}^{(1)} =  - \displaystyle\frac{ f_{i}^{(2)} }{ \tau } \quad\quad \rightarrow \quad\quad  f_{i}^{(1)}\textrm{ Formulation} \\
    \textrm{ou}& \\
    (Kn^{(2)}):& \partial_{t}^{(1)} f_{i}^{(0)} + \left( \partial_{t}^{(0)} + c_{i,\alpha}\partial_{\alpha} \right)^{2} \displaystyle\left(\frac{1}{2} - \tau\right) f_{i}^{(0)} =  - \displaystyle\frac{ f_{i}^{(2)} }{ \tau } \quad\quad \rightarrow  \quad\quad f_{i}^{(0)}\textrm{ Formulation} \\
    (Kn^{(3)}):& \partial_{t}^{(2)} f_{i}^{(0)} +\partial_{t}^{(1)} f_{i}^{(1)} + \partial_{t}^{(1)}\left( \left( \partial_{t}^{(0)} + c_{i,\alpha}\partial_{\alpha} \right)f_{i}^{(0)} \right) + \left( \partial_{t}^{(0)} + c_{i,\alpha}\partial_{\alpha} \right)f_{i}^{(2)} + \displaystyle\frac{\left( \partial_{t}^{(0)} + c_{i,\alpha}\partial_{\alpha} \right)^{2} f_{i}^{(1)}}{2} \\
    & +  \displaystyle\frac{\left( \partial_{t}^{(0)} + c_{i,\alpha}\partial_{\alpha} \right)^{3} f_{i}^{(0)}}{6} = - \displaystyle\frac{ f_{i}^{(3)} }{ \tau },  \\
    \textrm{ou}& \\
    (Kn^{(3)}):& \partial_{t}^{(2)} f_{i}^{(0)} + \left(1-\displaystyle\frac{1}{\tau}\right)\partial_{t}^{(1)} f_{i}^{(1)} + \displaystyle\left(\frac{1}{2} - \frac{1}{6\tau}\right) \left( \partial_{t}^{(0)} + c_{i,\alpha}\partial_{\alpha} \right)^{2} f_{i}^{(1)} + \left( \partial_{t}^{(0)} + c_{i,\alpha}\partial_{\alpha} \right)f_{i}^{(2)} =  - \displaystyle\frac{f_{i}^{(3)}}{\tau},
 \end{array}
\end{split}
\end{equation}
\sphinxAtStartPar
\sphinxstylestrong{Zero\sphinxhyphen{}Order Moment Balance}

\sphinxAtStartPar
To retrieve the balance equation, we sum the Eq. \eqref{equation:LBM-Study/Convective-Diffusive-Equation:Chap-Kn-cvdif-NotEq} for \(Kn^{(1)}\) and \(Kn^{(2)}\) over \(\sum_{i=0} \):
\begin{equation}\label{equation:LBM-Study/Convective-Diffusive-Equation:cvdif-Kn1-B0-all}
\begin{split}
\begin{array}{l}
(Kn^{(1)}): \displaystyle\sum_{i}\left( \partial_{t}^{(0)} + c_{i,\alpha}\partial_{\alpha} \right)f_{i}^{(0)}= - \displaystyle\sum_{i}\left(\displaystyle\frac{ f_{i}^{(1)} }{ \tau } \right),\\
(Kn^{(1)}): \partial_{t}^{(0)}  u + \partial_{\alpha} B_{\alpha} =0,
\end{array}
\end{split}
\end{equation}\begin{equation}\label{equation:LBM-Study/Convective-Diffusive-Equation:cvdif-Kn2-B0-all}
\begin{split}
\begin{array}{ll}
(Kn^{(2)}):& \displaystyle\sum_{i}\left[ \partial_{t}^{(1)} f_{i}^{(0)} + \displaystyle\left(1 - \frac{1}{2\tau}\right)\left( \partial_{t}^{(0)} + c_{i,\alpha}\partial_{\alpha} \right) f_{i}^{(1)}   \right]=  - \displaystyle\sum_{i}\left(\displaystyle\frac{ f_{i}^{(2)} }{ \tau } \right), \nonumber \\
(Kn^{(2)}):& \partial_{t}^{(1)} u + \displaystyle\left(1 - \frac{1}{2\tau}\right)\partial_{\alpha} m^{(1)}_{\alpha} = 0,
\end{array}
\end{split}
\end{equation}
\sphinxAtStartPar
To determine the unknow term \(m^{(1)}_{\alpha}\), we sum the Eq. \eqref{equation:LBM-Study/Convective-Diffusive-Equation:Chap-Kn-cvdif-NotEq} for \(Kn^{(1)}\) over the first\sphinxhyphen{}order moment:
\begin{equation}\label{equation:LBM-Study/Convective-Diffusive-Equation:cvdif-Kn1-B1-all}
\begin{split}
\begin{array}{l}
(Kn^{(1)}): \displaystyle\sum_{i}\left( \partial_{t}^{(0)} + c_{i,\beta}\partial_{\beta} \right)f_{i}^{(0)}c_{i,\alpha}= - \displaystyle\sum_{i}\left(\displaystyle\frac{ f_{i}^{(1)} }{ \tau } c_{i,\alpha} \right),\\
(Kn^{(1)}): \partial_{t}^{(0)}  B_{\alpha} + \partial_{\beta} \left(c_{s}^{2} C \delta_{\alpha\beta} + \Lambda_{\alpha\beta}\right) = - \displaystyle\frac{m^{(1)}_{\alpha}}{\tau} \quad \quad \rightarrow \quad \quad   m^{(1)}_{\alpha} =-\tau \partial_{\beta} \left(c_{s}^{2} C \delta_{\alpha\beta} \right).
\end{array}
\end{split}
\end{equation}
\sphinxAtStartPar
Extending the analysis up to \(Kn^{(3)}\)
\begin{equation}\label{equation:LBM-Study/Convective-Diffusive-Equation:cvdif-Kn3-B0-all}
\begin{split}
\begin{array}{ll}
(Kn^{(3)}):& \displaystyle\sum_{i}\Bigg[ \partial_{t}^{(2)} f_{i}^{(0)} + \left(1-\displaystyle\frac{1}{\tau}\right)\partial_{t}^{(1)} f_{i}^{(1)} + \displaystyle\left(\frac{1}{2} - \frac{1}{6\tau}\right) \left( \partial_{t}^{(0)} + c_{i,\alpha}\partial_{\alpha} \right)^{2} f_{i}^{(1)} \\
& + \left( \partial_{t}^{(0)} + c_{i,\alpha}\partial_{\alpha} \right)f_{i}^{(2)} \Bigg] =  - \displaystyle\sum_{i}\left[\displaystyle\frac{f_{i}^{(3)}}{\tau} \right], \\
(Kn^{(3)}):& \partial_{t}^{(2)} u + \partial_{t}^{(1)}\left(\partial_{\alpha} m^{(1)}_{\alpha} + \partial_{\beta} m^{(1)}_{\beta}\right) + \partial_{\alpha}\partial_{\beta} m^{(1)}_{\alpha\beta} + \partial_{\alpha} m^{(2)}_{\alpha} = 0.
\end{array}
\end{split}
\end{equation}
\sphinxAtStartPar
By summing the Eqs. \eqref{equation:LBM-Study/Convective-Diffusive-Equation:cvdif-Kn1-B0-all}, \eqref{equation:LBM-Study/Convective-Diffusive-Equation:cvdif-Kn2-B0-all}, and \eqref{equation:LBM-Study/Convective-Diffusive-Equation:cvdif-Kn3-B0-all}, and also extending to higher orders of Knudsen, we have
\begin{equation*}
\begin{split}
\partial_{t}  u + \partial_{\alpha} B_{\alpha} - \partial_{\alpha}\left(\nu \partial_{\alpha} C \right) + \underbrace{\partial_{\alpha} m^{(2)}_{\alpha} + \partial_{\alpha}\partial_{\beta} m^{(1)}_{\alpha\beta} + \partial_{t}^{(1)}\left(\partial_{\alpha} m^{(1)}_{\alpha} + \partial_{\beta} m^{(1)}_{\beta}\right) + ...  }_{\textrm{Numerical Errors}} =0,
\end{split}
\end{equation*}
\sphinxAtStartPar
where \(\nu=c_{s}^{2}(\tau-1/2)\). The numerical errors can be corrected using extra moments in the equilibrium distribution function and some ones can neglected due to its relevance decay exponentially with the decrease of \(\Delta x\) and \(\Delta t\).




\subsubsection{Boundary Conditions}
\label{\detokenize{LBM-Study/Convective-Diffusive-Equation:boundary-conditions}}
\sphinxAtStartPar
The boundary conditions for the lattices can be derived by solving a linear system of known moments.


\begin{savenotes}\sphinxattablestart
\sphinxthistablewithglobalstyle
\centering
\begin{tabulary}{\linewidth}[t]{TTTTTT}
\sphinxtoprule
\sphinxstyletheadfamily 
\sphinxAtStartPar
D1Q3:
&\sphinxstyletheadfamily 
\sphinxAtStartPar
Boundaries
&
\sphinxAtStartPar

&\sphinxstyletheadfamily 
\sphinxAtStartPar
Layers
&
\sphinxAtStartPar

&
\sphinxAtStartPar

\\
\sphinxmidrule
\sphinxtableatstartofbodyhook
\sphinxAtStartPar

&
\sphinxAtStartPar

&
\sphinxAtStartPar

&
\sphinxAtStartPar
West
&
\sphinxAtStartPar

&
\sphinxAtStartPar
East
\\
\sphinxhline
\sphinxAtStartPar

&
\sphinxAtStartPar
Unknown \(f_{i}\)
&
\sphinxAtStartPar

&
\sphinxAtStartPar
\(f_2=u_{e} - f_1 - f_0\)
&
\sphinxAtStartPar

&
\sphinxAtStartPar
\(f_1=u_{w} - f_2 - f_0\)
\\
\sphinxhline
\sphinxAtStartPar

&
\sphinxAtStartPar

&
\sphinxAtStartPar

&
\sphinxAtStartPar

&
\sphinxAtStartPar

&
\sphinxAtStartPar

\\
\sphinxbottomrule
\end{tabulary}
\sphinxtableafterendhook\par
\sphinxattableend\end{savenotes}


\section{1st Benchmark: Decay of Sinoidal Wave}
\label{\detokenize{LBM-Study/Convective-Diffusive-Equation:st-benchmark-decay-of-sinoidal-wave}}
\sphinxAtStartPar
Considering a 1D \sphinxstylestrong{convection–diffusion equation} with constant coefficients
\$\(
\frac{\partial u}{\partial t}+c\frac{\partial u}{\partial x} = \nu\frac{\partial^2 u}{\partial x^2}, \qquad x\in\mathbb{R},\ t\ge 0,
\)\$

\sphinxAtStartPar
with \sphinxstylestrong{sinusoidal initial condition}
\begin{equation*}
\begin{split}
u(x,0)=A\sin(kx+\phi),
\end{split}
\end{equation*}
\sphinxAtStartPar
where \(A\) is the initial amplitude, \(k\) the wavenumber, \(\phi\) a phase shift.

\sphinxAtStartPar
Then the analytical solution is
\begin{equation*}
\begin{split}
\boxed{u(x,t)=Ae^{-\nu k^{2}t}\sin\big(k(x-ct)+\phi\big) }
\end{split}
\end{equation*}
\sphinxAtStartPar
So:
\begin{itemize}
\item {} 
\sphinxAtStartPar
\sphinxstylestrong{Convection} shifts the wave: \(x \mapsto x-ct\) (wave travels at speed \(c\)).

\item {} 
\sphinxAtStartPar
\sphinxstylestrong{Diffusion} damps the amplitude exponentially: \(A(t)=Ae^{-\nu k^{2}t}\).

\end{itemize}

\sphinxAtStartPar
If you use a cosine initial condition \(u(x,0)=A\cos(kx+\phi)\), the solution is the same form:
\$\(
u(x,t)=Ae^{-\nu k^{2}t}\cos\big(k(x-ct)+\phi\big).
\)\$


\subsection{LBM D1Q3 Solution}
\label{\detokenize{LBM-Study/Convective-Diffusive-Equation:lbm-d1q3-solution}}
\begin{sphinxuseclass}{cell}
\begin{sphinxuseclass}{tag_hide-input}
\begin{sphinxuseclass}{tag_hide-output}
\end{sphinxuseclass}
\end{sphinxuseclass}
\end{sphinxuseclass}

\subsection{FDM Solution (FTCS\sphinxhyphen{}SD) (LBM Framework \protect\(\tau=1\protect\))}
\label{\detokenize{LBM-Study/Convective-Diffusive-Equation:fdm-solution-ftcs-sd-lbm-framework-tau-1}}
\begin{sphinxuseclass}{cell}
\begin{sphinxuseclass}{tag_hide-input}
\begin{sphinxuseclass}{tag_hide-output}
\end{sphinxuseclass}
\end{sphinxuseclass}
\end{sphinxuseclass}
\begin{sphinxuseclass}{cell}
\begin{sphinxuseclass}{tag_hide-input}\begin{sphinxVerbatimOutput}

\begin{sphinxuseclass}{cell_output}
\noindent\sphinxincludegraphics{{35b9af4b1bebc079297c735561c98297fae66219a93a0670699bf142cac5fbb8}.png}

\end{sphinxuseclass}\end{sphinxVerbatimOutput}

\end{sphinxuseclass}
\end{sphinxuseclass}

\subsection{FDM Solution (FTCS)}
\label{\detokenize{LBM-Study/Convective-Diffusive-Equation:fdm-solution-ftcs}}
\begin{sphinxuseclass}{cell}
\begin{sphinxuseclass}{tag_hide-input}\begin{sphinxVerbatimOutput}

\begin{sphinxuseclass}{cell_output}
\noindent\sphinxincludegraphics{{6c384a64662e5d31a2d1873b8040f1881b0c12b9fb25e5181597af986f84f0f9}.png}

\end{sphinxuseclass}\end{sphinxVerbatimOutput}

\end{sphinxuseclass}
\end{sphinxuseclass}

\subsection{FDM Solution (FTFSCS)}
\label{\detokenize{LBM-Study/Convective-Diffusive-Equation:fdm-solution-ftfscs}}
\begin{sphinxuseclass}{cell}
\begin{sphinxuseclass}{tag_hide-input}\begin{sphinxVerbatimOutput}

\begin{sphinxuseclass}{cell_output}
\noindent\sphinxincludegraphics{{157f9ae315ac4c268db31b1d29b58efe387501d8e58ca0e9ee0a7bf411dd8b88}.png}

\end{sphinxuseclass}\end{sphinxVerbatimOutput}

\end{sphinxuseclass}
\end{sphinxuseclass}
\begin{sphinxuseclass}{cell}
\begin{sphinxuseclass}{tag_hide-input}\begin{sphinxVerbatimOutput}

\begin{sphinxuseclass}{cell_output}
\noindent\sphinxincludegraphics{{e761eb9d778cc14706a1be27827e9a60f7c496c3c2961c0f2fb7c4c7f6b12497}.png}

\end{sphinxuseclass}\end{sphinxVerbatimOutput}

\end{sphinxuseclass}
\end{sphinxuseclass}

\section{2nd Benchmark: 1D Perturbation wave\sphinxhyphen{}decay}
\label{\detokenize{LBM-Study/Convective-Diffusive-Equation:nd-benchmark-1d-perturbation-wave-decay}}
\begin{sphinxuseclass}{cell}
\begin{sphinxuseclass}{tag_hide-input}\begin{sphinxVerbatimOutput}

\begin{sphinxuseclass}{cell_output}
\noindent\sphinxincludegraphics{{bd5de0ad3b16e4330e626a9a6d304a7c96f469d09f47e2616d2f6e3c44eabdb9}.png}

\end{sphinxuseclass}\end{sphinxVerbatimOutput}

\end{sphinxuseclass}
\end{sphinxuseclass}
\begin{sphinxuseclass}{cell}
\begin{sphinxuseclass}{tag_hide-input}\begin{sphinxVerbatimOutput}

\begin{sphinxuseclass}{cell_output}
\noindent\sphinxincludegraphics{{eef791156be27567701031b6548697fbf516acb96d41248c65e554192f329b12}.png}

\end{sphinxuseclass}\end{sphinxVerbatimOutput}

\end{sphinxuseclass}
\end{sphinxuseclass}
\begin{sphinxuseclass}{cell}
\begin{sphinxuseclass}{tag_hide-input}\begin{sphinxVerbatimOutput}

\begin{sphinxuseclass}{cell_output}
\noindent\sphinxincludegraphics{{58edb49af52dd291687fabed243ea341cc513836856569aeb086df52749e5c9e}.png}

\end{sphinxuseclass}\end{sphinxVerbatimOutput}

\end{sphinxuseclass}
\end{sphinxuseclass}
\sphinxstepscope


\chapter{Navier\sphinxhyphen{}Stokes Equation}
\label{\detokenize{Refs/LBM-Study/intro-nse-lbm/Navier-Stokes-Equation:navier-stokes-equation}}\label{\detokenize{Refs/LBM-Study/intro-nse-lbm/Navier-Stokes-Equation::doc}}
\sphinxAtStartPar
The Navier\sphinxhyphen{}Stokes equation is a non\sphinxhyphen{}linear PDE able to describe mass and momentum balance of fluid flow problems, commonly written as
\begin{equation*}
\begin{split}
\underbrace{\partial_t \rho + \partial_{\alpha}\rho u_{\alpha} = 0}_{\textrm{Mass Balance Equation}},
\end{split}
\end{equation*}\begin{equation}\label{equation:Refs/LBM-Study/intro-nse-lbm/Navier-Stokes-Equation:ns-eq-lb}
\begin{split}
\underbrace{\partial_t \rho u_{\alpha} + \partial_{\beta}\left(\rho u_{\alpha}\rho u_{\beta}\right) = \partial_{\beta}\left( -p \delta_{\alpha\beta} + \mu\left(\partial_{\alpha}u_{\beta}+\partial_{\beta}u_{\alpha}\right) \right) + \rho g_{\alpha}}_{\textrm{Momentum Balance Equation}} ,
\end{split}
\end{equation}
\sphinxAtStartPar
where \(\rho\) is the fluid density, \(u_{\alpha}\) is the \(\alpha\)\sphinxhyphen{}th component of the velocity vector field, \(p\) is the pressure field, \(\delta_{\alpha\beta}\) is the Kronecker delta, \(\mu\) is the dynamic viscosity (a constitutive parameter controlling the intensity of viscous diffusion), and \(g_{\alpha}\) is the \(\alpha\)\sphinxhyphen{}th component of the external body acceleration (e.g., gravity).


\section{Numerical Discretization: Lattice Boltzmann Equation}
\label{\detokenize{Refs/LBM-Study/intro-nse-lbm/Navier-Stokes-Equation:numerical-discretization-lattice-boltzmann-equation}}
\sphinxAtStartPar
Describing the problem through the BGK lattice Boltzmann equation:
\begin{equation}\label{equation:Refs/LBM-Study/intro-nse-lbm/Navier-Stokes-Equation:LB-ns-1st-Eq}
\begin{split}
\underbrace{f_{i}\left(x_{\alpha}+ c_{i,\alpha} \Delta t, t+\Delta t\right)- f_{i}(x_{\alpha}, t)= -\left( \frac{ f_{i} - f_{i}^{eq} }{ \tau} \right) }_{\textrm{1st order space-time discretization}}.
\end{split}
\end{equation}\begin{equation}\label{equation:Refs/LBM-Study/intro-nse-lbm/Navier-Stokes-Equation:LB-ns-2nd-Eq}
\begin{split}
\underbrace{\hat{f}_{i}\left(x_{\alpha}+ c_{i,\alpha} \Delta t, t+\Delta t\right)- \hat{f}_{i}(x_{\alpha}, t)= -\left( \frac{ \hat{f}_{i} - \hat{f}_{i}^{eq} }{ \tau } \right) + \left(1-\frac{1}{2\tau}F_{i} \right) }_{\textrm{2nd order space-time discretization}}.
\end{split}
\end{equation}
\sphinxAtStartPar
The equilibrium distribution function is defined by
\begin{equation*}
\begin{split}
f^{eq}_{i} = w_{i}\rho\left( 1 + \frac{c_{i,\alpha} u_{\alpha}}{c_{s}^{2}} + \frac{  u_{\alpha}u_{\beta}  \left( c_{i,\alpha}c_{i,\beta} - c_{s}^{2}\delta_{\alpha\beta} \right)}{2c_{s}^{4}}  \right).
\end{split}
\end{equation*}

\subsection{Lattice Direction Moments: D2Q9 Velocity\sphinxhyphen{}Space Discetization}
\label{\detokenize{Refs/LBM-Study/intro-nse-lbm/Navier-Stokes-Equation:lattice-direction-moments-d2q9-velocity-space-discetization}}
\begin{sphinxuseclass}{cell}
\begin{sphinxuseclass}{tag_hide-input}\begin{sphinxVerbatimOutput}

\begin{sphinxuseclass}{cell_output}
\noindent\sphinxincludegraphics{{5e80c4cd145c3ddfd698dc29a58326c9dd6de4352844f4fbfa38869da1c8afd6}.png}

\end{sphinxuseclass}\end{sphinxVerbatimOutput}

\end{sphinxuseclass}
\end{sphinxuseclass}

\subsubsection{Raw Moments}
\label{\detokenize{Refs/LBM-Study/intro-nse-lbm/Navier-Stokes-Equation:raw-moments}}
\begin{sphinxuseclass}{cell}
\begin{sphinxuseclass}{tag_hide-input}\begin{sphinxVerbatimOutput}

\begin{sphinxuseclass}{cell_output}\begin{equation*}
\begin{split}\displaystyle \begin{gathered}\underbrace{\sum_{i=0} f_{i}^{eq} =\sum_{i=0} f_{i} }_{\textrm{Zero-Order Moment}} =\rho,\quad \quad \underbrace{\sum_{i=0} f_{i}^{eq}c_{i,x} }_{\textrm{x-First-Order Moment}} =\rho u_{x},\quad \quad \underbrace{\sum_{i=0} f_{i}^{eq}c_{i,y} }_{\textrm{y-First-Order Moment}} =\rho u_{y}\\[10pt]\underbrace{\sum_{i=0} f_{i}^{eq} c_{i,x}c_{i,x} }_{\textrm{xx-Second-Order Moment}} =\rho \left(u_{x}^{2} + \frac{1}{3}\right),\quad \quad \underbrace{\sum_{i=0} f_{i}^{eq} c_{i,x}c_{i,y} }_{\textrm{xy-Second-Order Moment}} =\rho u_{x} u_{y},\quad \quad \underbrace{\sum_{i=0} f_{i}^{eq} c_{i,y}c_{i,y} }_{\textrm{yy-Second-Order Moment}} =\rho \left(u_{y}^{2} + \frac{1}{3}\right)\\[10pt]\underbrace{\sum_{i=0} f_{i}^{eq} c_{i,x}c_{i,x}c_{i,x} }_{\textrm{xxx-Third-Order Moment}} =\rho u_{x},\quad \quad \underbrace{\sum_{i=0} f_{i}^{eq} c_{i,x}c_{i,x}c_{i,y} }_{\textrm{xxy-Third-Order Moment}} =\frac{\rho u_{y}}{3},\quad \quad \underbrace{\sum_{i=0} f_{i}^{eq} c_{i,x}c_{i,y}c_{i,y} }_{\textrm{xyy-Third-Order Moment}} =\frac{\rho u_{x}}{3},\quad \quad \underbrace{\sum_{i=0} f_{i}^{eq} c_{i,y}c_{i,y}c_{i,y} }_{\textrm{yyy-Third-Order Moment}} =\rho u_{y}\end{gathered}\end{split}
\end{equation*}
\end{sphinxuseclass}\end{sphinxVerbatimOutput}

\end{sphinxuseclass}
\end{sphinxuseclass}

\subsubsection{Hermite Moments}
\label{\detokenize{Refs/LBM-Study/intro-nse-lbm/Navier-Stokes-Equation:hermite-moments}}
\begin{sphinxuseclass}{cell}
\begin{sphinxuseclass}{tag_hide-input}\begin{sphinxVerbatimOutput}

\begin{sphinxuseclass}{cell_output}\begin{equation*}
\begin{split}\displaystyle \begin{gathered}\underbrace{\sum_{i=0} f_{i}^{eq} =\sum_{i=0} f_{i} }_{\textrm{Zero-Order Hermite Moment}} =\rho,\quad \quad \underbrace{\sum_{i=0} f_{i}^{eq}c_{i,x} }_{\textrm{x-First-Order Hermite Moment}} =\rho u_{x},\quad \quad \underbrace{\sum_{i=0} f_{i}^{eq}c_{i,y} }_{\textrm{y-First-Order Hermite Moment}} =\rho u_{y}\\[10pt]\underbrace{\sum_{i=0} f_{i}^{eq} c_{i,x}c_{i,x} }_{\textrm{xx-Second-Order Moment}} =\rho u_{x}^{2},\quad \quad \underbrace{\sum_{i=0} f_{i}^{eq} c_{i,x}c_{i,y} }_{\textrm{xy-Second-Order Moment}} =\rho u_{x} u_{y},\quad \quad \underbrace{\sum_{i=0} f_{i}^{eq} c_{i,y}c_{i,y} }_{\textrm{yy-Second-Order Moment}} =\rho u_{y}^{2}\\[10pt]\underbrace{\sum_{i=0} f_{i}^{eq} c_{i,x}c_{i,x}c_{i,x} }_{\textrm{xxx-Third-Order Moment}} =0,\quad \quad \underbrace{\sum_{i=0} f_{i}^{eq} c_{i,x}c_{i,x}c_{i,y} }_{\textrm{xxy-Third-Order Moment}} =0,\quad \quad \underbrace{\sum_{i=0} f_{i}^{eq} c_{i,x}c_{i,y}c_{i,y} }_{\textrm{xyy-Third-Order Moment}} =0,\quad \quad \underbrace{\sum_{i=0} f_{i}^{eq} c_{i,y}c_{i,y}c_{i,y} }_{\textrm{yyy-Third-Order Moment}} =0\end{gathered}\end{split}
\end{equation*}
\end{sphinxuseclass}\end{sphinxVerbatimOutput}

\end{sphinxuseclass}
\end{sphinxuseclass}

\subsection{Chapmann\sphinxhyphen{}Enskog Analysis: 1st order Lattice Boltzmann Equation}
\label{\detokenize{Refs/LBM-Study/intro-nse-lbm/Navier-Stokes-Equation:chapmann-enskog-analysis-1st-order-lattice-boltzmann-equation}}
\begin{sphinxuseclass}{cell}
\begin{sphinxuseclass}{tag_hide-input}
\begin{sphinxuseclass}{tag_hide-output}
\end{sphinxuseclass}
\end{sphinxuseclass}
\end{sphinxuseclass}
\sphinxAtStartPar
Applying the Chapmann\sphinxhyphen{}Enskog procedure to LB equation, we expand the term \(f_{i}\left(\boldsymbol{x}+\boldsymbol{c}_{i,\alpha} \Delta t, t+\Delta t\right)\) in a Taylor series to recover the derivative form of the equation, i.e.,
\begin{equation}\label{equation:Refs/LBM-Study/intro-nse-lbm/Navier-Stokes-Equation:ns-EqExp-NotEq}
\begin{split}
f_{i}\left( x_{\alpha} + c_{i,\alpha} \Delta t, t+\Delta t\right)- f_{i}(x_{\alpha}, t)=\displaystyle\sum_{j=1}^{\infty}\frac{\Delta t^{j}}{j!}D_{t}^{j}f_{i}= -\left( \frac{ f_{i} - f_{i}^{eq} }{ \tau } \right),
\end{split}
\end{equation}
\sphinxAtStartPar
where \(D_{t}f_{i}=\partial_{t}f_{i} + c_{i,\alpha}\partial_{\alpha} f_{i}\).

\sphinxAtStartPar
Rescaling to the dimensionless form of the Eq. \eqref{equation:Refs/LBM-Study/intro-nse-lbm/Navier-Stokes-Equation:ns-EqExp-NotEq}, we rewrite it in terms of the Knudsen number (\(Kn\)) following the procedure presented in Kremer {[}\hyperlink{cite.Refs/Citations:id12}{Kre10}{]} Section 3.3.1:
\begin{equation*}
\begin{split}
\displaystyle \sum^{\infty}_{j=1}  \frac{Kn^{(j-1)}}{j!} D_{t}^{j} f_{i} = - \frac{1}{Kn}\left( \frac{ f_{i} - f_{eq,i} }{ \tau } \right) \quad \quad \rightarrow \quad \quad \displaystyle \sum^{\infty}_{j=1}  \frac{Kn^{(j)}}{j!} D_{t}^{j} f_{i} = -  \frac{ f_{i} - f_{i}^{eq} }{ \tau },
\end{split}
\end{equation*}
\sphinxAtStartPar
expanding the above equation up to fourth\sphinxhyphen{}order:

\begin{sphinxuseclass}{cell}
\begin{sphinxuseclass}{tag_hide-input}\begin{sphinxVerbatimOutput}

\begin{sphinxuseclass}{cell_output}\begin{equation*}
\begin{split}\displaystyle \begin{aligned}\frac{f_i^{eq} - f_{i}}{\tau} &= Kn^{(1)} \left(c_{i,\alpha} \partial_{\alpha}\,f_{i} + \partial_t\,f_{i}\right) \\[6pt] &\quad + Kn^{(2)} \left(\frac{c_{i,\alpha}^{2} \partial_{\alpha}^{2}\,f_{i}}{2} + c_{i,\alpha} \partial_{\alpha} \partial_t\,f_{i} + \frac{\partial_t^{2}\,f_{i}}{2}\right) \\[6pt] &\quad + Kn^{(3)} \left(\frac{c_{i,\alpha}^{3} \partial_{\alpha}^{3}\,f_{i}}{6} + \frac{c_{i,\alpha}^{2} \partial_{\alpha}^{2} \partial_t\,f_{i}}{2} + \frac{c_{i,\alpha} \partial_{\alpha} \partial_t^{2}\,f_{i}}{2} + \frac{\partial_t^{3}\,f_{i}}{6}\right) \\[6pt] &\quad + Kn^{(4)} \left(\frac{c_{i,\alpha}^{4} \partial_{\alpha}^{4}\,f_{i}}{24} + \frac{c_{i,\alpha}^{3} \partial_{\alpha}^{3} \partial_t\,f_{i}}{6} + \frac{c_{i,\alpha}^{2} \partial_{\alpha}^{2} \partial_t^{2}\,f_{i}}{4} + \frac{c_{i,\alpha} \partial_{\alpha} \partial_t^{3}\,f_{i}}{6} + \frac{\partial_t^{4}\,f_{i}}{24}\right)\end{aligned}\end{split}
\end{equation*}
\end{sphinxuseclass}\end{sphinxVerbatimOutput}

\end{sphinxuseclass}
\end{sphinxuseclass}
\sphinxAtStartPar
Applying the asymptotic expansion in both the distribution function (\(f_{i}=f_{i}^{(0)}+Kn f_{i}^{(1)} + Kn^{2} f_{i}^{(2)}+\cdots\)) and time partial derivative (\(\partial_{t}=\partial_{t}^{(0)}+ Kn \partial_{t}^{(1)}+Kn^{2} \partial_{t}^{(2)}+\cdots\)):

\begin{sphinxuseclass}{cell}
\begin{sphinxuseclass}{tag_hide-input}\begin{sphinxVerbatimOutput}

\begin{sphinxuseclass}{cell_output}\begin{equation*}
\begin{split}\displaystyle \small\begin{aligned}& - \frac{f_{i}^{(0)}}{\tau} + \frac{f_{i}^{eq}}{\tau} + - \frac{Kn f_{i}^{(1)}}{\tau} + - \frac{Kn^{2} f_{i}^{(2)}}{\tau} + - \frac{Kn^{3} f_{i}^{(3)}}{\tau} + - \frac{Kn^{4} f_{i}^{(4)}}{\tau} &= \end{aligned}\end{split}
\end{equation*}\begin{equation*}
\begin{split}\displaystyle \small\begin{aligned}& Kn \left(c_{i,\alpha} \partial_{\alpha}\,f_{i}^{(0)} + \partial_t^{(0)}\,f_{i}^{(0)}\right) \\[6pt] &\quad + Kn^{2} \left(\frac{c_{i,\alpha}^{2} \partial_{\alpha}^{2}\,f_{i}^{(0)}}{2} + c_{i,\alpha} \partial_{\alpha}\,f_{i}^{(1)} + c_{i,\alpha} \partial_{\alpha} \partial_t^{(0)}\,f_{i}^{(0)} + \frac{\left(\partial_t^{(0)}\right)^{2}\,f_{i}^{(0)}}{2} + \partial_t^{(1)}\,f_{i}^{(0)} + \partial_t^{(0)}\,f_{i}^{(1)}\right) \\[6pt] &\quad + Kn^{3} \left(\frac{c_{i,\alpha}^{3} \partial_{\alpha}^{3}\,f_{i}^{(0)}}{6} + \frac{c_{i,\alpha}^{2} \partial_{\alpha}^{2}\,f_{i}^{(1)}}{2} + \frac{c_{i,\alpha}^{2} \partial_{\alpha}^{2} \partial_t^{(0)}\,f_{i}^{(0)}}{2} + c_{i,\alpha} \partial_{\alpha}\,f_{i}^{(2)} + \frac{c_{i,\alpha} \partial_{\alpha} \left(\partial_t^{(0)}\right)^{2}\,f_{i}^{(0)}}{2} + c_{i,\alpha} \partial_{\alpha} \partial_t^{(1)}\,f_{i}^{(0)} + c_{i,\alpha} \partial_{\alpha} \partial_t^{(0)}\,f_{i}^{(1)} + \frac{\left(\partial_t^{(0)}\right)^{3}\,f_{i}^{(0)}}{6} + \partial_t^{(2)}\,f_{i}^{(0)} + \frac{\left(\partial_t^{(0)}\right)^{2}\,f_{i}^{(1)}}{2} + \partial_t^{(1)}\,f_{i}^{(1)} + \partial_t^{(0)}\,f_{i}^{(2)} + \partial_t^{(0)} \partial_t^{(1)}\,f_{i}^{(0)}\right) \\[6pt] &\quad + Kn^{4} \left(\frac{c_{i,\alpha}^{4} \partial_{\alpha}^{4}\,f_{i}^{(0)}}{24} + \frac{c_{i,\alpha}^{3} \partial_{\alpha}^{3}\,f_{i}^{(1)}}{6} + \frac{c_{i,\alpha}^{3} \partial_{\alpha}^{3} \partial_t^{(0)}\,f_{i}^{(0)}}{6} + \frac{c_{i,\alpha}^{2} \partial_{\alpha}^{2}\,f_{i}^{(2)}}{2} + \frac{c_{i,\alpha}^{2} \partial_{\alpha}^{2} \left(\partial_t^{(0)}\right)^{2}\,f_{i}^{(0)}}{4} + \frac{c_{i,\alpha}^{2} \partial_{\alpha}^{2} \partial_t^{(1)}\,f_{i}^{(0)}}{2} + \frac{c_{i,\alpha}^{2} \partial_{\alpha}^{2} \partial_t^{(0)}\,f_{i}^{(1)}}{2} + c_{i,\alpha} \partial_{\alpha}\,f_{i}^{(3)} + \frac{c_{i,\alpha} \partial_{\alpha} \left(\partial_t^{(0)}\right)^{3}\,f_{i}^{(0)}}{6} + c_{i,\alpha} \partial_{\alpha} \partial_t^{(2)}\,f_{i}^{(0)} + \frac{c_{i,\alpha} \partial_{\alpha} \left(\partial_t^{(0)}\right)^{2}\,f_{i}^{(1)}}{2} + c_{i,\alpha} \partial_{\alpha} \partial_t^{(1)}\,f_{i}^{(1)} + c_{i,\alpha} \partial_{\alpha} \partial_t^{(0)}\,f_{i}^{(2)} + c_{i,\alpha} \partial_{\alpha} \partial_t^{(0)} \partial_t^{(1)}\,f_{i}^{(0)} + \frac{\left(\partial_t^{(0)}\right)^{4}\,f_{i}^{(0)}}{24} + \frac{\left(\partial_t^{(1)}\right)^{2}\,f_{i}^{(0)}}{2} + \partial_t^{(3)}\,f_{i}^{(0)} + \frac{\left(\partial_t^{(0)}\right)^{3}\,f_{i}^{(1)}}{6} + \partial_t^{(2)}\,f_{i}^{(1)} + \frac{\left(\partial_t^{(0)}\right)^{2}\,f_{i}^{(2)}}{2} + \partial_t^{(1)}\,f_{i}^{(2)} + \partial_t^{(0)}\,f_{i}^{(3)} + \partial_t^{(0)} \partial_t^{(2)}\,f_{i}^{(0)} + \frac{\left(\partial_t^{(0)}\right)^{2} \partial_t^{(1)}\,f_{i}^{(0)}}{2} + \partial_t^{(0)} \partial_t^{(1)}\,f_{i}^{(1)}\right)\end{aligned}\end{split}
\end{equation*}
\end{sphinxuseclass}\end{sphinxVerbatimOutput}

\end{sphinxuseclass}
\end{sphinxuseclass}
\sphinxAtStartPar
Splitting the equation in orders up to the order \(Kn^{3}\) (References with High\sphinxhyphen{}Order expansion in Knudsen\sphinxhyphen{}order: Chai \sphinxstyleemphasis{et al.} {[}\hyperlink{cite.Refs/Citations:id13}{CHGS18}{]}):
\begin{equation}\label{equation:Refs/LBM-Study/intro-nse-lbm/Navier-Stokes-Equation:Chap-Kn-ns-e-NotEq}
\begin{split}
\scriptsize
\begin{array}{ll}
    (Kn^{(0)}):& f_{i}^{(0)}= f_{i}^{eq},\\
    (Kn^{(1)}):& \left( \partial_{t}^{(0)} + c_{i,\alpha}\partial_{\alpha} \right)f_{i}^{(0)}= - \displaystyle\frac{ f_{i}^{(1)} }{ \tau } , \\
    (Kn^{(2)}):& \partial_{t}^{(1)} f_{i}^{(0)} + \left( \partial_{t}^{(0)} + c_{i,\alpha}\partial_{\alpha} \right)f_{i}^{(1)} + \displaystyle\frac{\left( \partial_{t}^{(0)} + c_{i,\alpha}\partial_{\alpha} \right)^{2} f_{i}^{(0)}}{2}= - \displaystyle\frac{ f_{i}^{(2)} }{ \tau },\\
    \textrm{ou}\\
    (Kn^{(2)}):& \partial_{t}^{(1)} f_{i}^{(0)} + \left( \partial_{t}^{(0)} + c_{i,\alpha}\partial_{\alpha} \right)\displaystyle\left(1 - \frac{1}{2\tau}\right) f_{i}^{(1)} =  - \displaystyle\frac{ f_{i}^{(2)} }{ \tau } \quad\quad \rightarrow \quad\quad  f_{i}^{(1)}\textrm{ Formulation} \\
    \textrm{ou}\\
    (Kn^{(2)}):& \partial_{t}^{(1)} f_{i}^{(0)} + \left( \partial_{t}^{(0)} + c_{i,\alpha}\partial_{\alpha} \right)^{2} \displaystyle\left(\frac{1}{2} - \tau\right) f_{i}^{(0)} =  - \displaystyle\frac{ f_{i}^{(2)} }{ \tau } \quad\quad \rightarrow  \quad\quad f_{i}^{(0)}\textrm{ Formulation} \\
    (Kn^{(3)}):& \partial_{t}^{(2)} f_{i}^{(0)} +\partial_{t}^{(1)} f_{i}^{(1)} + \partial_{t}^{(1)}\left( \left( \partial_{t}^{(0)} + c_{i,\alpha}\partial_{\alpha} \right)f_{i}^{(0)} \right) + \left( \partial_{t}^{(0)} + c_{i,\alpha}\partial_{\alpha} \right)f_{i}^{(2)} + \displaystyle\frac{\left( \partial_{t}^{(0)} + c_{i,\alpha}\partial_{\alpha} \right)^{2} f_{i}^{(1)}}{2}  \\
    & +  \displaystyle\frac{\left( \partial_{t}^{(0)} + c_{i,\alpha}\partial_{\alpha} \right)^{3} f_{i}^{(0)}}{6} = - \displaystyle\frac{ f_{i}^{(3)} }{ \tau },  \\
    \textrm{ou}\\
    (Kn^{(3)}):& \partial_{t}^{(2)} f_{i}^{(0)} + \left(2-\displaystyle\frac{1}{\tau}\right)\partial_{t}^{(1)} f_{i}^{(1)} + \displaystyle\left(1 - \frac{1}{6\tau} - \tau\right) \left( \partial_{t}^{(0)} + c_{i,\alpha}\partial_{\alpha} \right)^{2} f_{i}^{(1)} =  - \displaystyle\frac{f_{i}^{(3)}}{\tau}, \quad\quad \rightarrow \quad\quad  f_{i}^{(1)}\textrm{ Formulation} \\
    \textrm{ou}\\
    (Kn^{(3)}):& \partial_{t}^{(2)} f_{i}^{(0)} + 2\left(\displaystyle\frac{1}{2}-\tau\right)\partial_{t}^{(1)}\left( \partial_{t}^{(0)} + c_{i,\alpha}\partial_{\alpha} \right)f_{i}^{(0)} + \displaystyle\left(\tau^{2} - \tau + \frac{1}{6}\right) \left( \partial_{t}^{(0)} + c_{i,\alpha}\partial_{\alpha} \right)^{3} f_{i}^{(0)} =  - \displaystyle\frac{f_{i}^{(3)}}{\tau}, \quad\quad \rightarrow \quad\quad  f_{i}^{(0)}\textrm{ Formulation} \\
 \end{array}
\end{split}
\end{equation}

\subsubsection{Zero\sphinxhyphen{}Order Moment Balance}
\label{\detokenize{Refs/LBM-Study/intro-nse-lbm/Navier-Stokes-Equation:zero-order-moment-balance}}
\sphinxAtStartPar
To retrieve the mass balance equation, we sum the Eq. \eqref{equation:Refs/LBM-Study/intro-nse-lbm/Navier-Stokes-Equation:Chap-Kn-ns-e-NotEq} for \(Kn^{(1)}\), \(Kn^{(2)}\), and \(Kn^{(3)}\) over \(\sum_{i=0} \):
\begin{equation}\label{equation:Refs/LBM-Study/intro-nse-lbm/Navier-Stokes-Equation:ns-e-Kn1-B0-all}
\begin{split}
\begin{array}{ll}
(Kn^{(1)}):& \displaystyle\sum_{i}\left( \partial_{t}^{(0)} + c_{i,\alpha}\partial_{\alpha} \right)f_{i}^{(0)}= - \underbrace{\displaystyle\sum_{i}\left(\displaystyle\frac{ f_{i}^{(1)} }{ \tau } \right)}_{=0},\\
(Kn^{(1)}):& \partial_{t}^{(0)}  \rho + \partial_{\alpha} \left( \rho u_{\alpha} \right) =0,
\end{array}
\end{split}
\end{equation}\begin{equation}\label{equation:Refs/LBM-Study/intro-nse-lbm/Navier-Stokes-Equation:ns-e-Kn2-B0-all}
\begin{split}
\begin{array}{ll}
(Kn^{(2)}):& \displaystyle\sum_{i}\left[ \partial_{t}^{(1)} f_{i}^{(0)} + \displaystyle\left(1 - \frac{1}{2\tau}\right)\left( \partial_{t}^{(0)} + c_{i,\alpha}\partial_{\alpha} \right) f_{i}^{(1)}   \right]=  - \displaystyle\sum_{i}\left(\displaystyle\frac{ f_{i}^{(2)} }{ \tau } \right), \nonumber \\
(Kn^{(2)}):& \partial_t^{(1)} \displaystyle\sum_i f_i^{(0)} + \underbrace{\left(1-\frac{1}{2\tau}\right) \partial_t^{(0)} \sum_i f_i^{(1)}}_{=0} + \underbrace{\left(1-\frac{1}{2\tau}\right) \partial_\alpha \sum_i c_{i,\alpha}f_i^{(1)}}_{=0} =\underbrace{-\frac{1}{\tau} \sum_i f_i^{(2)}}_{=0} \\
(Kn^{(2)}):& \partial_{t}^{(1)} \rho = 0,
\end{array}
\end{split}
\end{equation}
\sphinxAtStartPar
Extending the analysis up to \(Kn^{(3)}\)
\begin{equation}\label{equation:Refs/LBM-Study/intro-nse-lbm/Navier-Stokes-Equation:ns-e-Kn3-B0-all}
\begin{split}
\begin{array}{ll}
(Kn^{(3)}):& \displaystyle\sum_{i}\left[ \partial_{t}^{(2)} f_{i}^{(0)} + \left(2-\displaystyle\frac{1}{\tau}\right)\partial_{t}^{(1)} f_{i}^{(1)} + \displaystyle\left(1 - \frac{1}{6\tau} - \tau\right) \left( \partial_{t}^{(0)} + c_{i,\alpha}\partial_{\alpha} \right)^{2} f_{i}^{(1)} \right] =  - \displaystyle\sum_{i}\left[\displaystyle\frac{f_{i}^{(3)}}{\tau} \right], \\
(Kn^{(3)}):&  \partial_{t}^{(2)} \displaystyle\sum_{i}f_{i}^{(0)} + \underbrace{\left(2-\displaystyle\frac{1}{\tau}\right)\partial_{t}^{(1)} \sum_{i}f_{i}^{(1)} }_{=0} + \displaystyle\left(1 - \frac{1}{6\tau} - \tau\right) \left( \underbrace{\partial_{t}^{(0)}\partial_{t}^{(0)} \sum_{i}f_{i}^{(1)} }_{=0} + \underbrace{\partial_{\alpha}\partial_{t}^{(0)} \sum_{i}c_{i,\alpha}f_{i}^{(1)} }_{=0} + \underbrace{\partial_{\beta}\partial_{t}^{(0)} \sum_{i}c_{i,\beta}f_{i}^{(1)} }_{=0} \right) \\
& + \displaystyle\left(1 - \frac{1}{6\tau} - \tau\right) \left(\partial_{\alpha}\partial_{\beta} \sum_{i}c_{i,\alpha}c_{i,\beta}f_{i}^{(1)} \right)  =  - \displaystyle\frac{1}{\tau} \underbrace{ \displaystyle\sum_{i} f_{i}^{(3)} }_{=0}, \\
(Kn^{(3)}):& \partial_{t}^{(2)} \rho + \displaystyle\left(1 - \frac{1}{6\tau} - \tau\right)\partial_{\alpha}\partial_{\beta} m^{(1)}_{\alpha\beta}  = 0.
\end{array}
\end{split}
\end{equation}

\subsubsection{First\sphinxhyphen{}Order Moment Balance}
\label{\detokenize{Refs/LBM-Study/intro-nse-lbm/Navier-Stokes-Equation:first-order-moment-balance}}
\sphinxAtStartPar
To retrieve the momentum balance equation, we sum the Eq. \eqref{equation:Refs/LBM-Study/intro-nse-lbm/Navier-Stokes-Equation:Chap-Kn-ns-e-NotEq} for \(Kn^{(1)}\), \(Kn^{(2)}\), and \(Kn^{(3)}\) over \(\sum_{i=0}c_{i,\alpha} \):
\begin{equation}\label{equation:Refs/LBM-Study/intro-nse-lbm/Navier-Stokes-Equation:ns-e-Kn1-B1-all}
\begin{split}
\begin{array}{l}
(Kn^{(1)}): \displaystyle\sum_{i}\left( \partial_{t}^{(0)} + c_{i,\beta}\partial_{\beta} \right)f_{i}^{(0)}c_{i,\alpha}= - \underbrace{\displaystyle\sum_{i}\left(\displaystyle\frac{ f_{i}^{(1)} }{ \tau } \right)c_{i,\alpha} }_{=0},\\
(Kn^{(1)}): \partial_{t}^{(0)} \left( \rho u_{\alpha} \right) + \partial_{\beta} \left(\rho u_{\alpha}u_{\beta} + \rho c_{s}^{2}\delta_{\alpha\beta} \right) =0,
\end{array}
\end{split}
\end{equation}\begin{equation}\label{equation:Refs/LBM-Study/intro-nse-lbm/Navier-Stokes-Equation:ns-e-Kn2-B1-all}
\begin{split}
\begin{array}{ll}
(Kn^{(2)}):& \displaystyle\sum_{i}\left[ \partial_{t}^{(1)} f_{i}^{(0)} + \displaystyle\left(1 - \frac{1}{2\tau}\right)\left( \partial_{t}^{(0)} + c_{i,\beta}\partial_{\beta} \right) f_{i}^{(1)} \right]c_{i,\alpha} =  - \displaystyle\sum_{i}\left(\displaystyle\frac{ f_{i}^{(2)} }{ \tau } \right)c_{i,\alpha}, \nonumber \\
(Kn^{(2)}):& \partial_t^{(1)} \displaystyle\sum_i f_i^{(0)}c_{i,\alpha}  + \underbrace{\left(1-\frac{1}{2\tau}\right) \partial_t^{(0)} \sum_i f_i^{(1)}c_{i,\alpha} }_{=0} + \left(1-\frac{1}{2\tau}\right) \partial_\beta \sum_i f_i^{(1)}c_{i,\alpha}c_{i,\beta}  =\underbrace{-\frac{1}{\tau} \sum_i f_i^{(2)}c_{i,\alpha} }_{=0} \\
(Kn^{(2)}):& \partial_{t}^{(1)} \left(\rho u_{\alpha}\right) + \left(1-\displaystyle\frac{1}{2\tau}\right) \partial_\beta m^{(1)}_{\alpha\beta}  = 0,
\end{array}
\end{split}
\end{equation}
\sphinxAtStartPar
Extending the analysis up to \(Kn^{(3)}\)
\begin{equation}\label{equation:Refs/LBM-Study/intro-nse-lbm/Navier-Stokes-Equation:ns-e-Kn3-B1-all}
\begin{split}
\begin{array}{ll}
(Kn^{(3)}):& \displaystyle\sum_{i}\left[ \partial_{t}^{(2)} f_{i}^{(0)} + \left(2-\displaystyle\frac{1}{\tau}\right)\partial_{t}^{(1)} f_{i}^{(1)} + \displaystyle\left(1 - \frac{1}{6\tau} - \tau\right) \left( \partial_{t}^{(0)} + c_{i,\beta}\partial_{\beta} \right)^{2} f_{i}^{(1)} \right]c_{i,\alpha} =  - \displaystyle\sum_{i}\left[\displaystyle\frac{f_{i}^{(3)}}{\tau} \right]c_{i,\alpha}, \\
(Kn^{(3)}):&  \partial_{t}^{(2)} \displaystyle\sum_{i}f_{i}^{(0)}c_{i,\alpha} + \underbrace{\left(2-\displaystyle\frac{1}{\tau}\right)\partial_{t}^{(1)} \sum_{i}f_{i}^{(1)}c_{i,\alpha} }_{=0} + \displaystyle\left(1 - \frac{1}{6\tau} - \tau\right) \left( \underbrace{\partial_{t}^{(0)}\partial_{t}^{(0)} \sum_{i}f_{i}^{(1)}c_{i,\alpha} }_{=0} \right)\\
& + \displaystyle\left(1 - \frac{1}{6\tau} - \tau\right) \left(\partial_{t}^{(0)}\partial_{\beta} \sum_{i}f_{i}^{(1)}c_{i,\alpha}c_{i,\beta} + \partial_{t}^{(0)}\partial_{\gamma} \sum_{i}f_{i}^{(1)}c_{i,\alpha}c_{i,\gamma} + \partial_{\beta}\partial_{\gamma} \sum_{i}c_{i,\alpha}c_{i,\beta}c_{i,\gamma}f_{i}^{(1)} \right)  =  - \displaystyle\frac{1}{\tau} \underbrace{ \displaystyle\sum_{i} f_{i}^{(3)} }_{=0}, \\
(Kn^{(3)}):& \partial_{t}^{(2)} \left(\rho u_{\alpha}\right) + \displaystyle\left(1 - \frac{1}{6\tau} - \tau\right) \left( \partial_{t}^{(0)}\partial_{\beta} m^{(1)}_{\alpha\beta} + \partial_{t}^{(0)}\partial_{\gamma} m^{(1)}_{\alpha\gamma} + \partial_{\beta}\partial_{\gamma} m^{(1)}_{\alpha\beta\gamma} \right) = 0.
\end{array}
\end{split}
\end{equation}

\subsubsection{Second\sphinxhyphen{}Order Non\sphinxhyphen{}Equilibrium Moment}
\label{\detokenize{Refs/LBM-Study/intro-nse-lbm/Navier-Stokes-Equation:second-order-non-equilibrium-moment}}
\sphinxAtStartPar
To determine the unknow term \(m^{(1)}_{\alpha\beta}\), we sum the Eq. \eqref{equation:Refs/LBM-Study/intro-nse-lbm/Navier-Stokes-Equation:Chap-Kn-ns-e-NotEq} for \(Kn^{(1)}\) over the second\sphinxhyphen{}order moment:
\begin{equation*}
\begin{split}
\begin{array}{l}
(Kn^{(1)}): \displaystyle\sum_{i}\left( \partial_{t}^{(0)} + c_{i,\gamma}\partial_{\gamma} \right)f_{i}^{(0)}c_{i,\alpha}c_{i,\beta}= - \displaystyle\sum_{i}\left(\displaystyle\frac{ f_{i}^{(1)} }{ \tau } c_{i,\alpha}c_{i,\beta} \right),\\
(Kn^{(1)}): \partial_{t}^{(0)}  \left(\rho u_{\alpha}u_{\beta} + \rho c_{s}^{2}\delta_{\alpha\beta} \right) + \partial_{\gamma} \left(c_{s}^{2}\rho \left(u_{\alpha}\delta_{\beta\gamma} + u_{\beta}\delta_{\alpha\gamma} + u_{\gamma}\delta_{\alpha\beta} \right) \right) = - \displaystyle\frac{m^{(1)}_{\alpha\beta}}{\tau},
\end{array}
\end{split}
\end{equation*}
\sphinxAtStartPar
notice that third\sphinxhyphen{}order moment for the lattice D2Q9 is not correct recovered. Expanding the terms and using Eqs. \eqref{equation:Refs/LBM-Study/intro-nse-lbm/Navier-Stokes-Equation:ns-e-Kn1-B0-all} and \eqref{equation:Refs/LBM-Study/intro-nse-lbm/Navier-Stokes-Equation:ns-e-Kn1-B1-all} for numerical manipulation:
\begin{equation}\label{equation:Refs/LBM-Study/intro-nse-lbm/Navier-Stokes-Equation:ns-e-Kn1-B2-all}
\begin{split}
\small
\begin{array}{ll}
(Kn^{(1)}):& \partial_{t}^{(0)}  \left(\rho u_{\alpha}u_{\beta} + \rho c_{s}^{2}\delta_{\alpha\beta} \right) + \partial_{\gamma} \left(c_{s}^{2}\rho \left(u_{\alpha}\delta_{\beta\gamma} + u_{\beta}\delta_{\alpha\gamma} + u_{\gamma}\delta_{\alpha\beta} \right) \right) = - \displaystyle\frac{m^{(1)}_{\alpha\beta}}{\tau},\\
(Kn^{(1)}):& \partial_{t}^{(0)}  \left(\rho c_{s}^{2}\delta_{\alpha\beta} \right) + \partial_{t}^{(0)}  \left(\rho u_{\alpha}u_{\beta} \right) +  c_{s}^{2} \left(\partial_{\gamma} (\rho u_{\alpha})\delta_{\beta\gamma} + \partial_{\gamma}(\rho u_{\beta} )\delta_{\alpha\gamma} \right) + \partial_{\gamma} \left(c_{s}^{2}\rho u_{\gamma}\delta_{\alpha\beta} \right) = - \displaystyle\frac{m^{(1)}_{\alpha\beta}}{\tau},\\
(Kn^{(1)}):& c_{s}^{2} \partial_{t}^{(0)}  \left(\rho\right) \delta_{\alpha\beta} + \rho u_{\alpha} \partial_{t}^{(0)} \left(u_{\beta} \right) + u_{\beta}\partial_{t}^{(0)}  \left(\rho u_{\alpha}\right) +  c_{s}^{2} \left(\partial_{\beta} (\rho u_{\alpha}) + \partial_{\alpha} (\rho u_{\beta}) \right) + c_{s}^{2} \partial_{\gamma} \left(\rho u_{\gamma} \right) \delta_{\alpha\beta} = - \displaystyle\frac{m^{(1)}_{\alpha\beta}}{\tau},\\
(Kn^{(1)}):& -\textcolor{red}{c_{s}^{2} \partial_{\gamma} \left( \rho u_{\gamma} \right)\delta_{\alpha\beta}} + 
            \rho u_{\alpha} \partial_{t}^{(0)} \left(u_{\beta} \right) - 
            u_{\beta} \partial_{\gamma} \left(\rho u_{\alpha}u_{\gamma} + \rho c_{s}^{2}\delta_{\alpha\gamma} \right) +  
            c_{s}^{2} \left(\partial_{\beta} (\rho u_{\alpha}) + \partial_{\alpha} (\rho u_{\beta}) \right) + 
            \textcolor{red}{c_{s}^{2} \partial_{\gamma} \left(\rho u_{\gamma} \right) \delta_{\alpha\beta}} = 
            - \displaystyle\frac{m^{(1)}_{\alpha\beta}}{\tau},\\
(Kn^{(1)}):& u_{\alpha} \left( -u_{\beta}\partial_{t}^{(0)} \rho  - \partial_{\gamma} \left(\rho u_{\beta}u_{\gamma} + \rho c_{s}^{2}\delta_{\beta\gamma} \right)\right) - 
            u_{\beta} \partial_{\gamma} \left(\rho u_{\alpha}u_{\gamma} + \rho c_{s}^{2}\delta_{\alpha\gamma} \right) +  
            c_{s}^{2} \left(\partial_{\beta} (\rho u_{\alpha}) + \partial_{\alpha} (\rho u_{\beta}) \right) = 
            - \displaystyle\frac{m^{(1)}_{\alpha\beta}}{\tau},\\
(Kn^{(1)}):& u_{\alpha} \left( u_{\beta}\partial_{\gamma} (\rho u_{\gamma})  - \partial_{\gamma} \left(\rho u_{\beta}u_{\gamma} + \rho c_{s}^{2}\delta_{\beta\gamma} \right)\right) - 
            u_{\beta} \partial_{\gamma} \left(\rho u_{\alpha}u_{\gamma} + \rho c_{s}^{2}\delta_{\alpha\gamma} \right) +  
            c_{s}^{2} \left(\partial_{\beta} (\rho u_{\alpha}) + \partial_{\alpha} (\rho u_{\beta}) \right) = 
            - \displaystyle\frac{m^{(1)}_{\alpha\beta}}{\tau},\\
(Kn^{(1)}):& u_{\alpha} \left( -\rho u_{\gamma}\partial_{\gamma} u_{\beta}  - c_{s}^{2}\partial_{\beta} \rho \right) - 
            u_{\beta} \partial_{\gamma} \left(\rho u_{\alpha}u_{\gamma} + \rho c_{s}^{2}\delta_{\alpha\gamma} \right) +  
            c_{s}^{2} \left(\partial_{\beta} (\rho u_{\alpha}) + \partial_{\alpha} (\rho u_{\beta}) \right) = 
            - \displaystyle\frac{m^{(1)}_{\alpha\beta}}{\tau},\\
(Kn^{(1)}):&   -\rho u_{\alpha}u_{\gamma}\partial_{\gamma} u_{\beta}  - c_{s}^{2}u_{\alpha}\partial_{\beta} \rho - 
            u_{\beta} \partial_{\gamma}\rho u_{\alpha}u_{\gamma} - u_{\beta}c_{s}^{2} \partial_{\alpha}\rho  +  
            c_{s}^{2} \left(\partial_{\beta} (\rho u_{\alpha}) + \partial_{\alpha} (\rho u_{\beta}) \right) = 
            - \displaystyle\frac{m^{(1)}_{\alpha\beta}}{\tau},\\
(Kn^{(1)}):&   -\partial_{\gamma} (\rho u_{\alpha}u_{\beta}u_{\gamma}) 
            - c_{s}^{2} (u_{\alpha}\partial_{\beta} \rho + u_{\beta} \partial_{\alpha}\rho)  +  
            c_{s}^{2} \left(\partial_{\beta} (\rho u_{\alpha}) + \partial_{\alpha} (\rho u_{\beta}) \right) = 
            - \displaystyle\frac{m^{(1)}_{\alpha\beta}}{\tau},\\
(Kn^{(1)}):&   -\partial_{\gamma} (\rho u_{\alpha}u_{\beta}u_{\gamma}) 
            +c_{s}^{2} \rho\left(\partial_{\beta} u_{\alpha} + \partial_{\alpha} u_{\beta} \right) = 
            - \displaystyle\frac{m^{(1)}_{\alpha\beta}}{\tau},  \qquad \rightarrow \qquad  m^{(1)}_{\alpha\beta} = -c_{s}^{2} \tau \rho\left(\partial_{\beta} u_{\alpha} + \partial_{\alpha} u_{\beta} \right) + \tau \partial_{\gamma} (\rho u_{\alpha}u_{\beta}u_{\gamma})
\end{array}
\end{split}
\end{equation}

\subsubsection{Balance Equations}
\label{\detokenize{Refs/LBM-Study/intro-nse-lbm/Navier-Stokes-Equation:balance-equations}}
\sphinxAtStartPar
To recover the mass balance equation, we sum the Eqs. \eqref{equation:Refs/LBM-Study/intro-nse-lbm/Navier-Stokes-Equation:ns-e-Kn1-B0-all}, \eqref{equation:Refs/LBM-Study/intro-nse-lbm/Navier-Stokes-Equation:ns-e-Kn2-B0-all}, and \eqref{equation:Refs/LBM-Study/intro-nse-lbm/Navier-Stokes-Equation:ns-e-Kn3-B0-all}:
\begin{equation*}
\begin{split}
\begin{array}{l}
\partial_{t} \rho + \partial_{\alpha} (\rho u_{\alpha}) + \underbrace{ \left(1-\displaystyle\frac{1}{2\tau}\right) \partial_\beta m^{(1)}_{\alpha\beta} + ... + \mathcal{O}(Kn^{(\geq 3)}) }_{\textrm{Numerical Errors}} =0,\\
\partial_{t} \rho + \partial_{\alpha} (\rho u_{\alpha}) + \underbrace{ \left(\tau-\displaystyle\frac{1}{2}\right) \left(\partial_\beta (-c_{s}^{2} \rho\left(\partial_{\beta} u_{\alpha} + \partial_{\alpha} u_{\beta} \right) + \partial_{\gamma} (\rho u_{\alpha}u_{\beta}u_{\gamma}))\right) + ... + \mathcal{O}(Kn^{(> 3)}) }_{\textrm{Numerical Errors}} =0,
\end{array}
\end{split}
\end{equation*}
\sphinxAtStartPar
Repeting the process for momentum balance, we sum the Eqs. \eqref{equation:Refs/LBM-Study/intro-nse-lbm/Navier-Stokes-Equation:ns-e-Kn1-B1-all}, \eqref{equation:Refs/LBM-Study/intro-nse-lbm/Navier-Stokes-Equation:ns-e-Kn2-B1-all}, \eqref{equation:Refs/LBM-Study/intro-nse-lbm/Navier-Stokes-Equation:ns-e-Kn3-B1-all}, and \eqref{equation:Refs/LBM-Study/intro-nse-lbm/Navier-Stokes-Equation:ns-e-Kn1-B2-all}:
\begin{equation*}
\begin{split}
\begin{array}{l}
\partial_t \rho u_{\alpha} + \partial_{\beta}\left(\rho u_{\alpha}\rho u_{\beta}\right) =  \partial_{\beta}\left( -\rho c_{s}^{2} \delta_{\alpha\beta} + \left(\tau-\displaystyle\frac{1}{2}\right) \rho c_{s}^{2}\left(\partial_{\alpha}u_{\beta}+\partial_{\beta}u_{\alpha}\right) \right) \\
 \underbrace{ -\left(\tau-\displaystyle\frac{1}{2}\right)  \partial_{\gamma} (\rho u_{\alpha}u_{\beta}u_{\gamma}) - \displaystyle\left(1 - \frac{1}{6\tau} - \tau\right) \left( \partial_{t}^{(0)}\partial_{\beta} m^{(1)}_{\alpha\beta} + \partial_{t}^{(0)}\partial_{\gamma} m^{(1)}_{\alpha\gamma} + \partial_{\beta}\partial_{\gamma} m^{(1)}_{\alpha\beta\gamma} \right) + ... + \mathcal{O}(Kn^{(> 3)}) }_{\textrm{Numerical Errors}} =0,\\
\partial_t \rho u_{\alpha} + \partial_{\beta}\left(\rho u_{\alpha}\rho u_{\beta}\right) =  \partial_{\beta}\left( -p \delta_{\alpha\beta} + \mu\left(\partial_{\alpha}u_{\beta}+\partial_{\beta}u_{\alpha}\right) \right) \\
\underbrace{ -\left(\tau-\displaystyle\frac{1}{2}\right)  \partial_{\gamma} (\rho u_{\alpha}u_{\beta}u_{\gamma}) - \displaystyle\left(1 - \frac{1}{6\tau} - \tau\right) \left( \partial_{t}^{(0)}\partial_{\beta} m^{(1)}_{\alpha\beta} + \partial_{t}^{(0)}\partial_{\gamma} m^{(1)}_{\alpha\gamma} + \partial_{\beta}\partial_{\gamma} m^{(1)}_{\alpha\beta\gamma} \right) + ... + \mathcal{O}(Kn^{(> 3)}) }_{\textrm{Numerical Errors}} =0,\\
\end{array}
\end{split}
\end{equation*}
\sphinxAtStartPar
where \(\mu=\rho c_{s}^{2}(\tau-1/2)\).


\subsection{Lattice Direction Moments: (Multi\sphinxhyphen{}Speed Lattices) D2V17 Velocity\sphinxhyphen{}Space Discetization}
\label{\detokenize{Refs/LBM-Study/intro-nse-lbm/Navier-Stokes-Equation:lattice-direction-moments-multi-speed-lattices-d2v17-velocity-space-discetization}}
\begin{sphinxuseclass}{cell}
\begin{sphinxuseclass}{tag_hide-input}\begin{sphinxVerbatimOutput}

\begin{sphinxuseclass}{cell_output}
\noindent\sphinxincludegraphics{{3df92e67d5bd37c6948ff4e9fe72bdfd03512ea32d2a346073d7ba79cf469d1a}.png}

\end{sphinxuseclass}\end{sphinxVerbatimOutput}

\end{sphinxuseclass}
\end{sphinxuseclass}

\subsubsection{Raw Moments}
\label{\detokenize{Refs/LBM-Study/intro-nse-lbm/Navier-Stokes-Equation:id3}}
\begin{sphinxuseclass}{cell}
\begin{sphinxuseclass}{tag_hide-input}\begin{sphinxVerbatimOutput}

\begin{sphinxuseclass}{cell_output}\begin{equation*}
\begin{split}\displaystyle \begin{gathered}\underbrace{\sum_{i=0} f_{i}^{eq} =\sum_{i=0} f_{i} }_{\textrm{Zero-Order Moment}} =\rho,\quad \underbrace{\sum_{i=0} f_{i}^{eq}c_{i,x} }_{\textrm{x-First-Order Moment}} =\rho u^{*}_{x},\quad \underbrace{\sum_{i=0} f_{i}^{eq}c_{i,y} }_{\textrm{y-First-Order Moment}} =\rho u^{*}_{y}\end{gathered}\end{split}
\end{equation*}\begin{equation*}
\begin{split}\displaystyle \begin{gathered}\underbrace{\sum_{i=0} f_{i}^{eq} c_{i,x}c_{i,x} }_{\textrm{xx-Second-Order Moment}} =\rho c_{s}^{2} + \rho \left(u^{*}_{x}\right)^{2},\quad \underbrace{\sum_{i=0} f_{i}^{eq} c_{i,x}c_{i,y} }_{\textrm{xy-Second-Order Moment}} =\rho u^{*}_{x} u^{*}_{y},\quad \underbrace{\sum_{i=0} f_{i}^{eq} c_{i,y}c_{i,y} }_{\textrm{yy-Second-Order Moment}} =\rho c_{s}^{2} + \rho \left(u^{*}_{y}\right)^{2}\end{gathered}\end{split}
\end{equation*}\begin{equation*}
\begin{split}\displaystyle \begin{gathered}\underbrace{\sum_{i=0} f_{i}^{eq} c_{i,x}c_{i,x}c_{i,x} }_{\textrm{xxx-Third-Order Moment}} =3 \rho c_{s}^{2} u^{*}_{x} + \rho \left(u^{*}_{x}\right)^{3},\quad \underbrace{\sum_{i=0} f_{i}^{eq} c_{i,x}c_{i,x}c_{i,y} }_{\textrm{xxy-Third-Order Moment}} =\rho c_{s}^{2} u^{*}_{y} + \rho \left(u^{*}_{x}\right)^{2} u^{*}_{y},\\[10pt]\underbrace{\sum_{i=0} f_{i}^{eq} c_{i,x}c_{i,y}c_{i,y} }_{\textrm{xyy-Third-Order Moment}} =\rho c_{s}^{2} u^{*}_{x} + \rho u^{*}_{x} \left(u^{*}_{y}\right)^{2},\quad \underbrace{\sum_{i=0} f_{i}^{eq} c_{i,y}c_{i,y}c_{i,y} }_{\textrm{yyy-Third-Order Moment}} =3 \rho c_{s}^{2} u^{*}_{y} + \rho \left(u^{*}_{y}\right)^{3}\end{gathered}\end{split}
\end{equation*}
\end{sphinxuseclass}\end{sphinxVerbatimOutput}

\end{sphinxuseclass}
\end{sphinxuseclass}

\subsubsection{Hermite Moments}
\label{\detokenize{Refs/LBM-Study/intro-nse-lbm/Navier-Stokes-Equation:id4}}
\begin{sphinxuseclass}{cell}
\begin{sphinxuseclass}{tag_hide-input}\begin{sphinxVerbatimOutput}

\begin{sphinxuseclass}{cell_output}\begin{equation*}
\begin{split}\displaystyle \begin{gathered}\underbrace{\sum_{i=0} f_{i}^{eq} =\sum_{i=0} f_{i} }_{\textrm{Zero-Order Hermite Moment}} =\rho,\quad \underbrace{\sum_{i=0} f_{i}^{eq}c_{i,x} }_{\textrm{x-First-Order Hermite Moment}} =\rho u^{*}_{x},\quad \underbrace{\sum_{i=0} f_{i}^{eq}c_{i,y} }_{\textrm{y-First-Order Hermite Moment}} =\rho u^{*}_{y}\end{gathered}\end{split}
\end{equation*}\begin{equation*}
\begin{split}\displaystyle \begin{gathered}\underbrace{\sum_{i=0} f_{i}^{eq} c_{i,x}c_{i,x} }_{\textrm{xx-Second-Order Moment}} =\rho \left(u^{*}_{x}\right)^{2},\quad \underbrace{\sum_{i=0} f_{i}^{eq} c_{i,x}c_{i,y} }_{\textrm{xy-Second-Order Moment}} =\rho u^{*}_{x} u^{*}_{y},\quad \underbrace{\sum_{i=0} f_{i}^{eq} c_{i,y}c_{i,y} }_{\textrm{yy-Second-Order Moment}} =\rho \left(u^{*}_{y}\right)^{2}\end{gathered}\end{split}
\end{equation*}\begin{equation*}
\begin{split}\displaystyle \begin{gathered}\underbrace{\sum_{i=0} f_{i}^{eq} c_{i,x}c_{i,x}c_{i,x} }_{\textrm{xxx-Third-Order Moment}} =\rho \left(u^{*}_{x}\right)^{3},\quad \underbrace{\sum_{i=0} f_{i}^{eq} c_{i,x}c_{i,x}c_{i,y} }_{\textrm{xxy-Third-Order Moment}} =\rho \left(u^{*}_{x}\right)^{2} u^{*}_{y},\\[10pt]\underbrace{\sum_{i=0} f_{i}^{eq} c_{i,x}c_{i,y}c_{i,y} }_{\textrm{xyy-Third-Order Moment}} =\rho u^{*}_{x} \left(u^{*}_{y}\right)^{2},\quad \underbrace{\sum_{i=0} f_{i}^{eq} c_{i,y}c_{i,y}c_{i,y} }_{\textrm{yyy-Third-Order Moment}} =\rho \left(u^{*}_{y}\right)^{3}\end{gathered}\end{split}
\end{equation*}
\end{sphinxuseclass}\end{sphinxVerbatimOutput}

\end{sphinxuseclass}
\end{sphinxuseclass}
\sphinxAtStartPar
Doing the Chapman\sphinxhyphen{}Enskog analysis for the lattice D2V17, we obtain the same Eqs. \eqref{equation:Refs/LBM-Study/intro-nse-lbm/Navier-Stokes-Equation:ns-e-Kn1-B0-all}, \eqref{equation:Refs/LBM-Study/intro-nse-lbm/Navier-Stokes-Equation:ns-e-Kn2-B0-all}, \eqref{equation:Refs/LBM-Study/intro-nse-lbm/Navier-Stokes-Equation:ns-e-Kn3-B0-all}, \eqref{equation:Refs/LBM-Study/intro-nse-lbm/Navier-Stokes-Equation:ns-e-Kn1-B1-all}, \eqref{equation:Refs/LBM-Study/intro-nse-lbm/Navier-Stokes-Equation:ns-e-Kn2-B1-all}, and \eqref{equation:Refs/LBM-Study/intro-nse-lbm/Navier-Stokes-Equation:ns-e-Kn3-B1-all}. However, the calculation of the second\sphinxhyphen{}order moment of non\sphinxhyphen{}equilibrium \( m^{(1)}_{\alpha\beta}\) will change due to the D2V17 be able to correctly recover the third\sphinxhyphen{}order moments:
\begin{equation}\label{equation:Refs/LBM-Study/intro-nse-lbm/Navier-Stokes-Equation:ns-e-Kn1-B2-all-D2V17}
\begin{split}
\scriptsize
\begin{array}{ll}
(Kn^{(1)}):& \displaystyle\sum_{i}\left( \partial_{t}^{(0)} + c_{i,\gamma}\partial_{\gamma} \right)f_{i}^{(0)}c_{i,\alpha}c_{i,\beta}= - \displaystyle\sum_{i}\left(\displaystyle\frac{ f_{i}^{(1)} }{ \tau } c_{i,\alpha}c_{i,\beta} \right),\\
(Kn^{(1)}):& \partial_{t}^{(0)}  \left(\rho u_{\alpha}u_{\beta} + \rho c_{s}^{2}\delta_{\alpha\beta} \right) + \partial_{\gamma} \left(c_{s}^{2}\rho \left(u_{\alpha}\delta_{\beta\gamma} + u_{\beta}\delta_{\alpha\gamma} + u_{\gamma}\delta_{\alpha\beta} \right) \right) = - \displaystyle\frac{m^{(1)}_{\alpha\beta}}{\tau}, \\
(Kn^{(1)}):& \partial_{t}^{(0)}  \left(\rho u_{\alpha}u_{\beta} + \rho c_{s}^{2}\delta_{\alpha\beta} \right) + \partial_{\gamma} \left(c_{s}^{2}\rho \left(u_{\alpha}\delta_{\beta\gamma} + u_{\beta}\delta_{\alpha\gamma} + u_{\gamma}\delta_{\alpha\beta} \right) \right) + \partial_{\gamma}\left( u_{\alpha} u_{\beta} u_{\gamma}\right) = - \displaystyle\frac{m^{(1)}_{\alpha\beta}}{\tau},\\
(Kn^{(1)}):& \partial_{t}^{(0)}  \left(\rho c_{s}^{2}\delta_{\alpha\beta} \right) + \partial_{t}^{(0)}  \left(\rho u_{\alpha}u_{\beta} \right) +  c_{s}^{2} \left(\partial_{\gamma} (\rho u_{\alpha})\delta_{\beta\gamma} + \partial_{\gamma}(\rho u_{\beta} )\delta_{\alpha\gamma} \right) + \partial_{\gamma} \left(c_{s}^{2}\rho u_{\gamma}\delta_{\alpha\beta} \right) + \partial_{\gamma}\left( u_{\alpha} u_{\beta} u_{\gamma}\right) = - \displaystyle\frac{m^{(1)}_{\alpha\beta}}{\tau},\\
(Kn^{(1)}):& c_{s}^{2} \partial_{t}^{(0)}  \left(\rho\right) \delta_{\alpha\beta} + \rho u_{\alpha} \partial_{t}^{(0)} \left(u_{\beta} \right) + u_{\beta}\partial_{t}^{(0)}  \left(\rho u_{\alpha}\right) +  c_{s}^{2} \left(\partial_{\beta} (\rho u_{\alpha}) + \partial_{\alpha} (\rho u_{\beta}) \right) + c_{s}^{2} \partial_{\gamma} \left(\rho u_{\gamma} \right) \delta_{\alpha\beta} + \partial_{\gamma}\left( u_{\alpha} u_{\beta} u_{\gamma}\right) = - \displaystyle\frac{m^{(1)}_{\alpha\beta}}{\tau},\\
(Kn^{(1)}):& -\textcolor{red}{c_{s}^{2} \partial_{\gamma} \left( \rho u_{\gamma} \right)\delta_{\alpha\beta}} + 
            \rho u_{\alpha} \partial_{t}^{(0)} \left(u_{\beta} \right) - 
            u_{\beta} \partial_{\gamma} \left(\rho u_{\alpha}u_{\gamma} + \rho c_{s}^{2}\delta_{\alpha\gamma} \right) +  
            c_{s}^{2} \left(\partial_{\beta} (\rho u_{\alpha}) + \partial_{\alpha} (\rho u_{\beta}) \right) + 
            \textcolor{red}{c_{s}^{2} \partial_{\gamma} \left(\rho u_{\gamma} \right) \delta_{\alpha\beta}} + \partial_{\gamma}\left( u_{\alpha} u_{\beta} u_{\gamma}\right) = 
            - \displaystyle\frac{m^{(1)}_{\alpha\beta}}{\tau},\\
(Kn^{(1)}):& u_{\alpha} \left( -u_{\beta}\partial_{t}^{(0)} \rho  - \partial_{\gamma} \left(\rho u_{\beta}u_{\gamma} + \rho c_{s}^{2}\delta_{\beta\gamma} \right)\right) - 
            u_{\beta} \partial_{\gamma} \left(\rho u_{\alpha}u_{\gamma} + \rho c_{s}^{2}\delta_{\alpha\gamma} \right) +  
            c_{s}^{2} \left(\partial_{\beta} (\rho u_{\alpha}) + \partial_{\alpha} (\rho u_{\beta}) \right) 
            + \partial_{\gamma}\left( u_{\alpha} u_{\beta} u_{\gamma}\right) = 
            - \displaystyle\frac{m^{(1)}_{\alpha\beta}}{\tau},\\
(Kn^{(1)}):& u_{\alpha} \left( u_{\beta}\partial_{\gamma} (\rho u_{\gamma})  - \partial_{\gamma} \left(\rho u_{\beta}u_{\gamma} + \rho c_{s}^{2}\delta_{\beta\gamma} \right)\right) - 
            u_{\beta} \partial_{\gamma} \left(\rho u_{\alpha}u_{\gamma} + \rho c_{s}^{2}\delta_{\alpha\gamma} \right) +  
            c_{s}^{2} \left(\partial_{\beta} (\rho u_{\alpha}) + \partial_{\alpha} (\rho u_{\beta}) \right) 
            + \partial_{\gamma}\left( u_{\alpha} u_{\beta} u_{\gamma}\right) = 
            - \displaystyle\frac{m^{(1)}_{\alpha\beta}}{\tau},\\
(Kn^{(1)}):& u_{\alpha} \left( -\rho u_{\gamma}\partial_{\gamma} u_{\beta}  - c_{s}^{2}\partial_{\beta} \rho \right) - 
            u_{\beta} \partial_{\gamma} \left(\rho u_{\alpha}u_{\gamma} + \rho c_{s}^{2}\delta_{\alpha\gamma} \right) +  
            c_{s}^{2} \left(\partial_{\beta} (\rho u_{\alpha}) + \partial_{\alpha} (\rho u_{\beta}) \right) 
            + \partial_{\gamma}\left( u_{\alpha} u_{\beta} u_{\gamma}\right) = 
            - \displaystyle\frac{m^{(1)}_{\alpha\beta}}{\tau},\\
(Kn^{(1)}):&   -\rho u_{\alpha}u_{\gamma}\partial_{\gamma} u_{\beta}  - c_{s}^{2}u_{\alpha}\partial_{\beta} \rho - 
            u_{\beta} \partial_{\gamma}\rho u_{\alpha}u_{\gamma} - u_{\beta}c_{s}^{2} \partial_{\alpha}\rho  +  
            c_{s}^{2} \left(\partial_{\beta} (\rho u_{\alpha}) + \partial_{\alpha} (\rho u_{\beta}) \right) 
            + \partial_{\gamma}\left( u_{\alpha} u_{\beta} u_{\gamma}\right) = 
            - \displaystyle\frac{m^{(1)}_{\alpha\beta}}{\tau},\\
(Kn^{(1)}):&   \textcolor{red}{-\partial_{\gamma} (\rho u_{\alpha}u_{\beta}u_{\gamma})} 
            - c_{s}^{2} (u_{\alpha}\partial_{\beta} \rho + u_{\beta} \partial_{\alpha}\rho)  +  
            c_{s}^{2} \left(\partial_{\beta} (\rho u_{\alpha}) + \partial_{\alpha} (\rho u_{\beta}) \right) 
            \textcolor{red}{+ \partial_{\gamma}\left( u_{\alpha} u_{\beta} u_{\gamma}\right)} = 
            - \displaystyle\frac{m^{(1)}_{\alpha\beta}}{\tau},\\
(Kn^{(1)}):&  c_{s}^{2} \rho\left(\partial_{\beta} u_{\alpha} + \partial_{\alpha} u_{\beta} \right) = 
            - \displaystyle\frac{m^{(1)}_{\alpha\beta}}{\tau},  \qquad \rightarrow \qquad  m^{(1)}_{\alpha\beta} = -c_{s}^{2} \tau \rho\left(\partial_{\beta} u_{\alpha} + \partial_{\alpha} u_{\beta} \right)
\end{array}
\end{split}
\end{equation}
\sphinxAtStartPar
Repeting the process for momentum balance, we sum the Eqs. \eqref{equation:Refs/LBM-Study/intro-nse-lbm/Navier-Stokes-Equation:ns-e-Kn1-B1-all}, \eqref{equation:Refs/LBM-Study/intro-nse-lbm/Navier-Stokes-Equation:ns-e-Kn2-B1-all}, \eqref{equation:Refs/LBM-Study/intro-nse-lbm/Navier-Stokes-Equation:ns-e-Kn3-B1-all}, and \eqref{equation:Refs/LBM-Study/intro-nse-lbm/Navier-Stokes-Equation:ns-e-Kn1-B2-all-D2V17}:
\begin{equation*}
\begin{split}
\begin{array}{l}
\partial_t \rho u_{\alpha} + \partial_{\beta}\left(\rho u_{\alpha}\rho u_{\beta}\right) =  \partial_{\beta}\left( -\rho c_{s}^{2} \delta_{\alpha\beta} + \left(\tau-\displaystyle\frac{1}{2}\right) \rho c_{s}^{2}\left(\partial_{\alpha}u_{\beta}+\partial_{\beta}u_{\alpha}\right) \right) \\
 \underbrace{- \displaystyle\left(1 - \frac{1}{6\tau} - \tau\right) \left( \partial_{t}^{(0)}\partial_{\beta} m^{(1)}_{\alpha\beta} + \partial_{t}^{(0)}\partial_{\gamma} m^{(1)}_{\alpha\gamma} + \partial_{\beta}\partial_{\gamma} m^{(1)}_{\alpha\beta\gamma} \right) + ... + \mathcal{O}(Kn^{(> 3)}) }_{\textrm{Numerical Errors}} =0,\\
\partial_t \rho u_{\alpha} + \partial_{\beta}\left(\rho u_{\alpha}\rho u_{\beta}\right) =  \partial_{\beta}\left( -p \delta_{\alpha\beta} + \mu\left(\partial_{\alpha}u_{\beta}+\partial_{\beta}u_{\alpha}\right) \right) \\
\underbrace{- \displaystyle\left(1 - \frac{1}{6\tau} - \tau\right) \left( \partial_{t}^{(0)}\partial_{\beta} m^{(1)}_{\alpha\beta} + \partial_{t}^{(0)}\partial_{\gamma} m^{(1)}_{\alpha\gamma} + \partial_{\beta}\partial_{\gamma} m^{(1)}_{\alpha\beta\gamma} \right) + ... + \mathcal{O}(Kn^{(> 3)}) }_{\textrm{Numerical Errors}} =0,\\
\end{array}
\end{split}
\end{equation*}
\sphinxAtStartPar
where \(\mu=\rho c_{s}^{2}(\tau-1/2)\).


\section{Tutorials}
\label{\detokenize{Refs/LBM-Study/intro-nse-lbm/Navier-Stokes-Equation:tutorials}}
\sphinxstepscope


\part{Finite Difference Structure of Lattice Boltzmann Method}

\sphinxstepscope


\chapter{Intoduction to Difference Finite Scheme derived from Lattice Boltzmann Method: Diffusive Equation}
\label{\detokenize{Refs/LBM-Study/FDM-Belloti/FDM-LBM-Diffusive:intoduction-to-difference-finite-scheme-derived-from-lattice-boltzmann-method-diffusive-equation}}\label{\detokenize{Refs/LBM-Study/FDM-Belloti/FDM-LBM-Diffusive::doc}}
\sphinxAtStartPar
Due to the similarities in the grid structures employed by the lattice Boltzmann method and finite\sphinxhyphen{}difference schemes, several studies have investigated the direct relationships between these approaches and their mathematical interpretations {[}\hyperlink{cite.Refs/Citations:id5}{Anc94}, \hyperlink{cite.Refs/Citations:id11}{BGM22}, \hyperlink{cite.Refs/Citations:id8}{CCS23}, \hyperlink{cite.Refs/Citations:id7}{DLOT17}, \hyperlink{cite.Refs/Citations:id6}{Jun01}, \hyperlink{cite.Refs/Citations:id4}{Sug10}{]}. In this section, through mathematical manipulation of the lattice Boltzmann equations with single relaxation time, we derive a finite\sphinxhyphen{}difference interpretation for one of the simplest transport problems, involving purely diffusive transport of a scalar quantity. The present lattice Boltzmann formulation follows the analysis presented by Suga {[}\hyperlink{cite.Refs/Citations:id4}{Sug10}{]}.


\section{Diffusive Equation Interpretation with Single Relaxation Time}
\label{\detokenize{Refs/LBM-Study/FDM-Belloti/FDM-LBM-Diffusive:diffusive-equation-interpretation-with-single-relaxation-time}}
\sphinxAtStartPar
Given the LBM formulation for the diffusive equation:
\begin{equation}\label{equation:Refs/LBM-Study/FDM-Belloti/FDM-LBM-Diffusive:LB-fdm-df-Eq}
\begin{split}
f_i( x_{\alpha} + e_{i,\alpha} \delta t, t+\delta t) = f_i(x_{\alpha}, t)  -\left( \frac{ f_{i} - f_{i}^{eq} }{ \tau } \right) \qquad \rightarrow \qquad f_i( x_{\alpha} + e_{i,\alpha} \delta t, t+\delta t) = f_i(x_{\alpha}, t)  - \omega\left( f_{i} - f_{i}^{eq} \right), 
\end{split}
\end{equation}
\sphinxAtStartPar
where the equilibrium distribution function is defined by
\begin{equation}\label{equation:Refs/LBM-Study/FDM-Belloti/FDM-LBM-Diffusive:df-fdm-feq-Eq}
\begin{split}
f^{eq}_{i} = w_{i}\phi(x_{\alpha}, t),
\end{split}
\end{equation}
\sphinxAtStartPar
and the equilibrium moments are given by
\begin{equation}\label{equation:Refs/LBM-Study/FDM-Belloti/FDM-LBM-Diffusive:moments-fdm-fi-df}
\begin{split}
\displaystyle\sum_{i=0} f_{i}^{eq}=\phi, \quad \quad  \displaystyle\sum_{i}f_{i}^{eq} e_{i,\alpha} =0 \quad \quad \textrm{and} \quad \quad  \displaystyle\sum_{i}f_{i}^{eq} e_{i,\alpha} e_{i,\beta} =c_{s}^{2} \phi \delta_{\alpha\beta},
\end{split}
\end{equation}
\sphinxAtStartPar
where \(\phi(x_{\alpha}, t)\) is the scalar term transported by diffusion.




\subsection{FDM Interpretation for the Lattice D1Q3}
\label{\detokenize{Refs/LBM-Study/FDM-Belloti/FDM-LBM-Diffusive:fdm-interpretation-for-the-lattice-d1q3}}
\sphinxAtStartPar
For simplification we redefine the notation for \(f_i( x_{\alpha} + e_{i,\alpha} \delta t, t+\delta t)=f_{i,x+c_{i}}^{t+1}\). Using this notation, Eq. \eqref{equation:Refs/LBM-Study/FDM-Belloti/FDM-LBM-Diffusive:LB-fdm-df-Eq} can be reinterpreted by applying a shift in \(c_i\) to express the distribution function at the pre\sphinxhyphen{}streaming state:
\begin{equation}\label{equation:Refs/LBM-Study/FDM-Belloti/FDM-LBM-Diffusive:pre-fdm-LB-df-Eq}
\begin{split}
f_{i,x}^{t+1}  = f_{i,x-c_{i}}^{t}  - \omega \left( f_{i,x-c_{i}}^{t} - f_{i,x-c_{i}}^{eq,t} \right).
\end{split}
\end{equation}
\sphinxAtStartPar
Summing the Eq. \eqref{equation:Refs/LBM-Study/FDM-Belloti/FDM-LBM-Diffusive:pre-fdm-LB-df-Eq} over all lattice direction and replacing \(f_{i}^{eq}\) by Eq. \eqref{equation:Refs/LBM-Study/FDM-Belloti/FDM-LBM-Diffusive:df-fdm-feq-Eq}, we have:
\begin{equation*}
\begin{split}
\displaystyle\sum_{i}\Big( f_{i,x}^{t+1} \Big) = \displaystyle\sum_{i}\Big( f_{i,x-c_{i}}^{t}  - \omega \left( f_{i,x-c_{i}}^{t} - f_{i,x-c_{i}}^{eq,t} \right) \Big) \qquad \rightarrow  \qquad \phi_{x}^{t+1} = \displaystyle\sum_{i}\Big( (1-\omega) f_{i,x-c_{i}}^{t}  + \omega w_{i}\phi_{x-c_{i}}^{t} \Big),
\end{split}
\end{equation*}
\sphinxAtStartPar
using the Eq. \eqref{equation:Refs/LBM-Study/FDM-Belloti/FDM-LBM-Diffusive:moments-fdm-fi-df} to replace \(f_{i,x-c_{i}}^{t}\) in the above equation:
\begin{equation*}
\begin{split}
\phi_{x}^{t+1} = \displaystyle\sum_{i}\Big( (1-\omega) f_{i,x-c_{i}}^{t}  + \omega w_{i}\phi_{x-c_{i}}^{t} \Big) \qquad \rightarrow  \qquad \phi_{x}^{t+1} = \displaystyle\sum_{i}\Big( (1-\omega + \omega w_{i}) \phi_{x-c_{i}}^{t} - (1-\omega) \displaystyle\sum_{k\neq i} f_{k,x-c_{i}}^{t} \Big).
\end{split}
\end{equation*}
\sphinxAtStartPar
Spliting the terms in the above equation and using Eq. \eqref{equation:Refs/LBM-Study/FDM-Belloti/FDM-LBM-Diffusive:pre-fdm-LB-df-Eq} for \(f_{i,x-c_{i}}^{t}\):
\begin{equation*}
\begin{split}
\phi_{x}^{t+1} = \displaystyle\sum_{i}\Big( (1-\omega + \omega w_{i}) \phi_{x-c_{i}}^{t}\Big) - (1-\omega)\displaystyle\sum_{i}\left( \displaystyle\sum_{k\neq i} f_{k,x-c_{i}}^{t} \right),
\end{split}
\end{equation*}\begin{equation*}
\begin{split}
\phi_{x}^{t+1} = \displaystyle\sum_{i}\Big( (1-\omega + \omega w_{i}) \phi_{x-c_{i}}^{t}\Big) - (1-\omega)\displaystyle\sum_{i}\Bigg( \displaystyle\sum_{k\neq i} \left( f_{k,x-c_{i}-c_{k}}^{t-1}  - \omega \left( f_{k,x-c_{i}-c_{k}}^{t-1} - f_{k,x-c_{i}-c_{k}}^{eq,t-1} \right) \right) \Bigg),
\end{split}
\end{equation*}
\sphinxAtStartPar
doing the same initial procedure to the terms \(f_{k,x-c_{i}}^{t-1}\):
\begin{equation*}
\begin{split}
\phi_{x}^{t+1} = \displaystyle\sum_{i}\Big( (1-\omega + \omega w_{i}) \phi_{x-c_{i}}^{t}\Big) -(1-\omega)\displaystyle\sum_{i}\Bigg( \displaystyle\sum_{k\neq i} \left( f_{k,x-c_{i}-c_{k}}^{t-1}  - \omega \left( f_{k,x-c_{i}-c_{k}}^{t-1} - f_{k,x-c_{i}-c_{k}}^{eq,t-1} \right) \right) \Bigg)
\end{split}
\end{equation*}\begin{equation*}
\begin{split}
\phi_{x}^{t+1} = \displaystyle\sum_{i}\Big( (1-\omega + \omega w_{i}) \phi_{x-c_{i}}^{t}\Big) - (1-\omega) \displaystyle\sum_{i}\Bigg( \displaystyle\sum_{k\neq i} \left( \omega w_{k} \phi_{x-c_{i}-c_{k}}^{t-1} + (1-\omega) f_{k,x-c_{i}-c_{j}}^{t-1} \right) \Bigg),
\end{split}
\end{equation*}
\sphinxAtStartPar
notice that \(\displaystyle\sum_{i}\Bigg( \displaystyle\sum_{k\neq i} \left( \omega w_{k} \phi_{x-c_{i}-c_{k}}^{t-1} \right) \Bigg)=\displaystyle\sum_{i}\Bigg( \displaystyle\sum_{k\neq i} \left(w_{k}\right) \omega \phi_{x+c_{i}}^{t-1}  \Bigg)\) and using Eq. \eqref{equation:Refs/LBM-Study/FDM-Belloti/FDM-LBM-Diffusive:moments-fdm-fi-df} we have \(\displaystyle\sum_{i}\Bigg( \displaystyle\sum_{k\neq i} \left( f_{k,x-c_{i}-c_{j}}^{t-1} \right) \Bigg)= \displaystyle\sum_{i}\left(\phi_{x+c_{i}}^{t-1} - f_{i,x+c_{i}}^{t-1}\right) \). So rewriting the above equation base in the observations:
\begin{equation}\label{equation:Refs/LBM-Study/FDM-Belloti/FDM-LBM-Diffusive:last-fdm-form}
\begin{split}
\phi_{x}^{t+1} = \displaystyle\sum_{i}\Big( (1-\omega + \omega w_{i}) \phi_{x-c_{i}}^{t}\Big) - (1-\omega)\displaystyle\sum_{i}\Bigg( \left(1-\omega + \omega \displaystyle\sum_{k\neq i} \left(w_{k}\right) \right) \phi_{x+c_{i}}^{t-1} \Bigg) + (1-\omega)^{2}\displaystyle\sum_{i}\Bigg( f_{i,x+c_{i}}^{t-1} \Bigg).
\end{split}
\end{equation}
\sphinxAtStartPar
Interpreting \(f_{i,x+c_{i}}^{t-1}\) with Eq. \eqref{equation:Refs/LBM-Study/FDM-Belloti/FDM-LBM-Diffusive:pre-fdm-LB-df-Eq}:
\begin{equation*}
\begin{split}
f_{i,x+c_{i}}^{t-1}  = f_{i,x}^{t-2}  - \omega \left( f_{i,x}^{t-2} - f_{i,x}^{eq,t-2} \right) ,
\end{split}
\end{equation*}
\sphinxAtStartPar
and replaing it in Eq. \eqref{equation:Refs/LBM-Study/FDM-Belloti/FDM-LBM-Diffusive:last-fdm-form}, we obtained a four\sphinxhyphen{}level explicity scheme:
\begin{equation*}
\begin{split}
\phi_{x}^{t+1} = \displaystyle\sum_{i}\Big( (1-\omega + \omega w_{i}) \phi_{x-c_{i}}^{t}\Big) - (1-\omega)\displaystyle\sum_{i}\Bigg( \left(1-\omega + \omega \displaystyle\sum_{k\neq i} \left(w_{k}\right) \right) \phi_{x+c_{i}}^{t-1} \Bigg) + (1-\omega)^{2}\displaystyle\sum_{i}\Bigg(f_{i,x}^{t-2}  - \omega \left( f_{i,x}^{t-2} - f_{i,x}^{eq,t-2} \right)  \Bigg),
\end{split}
\end{equation*}\begin{equation}\label{equation:Refs/LBM-Study/FDM-Belloti/FDM-LBM-Diffusive:df-final-form-fdm}
\begin{split}
\phi_{x}^{t+1} = \displaystyle\sum_{i}\Big( (1-\omega + \omega w_{i}) \phi_{x-c_{i}}^{t}\Big) - (1-\omega)\displaystyle\sum_{i}\Bigg( \left(1-\omega + \omega \displaystyle\sum_{k\neq i} \left(w_{k}\right) \right) \phi_{x+c_{i}}^{t-1} \Bigg) + (1-\omega)^{2} \phi_{x}^{t-2}
\end{split}
\end{equation}
\sphinxAtStartPar
Expanding the sums:
\begin{equation*}
\begin{split}
\begin{array}{c}
\phi_{x}^{t+1} = (1-\omega + \omega w_{0}) \phi_{x}^{t} + (1-\omega + \omega w_{1}) \phi_{x-1}^{t} + (1-\omega + \omega w_{2}) \phi_{x+1}^{t} - (1-\omega)^{2} \Big( \phi_{x}^{t-1} + \phi_{x-1}^{t-1} + \phi_{x+1}^{t-1}  \Big) \\
-(1-\omega)\omega \Big( (w_{1}+w_{2})\phi_{x}^{t-1} + (w_{0}+w_{1})\phi_{x-1}^{t-1} + (w_{0}+w_{2})\phi_{x+1}^{t-1}  \Big) + (1-\omega)^{2} \phi_{x}^{t-2}
\end{array}
\end{split}
\end{equation*}

\section{1st Benchmark: Decay of Sinoidal Wave}
\label{\detokenize{Refs/LBM-Study/FDM-Belloti/FDM-LBM-Diffusive:st-benchmark-decay-of-sinoidal-wave}}
\sphinxAtStartPar
Considering a 1D \sphinxstylestrong{convection–diffusion equation} with constant coefficients
\$\(
\frac{\partial u}{\partial t}+c\frac{\partial u}{\partial x} = \nu\frac{\partial^2 u}{\partial x^2} \qquad x\in\mathbb{R},\ t\ge 0,
\)\$

\sphinxAtStartPar
with \sphinxstylestrong{sinusoidal initial condition}
\begin{equation*}
\begin{split}
u(x,0)=A\sin(kx+\phi),
\end{split}
\end{equation*}
\sphinxAtStartPar
where \(A\) is the initial amplitude, \(k\) the wavenumber, \(\phi\) a phase shift.

\sphinxAtStartPar
Then the analytical solution is
\$\(
\boxed{u(x,t)=Ae^{-\nu k^{2}t}\sin\big(k(x-ct)+\phi\big) }
\)\$

\sphinxAtStartPar
So:
\begin{itemize}
\item {} 
\sphinxAtStartPar
\sphinxstylestrong{Convection} shifts the wave: \(x \mapsto x-ct\) (wave travels at speed \(c\)).

\item {} 
\sphinxAtStartPar
\sphinxstylestrong{Diffusion} damps the amplitude exponentially: \(A(t)=Ae^{-\nu k^{2}t}\).

\end{itemize}

\sphinxAtStartPar
If you use a cosine initial condition \(u(x,0)=A\cos(kx+\phi)\), the solution is the same form:
\$\(
u(x,t)=Ae^{-\nu k^{2}t}\cos\big(k(x-ct)+\phi\big).
\)\$


\subsection{LBM D1Q3 Code \sphinxhyphen{} FDM Code (Only Diffusive Case \protect\(c=0\protect\))}
\label{\detokenize{Refs/LBM-Study/FDM-Belloti/FDM-LBM-Diffusive:lbm-d1q3-code-fdm-code-only-diffusive-case-c-0}}
\begin{sphinxuseclass}{cell}
\begin{sphinxuseclass}{tag_hide-input}
\begin{sphinxuseclass}{tag_hide-output}
\end{sphinxuseclass}
\end{sphinxuseclass}
\end{sphinxuseclass}
\begin{sphinxuseclass}{cell}
\begin{sphinxuseclass}{tag_hide-input}\begin{sphinxVerbatimOutput}

\begin{sphinxuseclass}{cell_output}
\noindent\sphinxincludegraphics{{69053fd6f9da8c2d1c7ebadd2831b9a51623ff3f1f4afc09d37df03bada78940}.png}

\end{sphinxuseclass}\end{sphinxVerbatimOutput}

\end{sphinxuseclass}
\end{sphinxuseclass}

\section{2nd Benchmark: Temperature ramp}
\label{\detokenize{Refs/LBM-Study/FDM-Belloti/FDM-LBM-Diffusive:nd-benchmark-temperature-ramp}}
\sphinxAtStartPar
Analyzing the accuracy of diffusive equation solver, we apply it to a one\sphinxhyphen{}dimensional analytical case. The geometry is described by a domain of length \(L\) initialized with a constant temperatura \(T(x,t=0)=0\). The boundary conditions are given \(T(x=0,t)=1\) and \(T(x=L,t)=0\).


\subsection{Steady\sphinxhyphen{}State Analytical Solution}
\label{\detokenize{Refs/LBM-Study/FDM-Belloti/FDM-LBM-Diffusive:steady-state-analytical-solution}}
\sphinxAtStartPar
The analytical solution of this problem is given by:
\begin{equation*}
\begin{split}
T(x,t\rightarrow \infty)= -\frac{x}{L} + 1.
\end{split}
\end{equation*}

\subsection{Transient Analytical Solution}
\label{\detokenize{Refs/LBM-Study/FDM-Belloti/FDM-LBM-Diffusive:transient-analytical-solution}}
\sphinxAtStartPar
For the 1D heat equation \(\partial_t T=\nu\,\partial_{xx}T\) on \(0<x<L\) with
\begin{equation*}
\begin{split}
T(0,t)=1,\qquad T(L,t)=0,
\end{split}
\end{equation*}
\sphinxAtStartPar
the steady state is \(T_{\text{ss}}(x)=1-\tfrac{x}{L}\). With \(u=T-T_{\text{ss}}\) (so \(u(0,t)=u(L,t)=0\)), the solution is
\begin{equation*}
\begin{split}
T(x,t)=1-\frac{x}{L}+\sum_{n=1}^{\infty} B_n\,e^{-\nu (n\pi/L)^2 t}\,\sin\!\Big(\frac{n\pi x}{L}\Big),
\quad
B_n=\frac{2}{L}\int_0^L\!\Big[T_0(\xi)-\Big(1-\frac{\xi}{L}\Big)\Big]\sin\!\Big(\frac{n\pi \xi}{L}\Big)\,d\xi.
\end{split}
\end{equation*}
\sphinxAtStartPar
Special case \(T_0(x)=0\): \(B_n=-\dfrac{2}{n\pi}\), so
\begin{equation*}
\begin{split}
T(x,t)=1-\frac{x}{L}-\frac{2}{\pi}\sum_{n=1}^{\infty}\frac{1}{n}\,e^{-\nu (n\pi/L)^2 t}\,\sin\!\Big(\frac{n\pi x}{L}\Big).
\end{split}
\end{equation*}
\sphinxAtStartPar
Here’s a ready\sphinxhyphen{}to\sphinxhyphen{}run Python snippet (with a single plot of \(T(x,t)\) vs \(x\) for multiple times):


\subsection{LBM D1Q3 Code \sphinxhyphen{} FDM Code}
\label{\detokenize{Refs/LBM-Study/FDM-Belloti/FDM-LBM-Diffusive:lbm-d1q3-code-fdm-code}}
\begin{sphinxuseclass}{cell}
\begin{sphinxuseclass}{tag_hide-input}
\begin{sphinxuseclass}{tag_hide-output}
\end{sphinxuseclass}
\end{sphinxuseclass}
\end{sphinxuseclass}
\begin{sphinxuseclass}{cell}
\begin{sphinxuseclass}{tag_hide-input}\begin{sphinxVerbatimOutput}

\begin{sphinxuseclass}{cell_output}
\noindent\sphinxincludegraphics{{39d476457e161e0f1e8598e65ae1e5d039de6d3b85f309ef1e58a0b6d3064822}.png}

\end{sphinxuseclass}\end{sphinxVerbatimOutput}

\end{sphinxuseclass}
\end{sphinxuseclass}

\section{Convergence Errors: Fourth\sphinxhyphen{}order Diffusive Equation}
\label{\detokenize{Refs/LBM-Study/FDM-Belloti/FDM-LBM-Diffusive:convergence-errors-fourth-order-diffusive-equation}}
\sphinxAtStartPar
The finite\sphinxhyphen{}difference interpretation of the lattice Boltzmann equation allows for a precise analysis of the numerical errors associated with the spatial and temporal discretization of the method. Following the procedure presented below, these numerical errors can be identified at different orders of \(\Delta x\) and \(\Delta t\) and, when possible, eliminated through an appropriate choice of model parameters.

\sphinxAtStartPar
Expanding the sums:
\begin{equation*}
\begin{split}
\begin{array}{c}
\phi_{x}^{t+1} = (1-\omega + \omega w_{0}) \phi_{x}^{t} + (1-\omega + \omega w_{1}) \phi_{x-1}^{t} + (1-\omega + \omega w_{2}) \phi_{x+1}^{t} - (1-\omega)^{2} \Big( \phi_{x}^{t-1} + \phi_{x-1}^{t-1} + \phi_{x+1}^{t-1}  \Big) \\
-(1-\omega)\omega \Big( (w_{1}+w_{2})\phi_{x}^{t-1} + (w_{0}+w_{1})\phi_{x-1}^{t-1} + (w_{0}+w_{2})\phi_{x+1}^{t-1}  \Big) + (1-\omega)^{2} \phi_{x}^{t-2}
\end{array}
\end{split}
\end{equation*}
\sphinxAtStartPar
Considering
\$\(
\alpha_{1}=\Omega + a\omega, \quad \alpha_{2}=\Omega + (1-2a)\omega, \quad \beta_{1}=-\Omega(\Omega + (1-a)\omega), \quad \textrm{and} \quad \beta_{2}=-\Omega(\Omega + 2a\omega),
\)\$

\sphinxAtStartPar
where \(\Omega=(1-\omega)\) and \(a=w_{1}=w_{2}\). Replacing in the Eq. \eqref{equation:Refs/LBM-Study/FDM-Belloti/FDM-LBM-Diffusive:df-final-form-fdm}:
\begin{equation*}
\begin{split}
\begin{array}{c}
\phi_{x}^{t+1} = \alpha_{1} \phi_{x-1}^{t} + \alpha_{2} \phi_{x}^{t} + \alpha_{1} \phi_{x+1}^{t} + \beta_{1}\phi_{x-1}^{t-1} + \beta_{2}\phi_{x}^{t-1} + \beta_{1}\phi_{x+1}^{t-1}  + (1-\omega)^{2} \phi_{x}^{t-2}
\end{array}
\end{split}
\end{equation*}
\sphinxAtStartPar
To recover the partial differetial equation described by the above equation, we expand each \(\phi\) term in Taylor series:
\begin{equation*}
\begin{split}
\phi_{x}^{t\pm 1} = \phi \pm \Delta_{t}\frac{\partial\phi}{\partial t} + \frac{\Delta_{t}^{2}}{2}\frac{\partial^{2}\phi}{\partial t^{2}} \pm \frac{\Delta_{t}^{3}}{6}\frac{\partial^{3}\phi}{\partial t^{3}} + \frac{\Delta_{t}^{4}}{24}\frac{\partial^{4}\phi}{\partial t^{4}} \pm  \mathcal{O}(\Delta t)^{4},
\end{split}
\end{equation*}\begin{equation*}
\begin{split}
\phi_{x\pm 1}^{t} = \phi \pm \Delta_{x}\frac{\partial\phi}{\partial x} + \frac{\Delta_{x}^{2}}{2}\frac{\partial^{2}\phi}{\partial x^{2}} \pm \frac{\Delta_{x}^{3}}{6}\frac{\partial^{3}\phi}{\partial x^{3}} + \frac{\Delta_{x}^{4}}{24}\frac{\partial^{4}\phi}{\partial x^{4}} \pm \mathcal{O}(\Delta x)^{4},
\end{split}
\end{equation*}\begin{equation*}
\begin{split}
\phi_{x}^{t - 2} = \phi - 2\Delta_{t}\frac{\partial\phi}{\partial t} + \frac{4\Delta_{t}^{2}}{2}\frac{\partial^{2}\phi}{\partial t^{2}} - \frac{8\Delta_{t}^{3}}{6}\frac{\partial^{3}\phi}{\partial t^{3}} + \frac{16\Delta_{t}^{4}}{24}\frac{\partial^{4}\phi}{\partial t^{4}} \pm  \mathcal{O}(\Delta t)^{4},
\end{split}
\end{equation*}\begin{equation*}
\begin{split}
\phi_{x\pm 1}^{t-1} = \phi \pm \Delta_{x}\frac{\partial\phi}{\partial x} - \Delta_{t}\frac{\partial\phi}{\partial t} + \frac{\Delta_{x}^{2}}{2}\frac{\partial^{2}\phi}{\partial x^{2}} \mp \Delta_{t}\Delta_{x} \frac{\partial^{2}\phi}{\partial x\partial t} + \frac{\Delta_{t}^{2}}{2}\frac{\partial^{2}\phi}{\partial t^{2}} \pm \frac{\Delta_{x}^{3}}{6}\frac{\partial^{3}\phi}{\partial x^{3}} \mp \frac{\Delta_{x}^{2}\Delta_{t}}{2}\frac{\partial^{3}\phi}{\partial x^{2}\partial t}  \pm \frac{\Delta_{x}\Delta_{t}^{2}}{2}\frac{\partial^{3}\phi}{\partial x\partial t^{2}} - \frac{\Delta_{t}^{3}}{6}\frac{\partial^{3}\phi}{\partial t^{3}} + \mathcal{O}(\Delta x \Delta t)^{3}.
\end{split}
\end{equation*}
\sphinxAtStartPar
Replacing the expanded terms and simplifying the term (this link provide a sympy code that reach the above equation \DUrole{xref,std,std-doc}{appendix/sympy\sphinxhyphen{}code\sphinxhyphen{}fdm\sphinxhyphen{}diffusive\sphinxhyphen{}expansion}), we have:
\begin{equation*}
\begin{split}
0=\Delta_{t} \omega^{2} \frac{\partial\phi}{\partial t} + \Delta_{x}^{2} a \omega \left(\omega - 2\right) \frac{\partial^{2}\phi}{\partial x^{2}}  + \frac{\Delta_{t}^{2} \omega \left(4 - 3 \omega\right) }{2} \frac{\partial^{2}\phi}{\partial t^{2}} - \Delta_{t} \Delta_{x}^{2} \left(a \omega^{2} - a \omega - \omega + 1\right) \frac{\partial^{3}\phi}{\partial x^{2}\partial t} + \frac{\Delta_{x}^{4} a \omega \left(\omega - 2\right) }{12}\frac{\partial^{4}\phi}{\partial x^{4}}
\end{split}
\end{equation*}\begin{equation*}
\begin{split}
\Delta_{t} \omega^{2} \frac{\partial\phi}{\partial t} = - \Delta_{x}^{2} a \omega \left(\omega - 2\right) \frac{\partial^{2}\phi}{\partial x^{2}}  - \frac{\Delta_{t}^{2} \omega \left(4 - 3 \omega\right) }{2} \frac{\partial^{2}\phi}{\partial t^{2}} + \Delta_{t} \Delta_{x}^{2} \left(a \omega^{2} - a \omega - \omega + 1\right) \frac{\partial^{3}\phi}{\partial x^{2}\partial t} - \frac{\Delta_{x}^{4} a \omega \left(\omega - 2\right) }{12}\frac{\partial^{4}\phi}{\partial x^{4}}
\end{split}
\end{equation*}\begin{equation*}
\begin{split}
\frac{\partial\phi}{\partial t} = \frac{a \left(2 - \omega \right)}{\omega} \frac{\Delta_{x}^{2}}{\Delta_{t}} \frac{\partial^{2}\phi}{\partial x^{2}}  + \frac{ \left(3 \omega - 4\right) \Delta_{t}}{2\omega} \frac{\partial^{2}\phi}{\partial t^{2}} + \frac{\left(a \omega^{2} - a \omega - \omega + 1\right)\Delta_{x}^{2}}{\omega^{2}} \frac{\partial^{3}\phi}{\partial x^{2}\partial t} + \frac{ a \left(2 - \omega\right) }{12\omega} \frac{\Delta_{x}^{4}}{\Delta_{t}} \frac{\partial^{4}\phi}{\partial x^{4}}
\end{split}
\end{equation*}
\sphinxAtStartPar
where \(\nu=\frac{a \left(2 - \omega \right)}{\omega} \frac{\Delta_{x}^{2}}{\Delta_{t}}\) is the same term recovered in the Chapman\sphinxhyphen{}Enskog Analysis. To enforce fourth\sphinxhyphen{}order \(\Delta_{x}\) convergence rate, employ the relations
\begin{equation*}
\begin{split}
\frac{\partial\phi}{\partial t} = \nu \frac{\partial^{2}\phi}{\partial x^{2}}, \qquad \frac{\partial^{2}\phi}{\partial t^{2}} = \nu^{2} \frac{\partial^{4}\phi}{\partial x^{4}}, \qquad \frac{\partial^{3}\phi}{\partial x^{2}\partial t} = \nu \frac{\partial^{4}\phi}{\partial x^{4}}
\end{split}
\end{equation*}\begin{equation*}
\begin{split}
\frac{\partial\phi}{\partial t} = \nu\frac{\partial^{2}\phi}{\partial x^{2}}  + \frac{a \left(2 - \omega \right)}{\omega} \frac{ \left(3 \omega - 4\right) \nu\Delta_{x}^{2}}{2\omega} \frac{\partial^{4}\phi}{\partial x^{4}} + \frac{\left(a \omega^{2} - a \omega - \omega + 1\right)\nu\Delta_{x}^{2}}{\omega^{2}} \frac{\partial^{4}\phi}{\partial x^{4}} + \frac{\nu\Delta_{x}^{2}}{12} \frac{\partial^{4}\phi}{\partial x^{4}}
\end{split}
\end{equation*}\begin{equation*}
\begin{split}
\frac{\partial\phi}{\partial t} = \nu\frac{\partial^{2}\phi}{\partial x^{2}}  + \frac{ a\left(6 \omega - 8 - 3\omega^{2} + 4 \omega \right) \nu\Delta_{x}^{2}}{2\omega^{2}} \frac{\partial^{4}\phi}{\partial x^{4}} + \frac{\left(a \omega^{2} - a \omega - \omega + 1\right)\nu\Delta_{x}^{2}}{\omega^{2}} \frac{\partial^{4}\phi}{\partial x^{4}} + \frac{\omega^{2}\nu\Delta_{x}^{2}}{12\omega^{2}} \frac{\partial^{4}\phi}{\partial x^{4}}
\end{split}
\end{equation*}\begin{equation*}
\begin{split}
\frac{\partial\phi}{\partial t} = \nu\frac{\partial^{2}\phi}{\partial x^{2}}  + \frac{\nu\Delta_{x}^{2}}{2\omega^{2}}\left( a\left(10 \omega - 8 - 3\omega^{2} \right)  + 2\left(a \omega^{2} - a \omega - \omega + 1\right) + \frac{\omega^{2}}{6} \right) \frac{\partial^{4}\phi}{\partial x^{4}} 
\end{split}
\end{equation*}\begin{equation*}
\begin{split}
\frac{\partial\phi}{\partial t} = \nu\frac{\partial^{2}\phi}{\partial x^{2}}  + \frac{\nu\Delta_{x}^{2}}{2\omega^{2}}\left( a\left(8 \omega - 8 - \omega^{2} \right)  + 2\left(1 - \omega \right) + \frac{\omega^{2}}{6} \right) \frac{\partial^{4}\phi}{\partial x^{4}} 
\end{split}
\end{equation*}\begin{equation*}
\begin{split}
\frac{\partial\phi}{\partial t} = \nu\frac{\partial^{2}\phi}{\partial x^{2}}  - \frac{\nu\Delta_{x}^{2}}{2\omega^{2}}\left( a\left(\omega^{2} - 8 \omega + 8 \right)  - \frac{\omega^{2} - 12\omega + 12 }{6} \right) \frac{\partial^{4}\phi}{\partial x^{4}} + \mathcal{O}(\Delta x)^{4}
\end{split}
\end{equation*}
\begin{sphinxuseclass}{cell}
\begin{sphinxuseclass}{tag_hide-input}\begin{sphinxVerbatimOutput}

\begin{sphinxuseclass}{cell_output}
\noindent\sphinxincludegraphics{{3a440122c391caca3638e8d174f3feca89bc6a1f796dd3a778e5d49e763311a3}.png}

\end{sphinxuseclass}\end{sphinxVerbatimOutput}

\end{sphinxuseclass}
\end{sphinxuseclass}
\begin{sphinxuseclass}{cell}
\begin{sphinxuseclass}{tag_hide-input}
\begin{sphinxuseclass}{tag_hide-output}
\end{sphinxuseclass}
\end{sphinxuseclass}
\end{sphinxuseclass}
\begin{sphinxuseclass}{cell}
\begin{sphinxuseclass}{tag_hide-input}\begin{sphinxVerbatimOutput}

\begin{sphinxuseclass}{cell_output}
\noindent\sphinxincludegraphics{{dd87cc115dae78bb74ff29c3aa999dd5d8cdaea33236e522a771dbcd27a44687}.png}

\end{sphinxuseclass}\end{sphinxVerbatimOutput}

\end{sphinxuseclass}
\end{sphinxuseclass}
\begin{sphinxuseclass}{cell}
\begin{sphinxuseclass}{tag_hide-input}
\begin{sphinxuseclass}{tag_hide-output}
\end{sphinxuseclass}
\end{sphinxuseclass}
\end{sphinxuseclass}
\begin{sphinxuseclass}{cell}
\begin{sphinxuseclass}{tag_hide-input}
\end{sphinxuseclass}
\end{sphinxuseclass}
\begin{sphinxuseclass}{cell}
\begin{sphinxuseclass}{tag_hide-input}\begin{sphinxVerbatimOutput}

\begin{sphinxuseclass}{cell_output}
\noindent\sphinxincludegraphics{{1f2930681a48377c82844249177836919539ffa5f4c574a0356a3d73d2f8e37d}.png}

\end{sphinxuseclass}\end{sphinxVerbatimOutput}

\end{sphinxuseclass}
\end{sphinxuseclass}

\section{Extension to Lattices with \protect\(d\protect\)\sphinxhyphen{}Dimensions and \protect\(q\protect\)\sphinxhyphen{}Velocities (\protect\(DdQq\protect\)):}
\label{\detokenize{Refs/LBM-Study/FDM-Belloti/FDM-LBM-Diffusive:extension-to-lattices-with-d-dimensions-and-q-velocities-ddqq}}
\sphinxAtStartPar
Through a general mathematical manipulation of the single relaxation time lattice Boltzmann equation, it is possible to derive a general expansion that expresses the lattice framework in a finite\sphinxhyphen{}difference form.

\sphinxAtStartPar
Building on the earlier analysis presented in Suga {[}\hyperlink{cite.Refs/Citations:id4}{Sug10}{]}, Suga {[}\hyperlink{cite.Refs/Citations:id9}{Sug14}{]} proposed a generalized formulation for diffusive transport on \(DdQq\) lattices, considering both isotropic and anisotropic diffusion. The resulting equations are presented in Eq. \eqref{equation:Refs/LBM-Study/FDM-Belloti/FDM-LBM-Diffusive:Suga-2014-Eqs}, with truncation\sphinxhyphen{}error terms highlighted in red and retained through fourth order in \(\Delta x\). Notice that additional second\sphinxhyphen{}order errors may arise when extending the formulation to higher dimensions:
\begin{equation}\label{equation:Refs/LBM-Study/FDM-Belloti/FDM-LBM-Diffusive:Suga-2014-Eqs}
\begin{split}
\begin{array}{ll}
D1Q3-Isotropic: \qquad & \displaystyle\frac{\partial\phi}{\partial t} = \nu\frac{\partial^{2}\phi}{\partial x^{2}}  - \textcolor{red} {\frac{\nu\Delta_{x}^{2}}{2\omega^{2}}\left( a\left(\omega^{2} - 8 \omega + 8 \right)  - \frac{\omega^{2} - 12\omega + 12 }{6} \right) \frac{\partial^{4}\phi}{\partial x^{4}} } + \mathcal{O}(\Delta x)^{4}\\
D2Q9-Isotropic: \qquad & \displaystyle\frac{\partial\phi}{\partial t} = \nu\left(\frac{\partial^{2}\phi}{\partial x^{2}} + \frac{\partial^{2}\phi}{\partial y^{2}}\right)  - \textcolor{red} { \Delta_{x}^{2} R1 \left( \frac{\partial^{4}\phi}{\partial x^{4}} + \frac{\partial^{4}\phi}{\partial y^{4}} \right) } - \textcolor{red} { \Delta_{x}^{2} R12 \left( \frac{\partial^{4}\phi}{\partial x^{2}\partial y^{2}} \right) } + \mathcal{O}(\Delta x)^{4}\\
D2Q13-Anisotropic: \qquad & \displaystyle\frac{\partial\phi}{\partial t} = \nu_{x}\frac{\partial^{2}\phi}{\partial x^{2}} + \nu_{y}\frac{\partial^{2}\phi}{\partial y^{2}}  - \textcolor{red} { \Delta_{x}^{2} R1 \frac{\partial^{4}\phi}{\partial x^{4}} + \Delta_{x}^{2} R2\frac{\partial^{4}\phi}{\partial y^{4}}  } - \textcolor{red} { \Delta_{x}^{2} R12 \left( \frac{\partial^{4}\phi}{\partial x^{2}\partial y^{2}} \right) } + \mathcal{O}(\Delta x)^{4}\\
D3Q19-Isotropic: \qquad & \displaystyle\frac{\partial\phi}{\partial t} = \nu\left(\frac{\partial^{2}\phi}{\partial x^{2}} + \frac{\partial^{2}\phi}{\partial y^{2}} + \frac{\partial^{2}\phi}{\partial z^{2}}\right)  - \textcolor{red} { \Delta_{x}^{2} R1 \left( \frac{\partial^{4}\phi}{\partial x^{4}} + \frac{\partial^{4}\phi}{\partial y^{4}} \right) } - \textcolor{red} { \Delta_{x}^{2} R12 \left( \frac{\partial^{4}\phi}{\partial x^{2}\partial y^{2}} + \frac{\partial^{4}\phi}{\partial x^{2}\partial z^{2}} + \frac{\partial^{4}\phi}{\partial y^{2}\partial z^{2}} \right) } + \mathcal{O}(\Delta x)^{4}\\
D3Q25-Anisotropic: \qquad & \displaystyle\frac{\partial\phi}{\partial t} = \nu_{x}\frac{\partial^{2}\phi}{\partial x^{2}} + \nu_{y}\frac{\partial^{2}\phi}{\partial y^{2}} + \nu_{z}\frac{\partial^{2}\phi}{\partial z^{2}}  - \textcolor{red} { \Delta_{x}^{2} \left( R1\frac{\partial^{4}\phi}{\partial x^{4}} + R2\frac{\partial^{4}\phi}{\partial y^{4}} + R3\frac{\partial^{4}\phi}{\partial y^{4}} \right) } - \textcolor{red} { \Delta_{x}^{2} \left( R12\frac{\partial^{4}\phi}{\partial x^{2}\partial y^{2}} + R13\frac{\partial^{4}\phi}{\partial x^{2}\partial z^{2}} + R23\frac{\partial^{4}\phi}{\partial y^{2}\partial z^{2}} \right) } + \mathcal{O}(\Delta x)^{4}
\end{array}
\end{split}
\end{equation}
\sphinxAtStartPar
In the Section XX we will present a generalized expansion of LB equation with single relaxation time in FDM framewrok for any equilibrium distribution function with \(n\) conservative terms in \(DdQq\) lattice structures.


\section{Diffusive Equation Interpretation with Multi\sphinxhyphen{}Relaxation Time (MRT\sphinxhyphen{}Framework)}
\label{\detokenize{Refs/LBM-Study/FDM-Belloti/FDM-LBM-Diffusive:diffusive-equation-interpretation-with-multi-relaxation-time-mrt-framework}}
\sphinxAtStartPar
The multiple\sphinxhyphen{}relaxation\sphinxhyphen{}time (MRT) collision model is a well\sphinxhyphen{}established scheme in the LBM literature that allows different non\sphinxhyphen{}equilibrium moments to relax independently. This flexibility enables various improvements in the numerical representation of partial differential equations. In the present section, we extend the FDM interpretation to the MRT framework of diffusive equation, following the analysis presented by Lin \sphinxstyleemphasis{et al.} {[}\hyperlink{cite.Refs/Citations:id10}{LHSC21}{]}.

\sphinxAtStartPar
Given the LB equation with MRT for the diffusive equation:
\begin{equation}\label{equation:Refs/LBM-Study/FDM-Belloti/FDM-LBM-Diffusive:MRT-LB-fdm-diff-Eq}
\begin{split}
f_i( x_{\alpha} + e_{i,\alpha} \delta t, t+\delta t) = f_i(x_{\alpha}, t)  - (\textbf{M}^{-1}\textbf{S}\textbf{M})_{ik}\left( f_{k} - f_{k}^{eq}  \right), 
\end{split}
\end{equation}
\sphinxAtStartPar
where \(\textbf{M}\) is matrix of moments, \(\textbf{S}\) is the diagonal matrix of relaxation times, and \(\textbf{M}^{-1}\) is the inverse matrix of \(\textbf{M}\). Considering the problem description in a lattice D1Q3, we have following a Gram\sphinxhyphen{}Schmidt procedure:
\begin{equation*}
\begin{split}
\textbf{M} =\begin{pmatrix} 1 \\ e_{i,x} \\ 3e_{i,x}^{2}-2 \end{pmatrix} =\begin{pmatrix} 1 & 1 & 1 \\ 0 & 1   & -1 \\ -2  & 1 & 1 \end{pmatrix},  \qquad \qquad  \textbf{S} = \begin{pmatrix} 0 & 0   & 0 \\ 0   & s_1 & 0 \\ 0   & 0   & s_2 \end{pmatrix},  \qquad \qquad \textbf{M}^{-1}=\begin{pmatrix}\dfrac{1}{3} & 0 & -\dfrac{1}{3} \\ \dfrac{1}{3} & \dfrac{1}{2} & \dfrac{1}{6} \\ \dfrac{1}{3} & -\dfrac{1}{2} & \dfrac{1}{6} \end{pmatrix}.
\end{split}
\end{equation*}
\begin{sphinxuseclass}{cell}
\begin{sphinxuseclass}{tag_remove-input}
\end{sphinxuseclass}
\end{sphinxuseclass}
\sphinxAtStartPar
Applying the same shift scheme employed in Eq. \eqref{equation:Refs/LBM-Study/FDM-Belloti/FDM-LBM-Diffusive:pre-fdm-LB-df-Eq} to Eq. \eqref{equation:Refs/LBM-Study/FDM-Belloti/FDM-LBM-Diffusive:MRT-LB-fdm-diff-Eq}, we obtain

\begin{sphinxuseclass}{cell}
\begin{sphinxuseclass}{tag_remove-input}\begin{sphinxVerbatimOutput}

\begin{sphinxuseclass}{cell_output}\begin{equation*}
\begin{split}\displaystyle \left[\begin{matrix}f_{0,x}^{t+1}\\f_{1,x}^{t+1}\\f_{2,x}^{t+1}\end{matrix}\right] = \left[\begin{matrix}\phi^{t}_{x} s_{2} w_{0} - \frac{\phi^{t}_{x} s_{2}}{3} - \frac{2 f_{0,x}^{t} s_{2}}{3} + f_{0,x}^{t} + \frac{f_{1,x}^{t} s_{2}}{3} + \frac{f_{2,x}^{t} s_{2}}{3}\\- \frac{\phi_{x-1}^t s_{2} w_{0}}{2} + \frac{\phi_{x-1}^t s_{2}}{6} + \frac{f_{0,x-1}^{t} s_{2}}{3} - \frac{f_{1,x-1}^{t} s_{1}}{2} - \frac{f_{1,x-1}^{t} s_{2}}{6} + f_{1,x-1}^{t} + \frac{f_{2,x-1}^{t} s_{1}}{2} - \frac{f_{2,x-1}^{t} s_{2}}{6}\\- \frac{\phi_{x+1}^t s_{2} w_{0}}{2} + \frac{\phi_{x+1}^t s_{2}}{6} + \frac{f_{0,x+1}^{t} s_{2}}{3} + \frac{f_{1,x+1}^{t} s_{1}}{2} - \frac{f_{1,x+1}^{t} s_{2}}{6} - \frac{f_{2,x+1}^{t} s_{1}}{2} - \frac{f_{2,x+1}^{t} s_{2}}{6} + f_{2,x+1}^{t}\end{matrix}\right],\end{split}
\end{equation*}
\end{sphinxuseclass}\end{sphinxVerbatimOutput}

\end{sphinxuseclass}
\end{sphinxuseclass}
\sphinxAtStartPar
expanding each term and replacing the conservative terms:

\begin{sphinxuseclass}{cell}
\begin{sphinxuseclass}{tag_remove-input}\begin{sphinxVerbatimOutput}

\begin{sphinxuseclass}{cell_output}\begin{equation*}
\begin{split}\displaystyle f_{0,x}^{t+1} = \phi^{t}_{x} s_{2} w_{0} + f_{0,x}^{t} \left(1 - s_{2}\right),\end{split}
\end{equation*}
\end{sphinxuseclass}\end{sphinxVerbatimOutput}

\end{sphinxuseclass}
\end{sphinxuseclass}
\begin{sphinxuseclass}{cell}
\begin{sphinxuseclass}{tag_remove-input}\begin{sphinxVerbatimOutput}

\begin{sphinxuseclass}{cell_output}\begin{equation*}
\begin{split}\displaystyle f_{1,x}^{t+1} = - \frac{\phi_{x-1}^t s_{2} w_{0}}{2} + \frac{f_{0,x-1}^{t} s_{2}}{2} + f_{1,x-1}^{t} + \frac{s_{1} \left(- f_{1,x-1}^{t} + f_{2,x-1}^{t}\right)}{2},\end{split}
\end{equation*}
\end{sphinxuseclass}\end{sphinxVerbatimOutput}

\end{sphinxuseclass}
\end{sphinxuseclass}
\begin{sphinxuseclass}{cell}
\begin{sphinxuseclass}{tag_remove-input}\begin{sphinxVerbatimOutput}

\begin{sphinxuseclass}{cell_output}\begin{equation*}
\begin{split}\displaystyle f_{2,x}^{t+1} = - \frac{\phi_{x+1}^t s_{2} w_{0}}{2} + \frac{f_{0,x+1}^{t} s_{2}}{2} + f_{2,x+1}^{t} + \frac{s_{1} \left(f_{1,x+1}^{t} - f_{2,x+1}^{t}\right)}{2}.\end{split}
\end{equation*}
\end{sphinxuseclass}\end{sphinxVerbatimOutput}

\end{sphinxuseclass}
\end{sphinxuseclass}
\sphinxAtStartPar
Summing over all lattice direction, we have:

\begin{sphinxuseclass}{cell}
\begin{sphinxuseclass}{tag_remove-input}\begin{sphinxVerbatimOutput}

\begin{sphinxuseclass}{cell_output}\begin{equation*}
\begin{split}\displaystyle \phi_{x+1}^t = f_{0,x}^{t} \left(1 - s_{2}\right) + f_{1,x-1}^{t} + f_{2,x+1}^{t} + \frac{s_{1} \left(f_{1,x+1}^{t} - f_{1,x-1}^{t} - f_{2,x+1}^{t} + f_{2,x-1}^{t}\right)}{2} + s_{2} w_{0} \left(\phi^{t}_{x} - \frac{\phi_{x+1}^t}{2} - \frac{\phi_{x-1}^t}{2}\right) + s_{2} \left(\frac{f_{0,x+1}^{t}}{2} + \frac{f_{0,x-1}^{t}}{2}\right),\end{split}
\end{equation*}
\end{sphinxuseclass}\end{sphinxVerbatimOutput}

\end{sphinxuseclass}
\end{sphinxuseclass}
\sphinxAtStartPar
Using the shift interpretation for the distribution functions, we obtain
\begin{equation}\label{equation:Refs/LBM-Study/FDM-Belloti/FDM-LBM-Diffusive:rel-f-1}
\begin{split}
f_{0,x}^{t}+f_{1,x+1}^{t}+f_{2,x-1}^{t}=\phi^{t-1}_{x},
\end{split}
\end{equation}
\sphinxAtStartPar
applying this relation to the equation above:

\begin{sphinxuseclass}{cell}
\begin{sphinxuseclass}{tag_remove-input}\begin{sphinxVerbatimOutput}

\begin{sphinxuseclass}{cell_output}\begin{equation*}
\begin{split}\displaystyle \phi_{x+1}^t = f_{0,x}^{t} \left(1 - s_{2}\right) + f_{1,x-1}^{t} + f_{2,x+1}^{t} + \frac{s_{1} \left(\phi_{x}^{t-1} - f_{0,x}^{t} - f_{1,x-1}^{t} - f_{2,x+1}^{t}\right)}{2} + s_{2} w_{0} \left(\phi^{t}_{x} - \frac{\phi_{x+1}^t}{2} - \frac{\phi_{x-1}^t}{2}\right) + s_{2} \left(\frac{f_{0,x+1}^{t}}{2} + \frac{f_{0,x-1}^{t}}{2}\right).\end{split}
\end{equation*}
\end{sphinxuseclass}\end{sphinxVerbatimOutput}

\end{sphinxuseclass}
\end{sphinxuseclass}
\sphinxAtStartPar
To repalce the term \(f_{0,x}^{t}+f_{1,x-1}^{t}+f_{2,x+1}^{t}\), in the above equation we use the \(\sum_{i}f_{i,x}^{t}=\phi_{x}^{t}\):

\begin{sphinxuseclass}{cell}
\begin{sphinxuseclass}{tag_remove-input}\begin{sphinxVerbatimOutput}

\begin{sphinxuseclass}{cell_output}\begin{equation*}
\begin{split}\displaystyle f_{0,x}^{t} + f_{1,x-1}^{t} + f_{2,x+1}^{t} = \phi_{x+1}^t + \phi_{x-1}^t - \phi_{x}^{t-1} - f_{0,x+1}^{t} - f_{0,x-1}^{t} + 2 f_{0,x}^{t},\end{split}
\end{equation*}
\end{sphinxuseclass}\end{sphinxVerbatimOutput}

\end{sphinxuseclass}
\end{sphinxuseclass}
\sphinxAtStartPar
where
\begin{equation}\label{equation:Refs/LBM-Study/FDM-Belloti/FDM-LBM-Diffusive:rel-pre-shift-2}
\begin{split}
\phi_{x+1}^t = \phi^{t}_{x} s_{2} w_{0} + \phi_{x+1}^t \left(- \frac{s_{1}}{2} - \frac{s_{2} w_{0}}{2} + 1\right) + \phi_{x-1}^t \left(- \frac{s_{1}}{2} - \frac{s_{2} w_{0}}{2} + 1\right) + \phi_{x}^{t-1} \left(s_{1} - 1\right) + \left(f_{0,x+1}^{t} + f_{0,x-1}^{t} - 2 f_{0,x}^{t}\right) \left(\frac{s_{1}}{2} + \frac{s_{2}}{2} - 1\right)
\end{split}
\end{equation}
\sphinxAtStartPar
To replace the distribution function in the above relation, we manipulate:

\begin{sphinxuseclass}{cell}
\begin{sphinxuseclass}{tag_remove-input}\begin{sphinxVerbatimOutput}

\begin{sphinxuseclass}{cell_output}\begin{equation*}
\begin{split}\displaystyle f_{0,x+1}^{t} + f_{0,x-1}^{t} - 2 f_{0,x}^{t} = - 2 \phi^{t}_{x} + f_{0,x+1}^{t} + f_{0,x-1}^{t} + 2 f_{1,x}^{t} + 2 f_{2,x}^{t},\end{split}
\end{equation*}\begin{equation*}
\begin{split}\displaystyle f_{0,x+1}^{t} + f_{0,x-1}^{t} - 2 f_{0,x}^{t} = - 2 \phi^{t}_{x} + f_{0,x+1}^{t-1} + f_{0,x-1}^{t-1} + 2 f_{1,x-1}^{t-1} + 2 f_{2,x+1}^{t-1} + s_{1} \left(f_{1,x+1}^{t-1} - f_{1,x-1}^{t-1} - f_{2,x+1}^{t-1} + f_{2,x-1}^{t-1}\right),\end{split}
\end{equation*}\begin{equation*}
\begin{split}\displaystyle f_{0,x+1}^{t} + f_{0,x-1}^{t} - 2 f_{0,x}^{t} = - 2 \phi^{t}_{x} + 2 \phi_{x+1}^{t-1} + 2 \phi_{x-1}^{t-1} - f_{0,x+1}^{t-1} - f_{0,x-1}^{t-1} - 2 f_{1,x+1}^{t-1} - 2 f_{2,x-1}^{t-1} + s_{1} \left(- \phi_{x+1}^{t-1} - \phi_{x-1}^{t-1} + f_{0,x+1}^{t-1} + f_{0,x-1}^{t-1} + 2 f_{1,x+1}^{t-1} + 2 f_{2,x-1}^{t-1}\right),\end{split}
\end{equation*}\begin{equation*}
\begin{split}\displaystyle f_{0,x+1}^{t} + f_{0,x-1}^{t} - 2 f_{0,x}^{t} = - 2 \phi^{t}_{x} - \phi_{x+1}^{t-1} \left(s_{1} - 2\right) - \phi_{x-1}^{t-1} \left(s_{1} - 2\right) + \left(s_{1} - 1\right) \left(2 \phi_{x}^{t-2} + f_{0,x+1}^{t-1} + f_{0,x-1}^{t-1} - 2 f_{0,x}^{t-1}\right).\end{split}
\end{equation*}
\end{sphinxuseclass}\end{sphinxVerbatimOutput}

\end{sphinxuseclass}
\end{sphinxuseclass}
\sphinxAtStartPar
Replacing the above relation in \(\phi_{x}^{t+1}\) term:

\begin{sphinxuseclass}{cell}
\begin{sphinxuseclass}{tag_remove-input}\begin{sphinxVerbatimOutput}

\begin{sphinxuseclass}{cell_output}\begin{equation*}
\begin{split}\displaystyle \phi_{x}^{t+1} = - \phi^{t}_{x} \left(s_{1} - s_{2} w_{0} + s_{2} - 2\right) + \phi_{x}^{t-1} \left(s_{1} - 1\right) + \left(- \frac{\phi_{x+1}^t}{2} - \frac{\phi_{x-1}^t}{2}\right) \left(s_{1} + s_{2} w_{0} - 2\right) + \left(- \frac{\phi_{x+1}^{t-1}}{2} - \frac{\phi_{x-1}^{t-1}}{2}\right) \left(s_{1}^{2} + s_{1} s_{2} - 4 s_{1} - 2 s_{2} + 4\right) + \left(s_{1} - 1\right) \left(s_{1} + s_{2} - 2\right) \left(\phi_{x}^{t-2} + \frac{f_{0,x+1}^{t-1}}{2} + \frac{f_{0,x-1}^{t-1}}{2} - f_{0,x}^{t-1}\right).\end{split}
\end{equation*}
\end{sphinxuseclass}\end{sphinxVerbatimOutput}

\end{sphinxuseclass}
\end{sphinxuseclass}
\sphinxAtStartPar
Applying a time shift in the Eq. \eqref{equation:Refs/LBM-Study/FDM-Belloti/FDM-LBM-Diffusive:rel-pre-shift-2} to replace \(f^{t-1}_{0,x+1}/2+f^{t-1}_{0,x-1}/2-f^{t-1}_{0,x}\) in the above equation, we obtain in the final form \(\phi^{t+1}_{x}\) only dependent of previous moments:

\begin{sphinxuseclass}{cell}
\begin{sphinxuseclass}{tag_remove-input}\begin{sphinxVerbatimOutput}

\begin{sphinxuseclass}{cell_output}\begin{equation*}
\begin{split}\displaystyle \phi_{x}^{t+1} = \phi^{t}_{x} \left(s_{2} w_{0} - s_{2} + 1\right) + \phi_{x+1}^t \left(- \frac{s_{1}}{2} - \frac{s_{2} w_{0}}{2} + 1\right) + \phi_{x+1}^{t-1} \left(\frac{s_{1} s_{2} w_{0}}{2} - \frac{s_{1} s_{2}}{2} + \frac{s_{1}}{2} - \frac{s_{2} w_{0}}{2} + s_{2} - 1\right) + \phi_{x-1}^t \left(- \frac{s_{1}}{2} - \frac{s_{2} w_{0}}{2} + 1\right) + \phi_{x-1}^{t-1} \left(\frac{s_{1} s_{2} w_{0}}{2} - \frac{s_{1} s_{2}}{2} + \frac{s_{1}}{2} - \frac{s_{2} w_{0}}{2} + s_{2} - 1\right) + \phi_{x}^{t-1} \left(- s_{1} s_{2} w_{0} + s_{1} + s_{2} w_{0} - 1\right) + \phi_{x}^{t-2} \left(s_{1} s_{2} - s_{1} - s_{2} + 1\right).\end{split}
\end{equation*}
\end{sphinxuseclass}\end{sphinxVerbatimOutput}

\end{sphinxuseclass}
\end{sphinxuseclass}

\subsection{Fourth\sphinxhyphen{}order MRT Diffusive Equation}
\label{\detokenize{Refs/LBM-Study/FDM-Belloti/FDM-LBM-Diffusive:fourth-order-mrt-diffusive-equation}}
\sphinxAtStartPar
To recover the partial differetial equation with Fourth\sphinxhyphen{}order of convergence error, we expand each \(\phi\) term in Taylor series (this link provide a sympy code that reach the above equation \DUrole{xref,std,std-doc}{appendix/sympy\sphinxhyphen{}code\sphinxhyphen{}mrt\sphinxhyphen{}fdm\sphinxhyphen{}diffusive\sphinxhyphen{}expansion\sphinxhyphen{}4th}):
\begin{equation*}
\begin{split}
0=\frac{\Delta_{t}^{2} \left(- 3 s_{1} s_{2} + 2 s_{1} + 2 s_{2}\right)}{2} \frac{\partial^{2}\phi}{\partial t^{2}}  + \frac{\Delta_{t} \Delta_{x}^{2} \left(s_{1} s_{2} w_{0} - s_{1} s_{2} + s_{1} - s_{2} w_{0} + 2 s_{2} - 2\right)}{2} \frac{\partial^{3}\phi}{\partial x^{2}\partial t} + \Delta_{t} s_{1} s_{2} \frac{\partial \phi}{\partial t} + \frac{\Delta_{x}^{4} s_{2} \left(- s_{1} w_{0} + s_{1} + 2 w_{0} - 2\right)}{24} \frac{\partial^{4}\phi}{\partial x^{4}} + \frac{\Delta_{x}^{2} s_{2} \left(- s_{1} w_{0} + s_{1} + 2 w_{0} - 2\right)}{2} \frac{\partial^{2}\phi}{\partial x^{2}}
\end{split}
\end{equation*}\begin{equation*}
\begin{split}
\Delta_{t} s_{1} s_{2} \frac{\partial \phi}{\partial t} = - \frac{\Delta_{x}^{2} s_{2} \left(- s_{1} w_{0} + s_{1} + 2 w_{0} - 2\right)}{2} \frac{\partial^{2}\phi}{\partial x^{2}} -\frac{\Delta_{t}^{2} \left(- 3 s_{1} s_{2} + 2 s_{1} + 2 s_{2}\right)}{2} \frac{\partial^{2}\phi}{\partial t^{2}}  - \frac{\Delta_{t} \Delta_{x}^{2} \left(s_{1} s_{2} w_{0} - s_{1} s_{2} + s_{1} - s_{2} w_{0} + 2 s_{2} - 2\right)}{2} \frac{\partial^{3}\phi}{\partial x^{2}\partial t} - \frac{\Delta_{x}^{4} s_{2} \left(- s_{1} w_{0} + s_{1} + 2 w_{0} - 2\right)}{24} \frac{\partial^{4}\phi}{\partial x^{4}}
\end{split}
\end{equation*}\begin{equation*}
\begin{split}
\frac{\partial \phi}{\partial t} = \frac{\Delta_{x}^{2} (1-w_{0})\left(2 - s_{1}\right)}{2\Delta_{t} s_{1} } \frac{\partial^{2}\phi}{\partial x^{2}} - \frac{\Delta_{t} \left(- 3 s_{1} s_{2} + 2 s_{1} + 2 s_{2}\right)}{2 s_{1} s_{2}} \frac{\partial^{2}\phi}{\partial t^{2}}  - \frac{\Delta_{x}^{2} \left(s_{1} s_{2} w_{0} - s_{1} s_{2} + s_{1} - s_{2} w_{0} + 2 s_{2} - 2\right)}{2 s_{1} s_{2}} \frac{\partial^{3}\phi}{\partial x^{2}\partial t} + \frac{\Delta_{x}^{4} (1-w_{0})\left(2 - s_{1}\right)}{24\Delta_{t} s_{1}} \frac{\partial^{4}\phi}{\partial x^{4}}
\end{split}
\end{equation*}
\sphinxAtStartPar
Considering diffusive coeficient \(\nu=(1-w_{0})\left(\frac{1}{s_{1}} - \frac{1}{2}\right)\frac{\Delta_{x}^{2}}{\Delta_{t}}\), we can rewrite the equation
\begin{equation*}
\begin{split}
\frac{\partial \phi}{\partial t} = \nu\frac{\partial^{2}\phi}{\partial x^{2}} -\frac{\Delta_{t} \left(- 3 s_{1} s_{2} + 2 s_{1} + 2 s_{2}\right)}{2 s_{1} s_{2}} \frac{\partial^{2}\phi}{\partial t^{2}}  - \frac{ \Delta_{x}^{2} \left(s_{1} s_{2} w_{0} - s_{1} s_{2} + s_{1} - s_{2} w_{0} + 2 s_{2} - 2\right)}{2 s_{1} s_{2}} \frac{\partial^{3}\phi}{\partial x^{2}\partial t} + \frac{\nu\Delta_{x}^{2}}{12} \frac{\partial^{4}\phi}{\partial x^{4}}
\end{split}
\end{equation*}
\sphinxAtStartPar
Using the relations:
\begin{equation*}
\begin{split}
\frac{\partial\phi}{\partial t} = \nu \frac{\partial^{2}\phi}{\partial x^{2}}, \qquad \frac{\partial^{2}\phi}{\partial t^{2}} = \nu^{2} \frac{\partial^{4}\phi}{\partial x^{4}}, \qquad \frac{\partial^{3}\phi}{\partial x^{2}\partial t} = \nu \frac{\partial^{4}\phi}{\partial x^{4}},
\end{split}
\end{equation*}
\sphinxAtStartPar
we can reduce the equation in terms of \(\frac{\partial^{4}\phi}{\partial x^{4}}\):
\begin{equation*}
\begin{split}
\frac{\partial \phi}{\partial t} = \nu\frac{\partial^{2}\phi}{\partial x^{2}} - \frac{\Delta_{x}^{2} (1-w_{0})\left(2 - s_{1}\right)}{2 s_{1} }\frac{\nu \left(- 3 s_{1} s_{2} + 2 s_{1} + 2 s_{2}\right)}{2 s_{1} s_{2}} \frac{\partial^{4}\phi}{\partial x^{4}}  - \frac{ \Delta_{x}^{2} \nu \left(s_{1} s_{2} w_{0} - s_{1} s_{2} + s_{1} - s_{2} w_{0} + 2 s_{2} - 2\right)}{2 s_{1} s_{2}} \frac{\partial^{4}\phi}{\partial x^{4}} + \frac{\nu\Delta_{x}^{2}}{12} \frac{\partial^{4}\phi}{\partial x^{4}}
\end{split}
\end{equation*}\begin{equation*}
\begin{split}
\frac{\partial \phi}{\partial t} = \nu\frac{\partial^{2}\phi}{\partial x^{2}} + \frac{\Delta_{x}^{2} \nu \left(3 s_{1}^{2} s_{2} w_{0} - 2 s_{1}^{2} s_{2} - 6 s_{1}^{2} w_{0} - 18 s_{1} s_{2} w_{0} + 12 s_{1} s_{2} + 12 s_{1} w_{0} + 12 s_{2} w_{0} - 12 s_{2}\right)}{12 s_{1}^{2} s_{2}} \frac{\partial^{4}\phi}{\partial x^{4}}
\end{split}
\end{equation*}
\sphinxAtStartPar
Making the term \(\left(3 s_{1}^{2} s_{2} w_{0} - 2 s_{1}^{2} s_{2} - 6 s_{1}^{2} w_{0} - 18 s_{1} s_{2} w_{0} + 12 s_{1} s_{2} + 12 s_{1} w_{0} + 12 s_{2} w_{0} - 12 s_{2}\right)\) equal to zero e can reach fourth order approximation. For \(w_{0}=2/3\) (weight specific from the lattice D1Q3) the relation reduce to \(4s_{1}^{2} -8s_{1}+4s_{2}\):

\begin{sphinxuseclass}{cell}
\begin{sphinxuseclass}{tag_hide-input}\begin{sphinxVerbatimOutput}

\begin{sphinxuseclass}{cell_output}
\noindent\sphinxincludegraphics{{f13f7eb2a540b6c4594c7fcb3a5d9844eee60e8ab3ff77f49c85390f2c34255e}.png}

\end{sphinxuseclass}\end{sphinxVerbatimOutput}

\end{sphinxuseclass}
\end{sphinxuseclass}
\begin{sphinxuseclass}{cell}
\begin{sphinxuseclass}{tag_hide-input}
\begin{sphinxuseclass}{tag_hide-output}
\end{sphinxuseclass}
\end{sphinxuseclass}
\end{sphinxuseclass}
\begin{sphinxuseclass}{cell}
\begin{sphinxuseclass}{tag_hide-input}\begin{sphinxVerbatimOutput}

\begin{sphinxuseclass}{cell_output}
\noindent\sphinxincludegraphics{{f5be7b93d6dc918ea2c651f011fcbdfc14d1d6fd28b4ab852bd497ed8bb26a7d}.png}

\end{sphinxuseclass}\end{sphinxVerbatimOutput}

\end{sphinxuseclass}
\end{sphinxuseclass}
\begin{sphinxuseclass}{cell}
\begin{sphinxuseclass}{tag_hide-input}
\begin{sphinxuseclass}{tag_hide-output}
\end{sphinxuseclass}
\end{sphinxuseclass}
\end{sphinxuseclass}
\begin{sphinxuseclass}{cell}
\begin{sphinxuseclass}{tag_hide-input}
\end{sphinxuseclass}
\end{sphinxuseclass}
\begin{sphinxuseclass}{cell}
\begin{sphinxuseclass}{tag_hide-input}\begin{sphinxVerbatimOutput}

\begin{sphinxuseclass}{cell_output}
\noindent\sphinxincludegraphics{{5ca224fa5aa399b45f991d1a4bec6f72314efc8c911485128239cced27c56558}.png}

\end{sphinxuseclass}\end{sphinxVerbatimOutput}

\end{sphinxuseclass}
\end{sphinxuseclass}
\sphinxstepscope


\chapter{Citations}
\label{\detokenize{Refs/Citations:citations}}\label{\detokenize{Refs/Citations::doc}}\subsubsection*{References}

\begin{sphinxthebibliography}{CHGS18}
\bibitem[Anc94]{Refs/Citations:id5}
\sphinxAtStartPar
MG Ancona. Fully\sphinxhyphen{}lagrangian and lattice\sphinxhyphen{}boltzmann methods for solving systems of conservation equations. \sphinxstyleemphasis{Journal of Computational Physics}, 115(1):107–120, 1994.
\bibitem[BGM22]{Refs/Citations:id11}
\sphinxAtStartPar
Thomas Bellotti, Benjamin Graille, and Marc Massot. Finite difference formulation of any lattice boltzmann scheme. \sphinxstyleemphasis{Numerische Mathematik}, 152(1):1–40, 2022.
\bibitem[CHGS18]{Refs/Citations:id13}
\sphinxAtStartPar
Zhenhua Chai, Nanzhong He, Zhaoli Guo, and Baochang Shi. Lattice Boltzmann model for high\sphinxhyphen{}order nonlinear partial differential equations. \sphinxstyleemphasis{Physical Review E}, 2018. \sphinxhref{https://arxiv.org/abs/0907.1720}{arXiv:0907.1720}, \sphinxhref{https://doi.org/10.1103/PhysRevE.97.013304}{doi:10.1103/PhysRevE.97.013304}.
\bibitem[CCS23]{Refs/Citations:id8}
\sphinxAtStartPar
Ying Chen, Zhenhua Chai, and Baochang Shi. Fourth\sphinxhyphen{}order multiple\sphinxhyphen{}relaxation\sphinxhyphen{}time lattice boltzmann model and equivalent finite\sphinxhyphen{}difference scheme for one\sphinxhyphen{}dimensional convection\sphinxhyphen{}diffusion equations. \sphinxstyleemphasis{Physical Review E}, 107(5):055305, 2023.
\bibitem[DLOT17]{Refs/Citations:id7}
\sphinxAtStartPar
François Dubois, Pierre Lallemand, Christian Obrecht, and Mohamed Mahdi Tekitek. Lattice boltzmann model approximated with finite difference expressions. \sphinxstyleemphasis{Computers \& Fluids}, 155:3–8, 2017.
\bibitem[Jun01]{Refs/Citations:id6}
\sphinxAtStartPar
Michael Junk. A finite difference interpretation of the lattice boltzmann method. \sphinxstyleemphasis{Numerical Methods for Partial Differential Equations: An International Journal}, 17(4):383–402, 2001.
\bibitem[Kre10]{Refs/Citations:id12}
\sphinxAtStartPar
Gilberto M Kremer. \sphinxstyleemphasis{An introduction to the Boltzmann equation and transport processes in gases}. Springer Science \& Business Media, 2010.
\bibitem[LHSC21]{Refs/Citations:id10}
\sphinxAtStartPar
Yuxin Lin, Ning Hong, Baochang Shi, and Zhenhua Chai. Multiple\sphinxhyphen{}relaxation\sphinxhyphen{}time lattice boltzmann model\sphinxhyphen{}based four\sphinxhyphen{}level finite\sphinxhyphen{}difference scheme for one\sphinxhyphen{}dimensional diffusion equations. \sphinxstyleemphasis{Physical Review E}, 104(1):015312, 2021.
\bibitem[Sug10]{Refs/Citations:id4}
\sphinxAtStartPar
Shinsuke Suga. An accurate multi\sphinxhyphen{}level finite difference scheme for 1d diffusion equations derived from the lattice boltzmann method. \sphinxstyleemphasis{Journal of Statistical Physics}, 140(3):494–503, 2010.
\bibitem[Sug14]{Refs/Citations:id9}
\sphinxAtStartPar
Shinsuke Suga. Accurate numerical schemes based on the lattice boltzmann methods for multi\sphinxhyphen{}dimensional diffusion equations. \sphinxstyleemphasis{International Journal of Modern Physics C}, 25(04):1350104, 2014.
\bibitem[Wu15]{Refs/Citations:id3}
\sphinxAtStartPar
Yu\sphinxhyphen{}Shu Wu. \sphinxstyleemphasis{Multiphase fluid flow in porous and fractured reservoirs}. Gulf professional publishing, 2015.
\bibitem[ZDH+25]{Refs/Citations:id2}
\sphinxAtStartPar
Xueyi Zhang, Zhi Dou, Mayumi Hamada, Pietro de Anna, and Joaquin Jimenez\sphinxhyphen{}Martinez. Enhanced reaction kinetics in stationary two\sphinxhyphen{}phase flow through porous media. \sphinxstyleemphasis{Environmental Science \& Technology}, 2025.
\end{sphinxthebibliography}







\renewcommand{\indexname}{Index}
\printindex
\end{document}